\chapter{Complex pluripotential prerequisites}
\index{plurisubharmonic function!complex}\index{nonpluripolar product!complex}\index{multiplier ideal}%
\index{Lelong number!complex normalization}\index{singularity type!ordinary}%
\index{singularity type!P-equivalence@\(P\)-equivalence}%
\index{singularity type!I-equivalence@\(\mathcal I\)-equivalence}%
\index{envelope!model (P-envelope)@model (\(P\)-envelope)}%
\index{envelope!ideal (I-envelope)@ideal (\(\mathcal I\)-envelope)}
\label{chap:complex-prerequisites}

This chapter fixes the complex-analytic conventions used throughout
\cref{part:complex}.  The point is not to reproduce all of global pluripotential
theory, but to put the exact interfaces needed by test curves,
transcendental non-Archimedean metrics, and traces in one place.
Normalization matters: changing either \(\ddc\) or the exponent in a
multiplier ideal changes every numerical Lelong formula below.

The chapter has a second, less obvious purpose.  Almost every theorem in
\cref{part:complex,part:comparisons} is a statement about a
\emph{singularity type}, and there are three inequivalent notions of
singularity type in circulation.  Two potentials may differ by a bounded
function; they may have the same model envelope; they may have the same
multiplier ideals.  The first relation is the finest, the third the
coarsest, and the implications between them run in one direction only.
The trace operator of \Cref{chap:operations-trace} is well defined only
for the coarsest of the three; the Duistermaat--Heckman measures of
\Cref{chap:transcendental-okounkov-dh} are invariants of the middle one;
and the no-go theorem \Cref{thm:verified-mass-no-go} says that no
invariant of the coarsest one can see nonpluripolar mass.  A reader who
carries a single undifferentiated notion of ``same singularities'' through
those chapters will find several of their statements either trivial or
false.  Distinguishing the three, and proving exactly which implications
hold, is therefore not bookkeeping; it is the load-bearing content of this
chapter.

\section{Orientation}
\label{sec:ch51-orientation}

\subsection{What this chapter must supply}
\label{sec:ch51-what}

The material below is complex analytic, and none of it is
non-Archimedean.  Its natural cross-references therefore point
\emph{forward}, into \cref{part:complex,part:comparisons}, and sideways to
\Cref{app:complex-reference}; they do not point back into
\crefrange{part:foundations}{part:bj}.  A reader arriving here from
\Cref{chap:bj-stability} should treat this chapter as the start of a new
book whose only contact with the previous four parts is the normalization
concordance of \Cref{app:normalizations}.

Six things must be in place before \Cref{chap:rays-test-curves} can begin.

\begin{enumerate}
\item\label{itm:ch51-need:currents}
      \emph{Potentials and currents.}  The dictionary between closed
      positive \((1,1)\)-currents in a fixed Bott--Chern class and
      \(\theta\)-psh functions modulo constants, the compactness of the
      normalized set of potentials \emph{for nets}, and the potential
      \(V_\theta\) of minimal singularities.  Nets, not sequences: the
      inverse systems of \Cref{chap:transcendental-metrics} are indexed by
      K\"ahler forms and have no countable cofinal subset.
\item\label{itm:ch51-need:lelong}
      \emph{Lelong numbers.}  A single normalization, its behaviour under
      the operations that occur later (sums, maxima, scaling, pullback,
      increasing limits), and the passage from points to divisors over
      \(X\).  The increasing-limit statement,
      \Cref{lem:ch51-lelong-increasing}, is used four times below and is
      not a formal consequence of the definition.
\item\label{itm:ch51-need:ideals}
      \emph{Multiplier ideals.}  The exponent convention, Skoda's
      integrability bounds, the strong openness theorem, and the exact
      valuative criterion that recovers every multiplier ideal from
      divisorial Lelong numbers.  Without the last item the
      \(\mathcal I\)-preorder would be a definition with no content.
\item\label{itm:ch51-need:npp}
      \emph{Nonpluripolar products and masses.}  The construction, the
      properties that are used without comment later (plurifine locality,
      no mass on pluripolar sets, monotonicity of total masses,
      bimeromorphic invariance), and an exact statement of which
      continuity rules are available.
\item\label{itm:ch51-need:envelopes}
      \emph{The two envelopes.}  \(P_\theta[\varphi]\) and
      \(P_\theta[\varphi]_{\mathcal I}\), their idempotence, the mass
      identity for the first, and the inequality between them.
\item\label{itm:ch51-need:comparison}
      \emph{The comparison hierarchy.}  The implications
      \(\preceq\Rightarrow\preceq_P\Rightarrow\preceq_{\mathcal I}\), with
      proofs, and an explicit example showing that the second implication
      is strict.  That example is \Cref{ex:ch51-cantor}, and it is the one
      the rest of \cref{part:complex} leans on.
\end{enumerate}

Everything else --- relative energy, the \(d_1\) and \(d_S\) metrics,
weak geodesics, the Ross--Witt Nystr\"om correspondence --- is either
proved in \Cref{chap:rays-test-curves} or recorded in
\Cref{app:complex-reference}.  This chapter states what those developments
consume and no more.  The division is recorded item by item in
\Cref{sec:ch51-interface}, so that a reader can check any later citation
against a single list.

\subsection{Four comparisons, and why there are four}
\label{sec:ch51-four}

Let \(\varphi\) and \(\psi\) be quasi-psh functions on \(X\).  Four
relations between them occur in this book, and each is used somewhere in
\cref{part:complex,part:comparisons}:
\begin{enumerate}
\item pointwise order, \(\varphi\leq\psi\);
\item ordinary singularity type, \(\varphi\preceq\psi\), meaning
      \(\varphi\leq\psi+O(1)\);
\item \(P\)-comparison, \(\varphi\preceq_P\psi\), meaning
      \(P_\theta[\varphi]\leq P_\theta[\psi]\);
\item \(\mathcal I\)-comparison, \(\varphi\preceq_{\mathcal I}\psi\),
      meaning \(\mathcal I(m\varphi)\subseteq\mathcal I(m\psi)\) for every
      \(m\geq1\).
\end{enumerate}
Each is coarser than the one before it, and the induced equivalence
relations \(=\), \(\sim\), \(\sim_P\), \(\sim_{\mathcal I}\) are
correspondingly coarser.  The four are genuinely distinct.  Bounded
perturbations separate the first from the second.  The function
\(-\log(-\log|z|^2)\) of \Cref{ex:appH-loglog} separates the second from
the third: it is unbounded, hence not ordinarily equivalent to \(0\), but
its \(P\)-envelope is \(0\).  \Cref{ex:ch51-cantor} separates the third
from the fourth, by exhibiting a potential all of whose Lelong numbers
vanish --- so that it is \(\mathcal I\)-equivalent to \(0\) --- but whose
nonpluripolar mass is \(1/2\) while the mass of \(0\) is \(1\).  The
\(P\)-projection theorem then forbids \(P\)-equivalence.

There is a useful mnemonic.  The relation \(\sim\) remembers the
potential up to a bounded error; \(\sim_P\) remembers exactly the
data that nonpluripolar masses and relative energies can see; and
\(\sim_{\mathcal I}\) remembers exactly the data that divisorial
valuations can see.  The last of these is what analytification produces
(\Cref{chap:valuative-analytification}), and it is also the datum on which
the analytic trace of \Cref{chap:operations-trace} depends.  The target of
that operator is nevertheless a \(P\)-class: every trace is
\(\mathcal I\)-good, and two choices are \(P\)-equivalent.  This upgrade
from \(\mathcal I\)-data at the source to \(P\)-uniqueness at the target is
a theorem about restricted quasi-equisingular approximants, not a formal
property of the three quotients.

\subsection{Standing hypotheses}
\label{sec:ch51-hypotheses}

Throughout, \(X\) is a connected compact K\"ahler manifold of complex
dimension \(n\), and \(\theta\) is a smooth closed real \((1,1)\)-form.
Three hypotheses are quantified explicitly at every statement rather than
assumed globally.

\emph{Positivity of the class.}  We distinguish pseudoeffective (some
closed positive current in \(\{\theta\}\)), big (some K\"ahler current),
nef, and K\"ahler.  These are not synonyms, and the difference is visible
already in \Cref{ex:ch51-blowup-mass}, where a big non-nef class on a
blown-up surface has \(\int_X\{\theta\}^2=0\) but nonpluripolar volume
\(1\).

\emph{Positivity of the mass.}  Several theorems below, in particular
\Cref{thm:complex-p-projection}, require
\(\int_X\langle\theta_\varphi^n\rangle>0\).  The book reserves the term
\emph{model potential} for the positive-mass case; the zero-mass regime is
handled in \Cref{chap:transcendental-metrics} by a directed system of
K\"ahler perturbations, and no statement of this chapter is silently
extended to it.

\emph{Smoothness of \(X\).}  The manifold is smooth.  Compact K\"ahler
\emph{spaces} appear from \Cref{chap:tropical-spectra-complex} onwards, and
the modifications needed there are introduced at that point.  Nothing in
this chapter is asserted for singular \(X\).

\subsection{What is imported and what is proved}
\label{sec:ch51-imported}

\begin{important}{The five imported statements}
\label{imp:ch51-imported}
Everything below is proved from the definitions except for the following
five inputs, which are quoted with exact references and whose
normalizations are checked against \Cref{sec:ch51-normalizations} at the
point of use.
\begin{enumerate}
\item\label{itm:ch51-import:bt}
      \emph{Bedford--Taylor theory} for locally bounded psh functions:
      the mixed products \(\ddc u_1\wedge\cdots\wedge\ddc u_p\) exist as
      closed positive currents, are continuous along monotone sequences,
      and are local for the plurifine topology
      \cite{BT76,BT82}.  This underlies the whole of
      \Cref{sec:ch51-npp}.
\item\label{itm:ch51-import:skoda}
      \emph{Skoda's integrability theorem}: \(\rme^{-\varphi}\) is locally
      integrable near \(x\) if \(\nu(\varphi,x)<1\) and fails to be
      integrable if \(\nu(\varphi,x)>n\) \cite{Sko72}, in the
      normalization \eqref{eq:complex-lelong-number}.
\item\label{itm:ch51-import:openness}
      \emph{The strong openness theorem}
      \(\mathcal I(\varphi)=\bigcup_{\epsilon>0}\mathcal I((1+\epsilon)\varphi)\)
      \cite[Theorem~1.1]{GZ15}, together with the valuative criterion for
      multiplier ideals \cite[Corollary~10.18]{Bou17}.
\item\label{itm:ch51-import:regularization}
      \emph{Demailly regularization} and its quasi-equisingular refinement
      \cite{Dem92reg,DPS01}, used only in
      \Cref{prop:complex-quasi-equisingular} and in the construction of
      the trace in \Cref{chap:operations-trace}.
\item\label{itm:ch51-import:monotone}
      \emph{Monotonicity of nonpluripolar masses}
      \cite[Theorem~1.1]{DDNL18mono}, with the small-unbounded-locus
      predecessor \cite[Theorem~1.16]{BEGZ10}, and the increasing-sequence
      continuity theorem for nonpluripolar products
      \cite[Theorem~2.3]{DDNL18mono}.
\end{enumerate}
Items \cref{itm:ch51-import:bt}--\cref{itm:ch51-import:monotone} are the
only places where an analytic theorem is used without proof.  Every other
statement in this chapter, including the two projection theorems and the
whole comparison hierarchy, is derived from them.
\end{important}

Three statements are proved here that are not in the supplied book in this
form, and they are flagged where they occur: the increasing-limit formula
for Lelong numbers (\Cref{lem:ch51-lelong-increasing}), the bidegree-one
identification of the nonpluripolar product
(\Cref{lem:ch51-npp-bidegree-one}), and the explicit separation of
\(\sim_P\) from \(\sim_{\mathcal I}\) with computed masses
(\Cref{ex:ch51-cantor}).  \Cref{sec:ch51-additions} lists them again with
their reasons.

\subsection{Roadmap}
\label{sec:ch51-roadmap}

\Cref{sec:ch51-normalizations} fixes \(\ddc\), the Lelong normalization,
and the squared-norm convention for metric potentials, and gives three
unit tests that detect every convention error known to the author.
\Cref{sec:ch51-qpsh} develops quasi-psh functions and closed positive
currents.  \Cref{sec:ch51-lelong} treats Lelong numbers, including their
behaviour under monotone limits, and
\Cref{sec:ch51-multiplier} treats multiplier ideals, with a monomial
example carried to explicit generators.  \Cref{sec:ch51-npp} constructs
nonpluripolar products, states the mass inequalities exactly, and computes
a mass in a big non-K\"ahler class.  \Cref{sec:ch51-p-envelope,%
sec:ch51-i-envelope} construct the two envelopes, and
\Cref{sec:ch51-p-vs-i} proves the implications between the associated
equivalence relations and separates them by example.
\Cref{sec:ch51-interface} is the import ledger,
\Cref{sec:ch51-dictionary} the chapter dictionary, and
\Cref{sec:ch51-source} the source and historical note.

\section{Normalizations, and how to test them}
\label{sec:ch51-normalizations}

Throughout this chapter \(X\) is a connected compact K\"ahler manifold of
dimension \(n\).  We use
\[
             \dc=\frac{\rmi}{4\pi}(\bar\partial-\partial),
             \qquad
             \ddc=\frac{\rmi}{2\pi}\partial\bar\partial .
\tag{51.1}\label{eq:complex-ddc-normalization}
\]
Forms are real unless another coefficient field is displayed.
Positive products are always nonpluripolar products.  A psh or
quasi-psh function is allowed to take the value \(-\infty\), but is not
identically \(-\infty\).

This is the convention of \Cref{sec:appJ-complex-ddc}, and it is binding.
Its immediate consequence is the Poincar\'e--Lelong formula in the form
\[
 \ddc\log|f|^2=[\operatorname{div}(f)],
 \qquad
 \ddc\log|f|=\tfrac12[\operatorname{div}(f)],
\]
for a nonzero meromorphic function \(f\).  A source using
\(\ddc=\rmi\partial\bar\partial/\pi\) obtains the same Poincar\'e--Lelong
formula for \(\log|f|\) rather than for \(\log|f|^2\); a source using
\(\ddc=\rmi\partial\bar\partial\) obtains it with a factor \(2\pi\).
Every formula imported below has been checked against
\eqref{eq:complex-ddc-normalization}, and every rescaling is displayed.

\subsection{Potentials of metrics: the squared-norm convention}
\label{sec:ch51-squared-norm}

If \(L\) is a holomorphic line bundle with smooth Hermitian metrics
\(h_0\) and \(h\), the relative potential \(\varphi\) is defined by
\[
 |s|_h^2=\rme^{-\varphi}|s|_{h_0}^2 ,
\]
so that
\[
 c_1(L,h)=c_1(L,h_0)+\ddc\varphi .
\]
This is \eqref{eq:appJ-complex-weight} and
\eqref{eq:appJ-complex-curvature-change}.  The additive weight of the
underlying \emph{norm} is \(\varphi/2\).  The distinction is invisible
inside \cref{part:complex} and fatal at its boundary: a non-Archimedean
metric weight in the sense of \eqref{eq:appJ-na-weight} is a weight of the
norm, not of the squared norm, so every comparison between a complex ray
and a non-Archimedean potential carries a factor \(2\) that must be
inserted or cancelled explicitly.  \Cref{sec:appJ-metrics} is the
authority; \Cref{chap:dxz-bj-comparison} is where the factor is actually
used.

\begin{warning}{Half-logarithms}
\label{warn:ch51-half-log}
Several standard sources define the Lelong number of \(\varphi\) at \(x\)
as the supremum of \(\gamma\) with
\(\varphi\leq\gamma\log|z-x|+O(1)\), and the multiplier ideal by
integrability of \(|f|^2\rme^{-2\varphi}\).  Both conventions differ from
ours by a factor \(2\), and the two factors do \emph{not} cancel in mixed
formulas such as \eqref{eq:complex-lelong-from-ideals}.  When a result is
quoted from such a source below, the substitution
\(\varphi\rightsquigarrow\varphi/2\) is displayed in the proof rather than
performed silently.
\end{warning}

\subsection{Three unit tests}
\label{sec:ch51-unit-tests}

The following three computations detect, between them, every convention
error that occurred during the writing of \cref{part:complex}.  They are
repeated in condensed form in \Cref{sec:appH-unit-tests}.

\emph{Test one: one logarithmic pole.}  On the unit disk, the function
\(u=\log|z|^2\) satisfies
\[
 \ddc u=\delta_0,\qquad \nu(u,0)=1,\qquad
 \mathcal I(u)=(z),\qquad \mathcal I(\lambda u)=(z^{\lfloor\lambda\rfloor})
 \ \ (\lambda\notin\mathbf Z).
\]
This tests simultaneously the factor in \(\ddc\), the normalization of the
Lelong number, and the exponent in the multiplier ideal.  A convention in
which \(\nu(\log|z|^2,0)=2\) is the half-logarithm convention of
\Cref{warn:ch51-half-log}.

\emph{Test two: the Fubini--Study class on \(\mathbf P^1\).}  With
\(\omega=\ddc\log(1+|z|^2)\) on \(\mathbf C\subset\mathbf P^1\) one has
\(\int_{\mathbf P^1}\omega=1\) and, for every \(R>0\),
\[
 \int_{\{|z|<R\}}\omega=\frac{R^2}{1+R^2}.
\]
A convention producing total mass \(2\pi\) has omitted the \(2\pi\) in
\eqref{eq:complex-ddc-normalization}.  The Green function normalized by
\(\sup_{\mathbf P^1}G_p=0\) and \(\omega+\ddc G_p=\delta_p\) is then
\[
 G_p(z)=\log\frac{|z-p|^2}{(1+|z|^2)(1+|p|^2)},
\]
the logarithm of the squared chordal distance; in particular \(G_p\leq0\)
with equality exactly at the antipode of \(p\).  This normalization is used
in \Cref{ex:ch51-p1-envelope,ex:ch51-cantor}.

\emph{Test three: adding a constant.}  For every \(c\in\mathbf R\),
\[
 \MA_\theta(\varphi+c)=\MA_\theta(\varphi),
 \qquad
 M_\theta(\varphi+c)=M_\theta(\varphi),
\]
while relative energy changes by \(cM_\theta(\phi)\) and not by \(c\).  Any
formula in which a constant shifts the energy by \(c\) has divided by the
mass; see \Cref{sec:appJ-energy}.

\begin{remark}[Where the conventions are recorded]
\label{rem:ch51-convention-map}
The definitive statements are:
\eqref{eq:appJ-ddc} for \(\ddc\);
\eqref{eq:appJ-complex-pl} for Poincar\'e--Lelong;
\eqref{eq:appJ-complex-weight} for metric potentials;
\eqref{eq:appJ-npp} for the nonpluripolar mass and the prohibition of
probability normalization in the zero-volume case.  The present chapter
restates them so that \cref{part:complex} can be read without turning to
the appendix, and does not modify any of them.
\end{remark}

\section{Quasi-psh functions and positive currents}
\label{sec:ch51-qpsh}

The objects of this chapter are of two kinds, and it is worth saying at
once that they are the same objects seen from two sides.  On one side
there are closed positive \((1,1)\)-currents \(T\) on \(X\); on the other
there are upper semicontinuous functions \(\varphi\) with a positivity
condition.  The \(\partial\bar\partial\)-lemma on a compact K\"ahler
manifold makes the correspondence between them exact, and the whole of
\cref{part:complex} is written on the potential side because that is the
side on which suprema, infima, maxima, and Legendre transforms are
available.

\subsection{Potentials of a fixed class}
\label{sec:ch51-potentials}

Let \(\theta\) be a smooth closed \((1,1)\)-form.  On a coordinate open
set \(U\), choose \(g\in C^\infty(U)\) with \(\theta|_U=\ddc g\).

\begin{definition}[Potentials]
\label{def:complex-theta-psh}
A function \(\varphi:X\to[-\infty,+\infty)\) is
\(\theta\)-plurisubharmonic if it is upper semicontinuous and
\(g+\varphi\) is psh on every such \(U\).  We write
\[
                   \varphi\in\PSH(X,\theta).
\]
It is quasi-psh if it belongs to \(\PSH(X,\theta')\) for some smooth
closed \((1,1)\)-form \(\theta'\).  We write
\[
       \theta_\varphi:=\theta+\ddc\varphi .
\tag{51.2}\label{eq:complex-current-potential}
\]
\end{definition}

Thus \(\theta_\varphi\) is a closed positive current.  Conversely, every
closed positive \((1,1)\)-current in \(\{\theta\}\) has this form.
If \(\theta'=\theta+\ddc f\), then
\[
 \varphi\in\PSH(X,\theta)
 \quad\Longleftrightarrow\quad
 \varphi-f\in\PSH(X,\theta'),
\tag{51.3}\label{eq:complex-change-representative}
\]
and both sides define the same current.

Both assertions deserve a word.  That \(\theta_\varphi\) is closed is
immediate, since \(\theta\) is closed and \(\ddc\varphi\) is
\(\ddc\)-exact in the sense of currents; that it is positive is the local
condition \(\ddc(g+\varphi)\geq0\), which is the definition of
plurisubharmonicity of \(g+\varphi\).  The converse uses the
\(\partial\bar\partial\)-lemma on the compact K\"ahler manifold \(X\): if
\(T\geq0\) is closed with \(\{T\}=\{\theta\}\), then \(T-\theta\) is a
\(\ddc\)-exact real \((1,1)\)-current, so \(T-\theta=\ddc\varphi\) for some
\(\varphi\in L^1(X)\), unique up to an additive constant, and
\(\varphi\) agrees almost everywhere with an upper semicontinuous
representative because \(g+\varphi\) is psh on each \(U\).  We always work
with the upper semicontinuous representative.  The map
\[
 \PSH(X,\theta)/\mathbf R
 \longrightarrow
 \{\text{closed positive currents in }\{\theta\}\},
 \qquad
 \varphi\longmapsto\theta_\varphi ,
\]
is therefore a bijection.

Formula \eqref{eq:complex-change-representative} says that the choice of
smooth representative is immaterial for the current but shifts every
potential by a fixed smooth function.  Since \(f\) is bounded on the
compact \(X\), this shift changes neither the singularity type
\(\preceq\), nor the Lelong numbers, nor the multiplier ideals, nor any
nonpluripolar product.  It does change \(V_\theta\) and
\(P_\theta[\varphi]\) as \emph{functions}, which is why those two symbols
always carry the subscript \(\theta\) and never appear without it.

\begin{remark}[What is not being assumed]
\label{rem:ch51-no-positivity-assumed}
\Cref{def:complex-theta-psh} makes no positivity assumption on
\(\{\theta\}\); the set \(\PSH(X,\theta)\) may well be empty, and it is
nonempty exactly when \(\{\theta\}\) is pseudoeffective.  Bigness,
nefness, and the K\"ahler property are extra hypotheses, imposed
individually.  In \Cref{chap:transcendental-metrics} the class is allowed
to be merely pseudoeffective and even of vanishing volume, which is the
reason for the perturbative inverse system used there.
\end{remark}

\subsection{Elementary operations}
\label{sec:ch51-operations}

The following closure properties are used constantly and are all local
consequences of the corresponding statements for psh functions.

\begin{proposition}\label{prop:ch51-psh-operations}
Let \(\theta,\theta'\) be smooth closed real \((1,1)\)-forms.
\begin{enumerate}
\item\label{itm:ch51-op:max}
      If \(\varphi,\psi\in\PSH(X,\theta)\), then
      \(\max\{\varphi,\psi\}\in\PSH(X,\theta)\).
\item\label{itm:ch51-op:convex}
      If \(\varphi,\psi\in\PSH(X,\theta)\) and \(t\in[0,1]\), then
      \((1-t)\varphi+t\psi\in\PSH(X,\theta)\).
\item\label{itm:ch51-op:sum}
      If \(\varphi\in\PSH(X,\theta)\) and \(\psi\in\PSH(X,\theta')\),
      then \(\varphi+\psi\in\PSH(X,\theta+\theta')\), and
      \(\lambda\varphi\in\PSH(X,\lambda\theta)\) for \(\lambda>0\).
\item\label{itm:ch51-op:sup}
      If \(\{\varphi_i\}\subset\PSH(X,\theta)\) is a family that is
      uniformly bounded above, then
      \(\uscsup_i\varphi_i\in\PSH(X,\theta)\), where the star denotes the
      upper semicontinuous regularization.  Moreover
      \(\uscsup_i\varphi_i=\sup_i\varphi_i\) outside a pluripolar set.
\item\label{itm:ch51-op:decreasing}
      If \((\varphi_j)\subset\PSH(X,\theta)\) is a decreasing sequence,
      then its pointwise limit is either identically \(-\infty\) or lies
      in \(\PSH(X,\theta)\).
\item\label{itm:ch51-op:pullback}
      If \(\pi:Y\to X\) is a holomorphic map from a compact K\"ahler
      manifold whose image is not contained in \(\{\varphi=-\infty\}\),
      then \(\pi^*\varphi\in\PSH(Y,\pi^*\theta)\).
\end{enumerate}
\end{proposition}

\begin{proof}
All six assertions are local, and locally \(\varphi\) is a psh function
minus a smooth one.  Items \cref{itm:ch51-op:max} and
\cref{itm:ch51-op:convex} are the corresponding statements for psh
functions, because \(g+\max\{\varphi,\psi\}=\max\{g+\varphi,g+\psi\}\) and
\(g+(1-t)\varphi+t\psi=(1-t)(g+\varphi)+t(g+\psi)\).
Item \cref{itm:ch51-op:sum} follows from
\((g+g')+(\varphi+\psi)=(g+\varphi)+(g'+\psi)\), a sum of psh functions,
and homogeneity of the psh condition under positive scaling.
Item \cref{itm:ch51-op:sup} is the classical fact that the upper
semicontinuous regularization of a locally bounded-above family of psh
functions is psh, together with the theorem of Bedford--Taylor that the
set where the regularization differs from the supremum is negligible,
hence pluripolar \cite[Theorem~7.1]{BT82}.
Item \cref{itm:ch51-op:decreasing} is the classical decreasing-limit
theorem.  Item \cref{itm:ch51-op:pullback} follows because
\(\pi^*(g+\varphi)\) is psh whenever \(g+\varphi\) is psh and
\(\pi^*\varphi\not\equiv-\infty\) on each component.
\end{proof}

\begin{warning}{Regularization is not free}
\label{warn:ch51-usc}
Item \cref{itm:ch51-op:sup} of \Cref{prop:ch51-psh-operations} produces \(\uscsup_i\varphi_i\), not
\(\sup_i\varphi_i\).  The two differ on a pluripolar set, which is
invisible to every nonpluripolar product but is emphatically visible to a
pointwise inequality.  All envelopes in this chapter --- \(V_\theta\), the
rooftop \(\varphi\wedge_\theta\psi\), \(P_\theta[\varphi]\), and
\(P_\theta[\varphi]_{\mathcal I}\) --- are defined with the
regularization.  The one place where inserting an extra regularization
would change a statement is the inverse Legendre transform of
\Cref{sec:appH-test-curves}, where the infimum is already upper
semicontinuous under the hypotheses of the Ross--Witt Nystr\"om
correspondence; see \Cref{thm:rays-rwn}.
\end{warning}

\subsection{Comparison of singularities}
\label{sec:ch51-ordinary-order}

\begin{definition}[Comparison up to a bounded error]
\label{def:complex-singularity-order}
For quasi-psh functions \(\varphi,\psi\), write
\[
 \varphi\preceq\psi
 \quad\Longleftrightarrow\quad
 \varphi\leq\psi+C\ \text{on }X\ \text{for some }C\in\mathbf R.
\tag{51.4}\label{eq:complex-singularity-order}
\]
Then \(\varphi\) is said to be more singular than \(\psi\).
Write \(\varphi\sim\psi\) if both comparisons hold.
\end{definition}

The direction in \eqref{eq:complex-singularity-order} is used without
exception: smaller potentials are more singular.  Adding a bounded
function does not change a singularity type.  The maximum
\(\varphi\vee\psi\) is less singular than either input.

Three comments on \Cref{def:complex-singularity-order} are worth making
before any theorem uses it.

First, \(\preceq\) is defined for quasi-psh functions on \(X\), not for
elements of a fixed \(\PSH(X,\theta)\); the definition compares functions,
and the reference forms play no role.  Consequently
\eqref{eq:complex-change-representative} does not affect it: replacing
\(\theta\) by \(\theta+\ddc f\) subtracts the bounded function \(f\) from
every potential.

Second, \(\preceq\) is a preorder and not an order.  Its associated
equivalence classes are the \emph{ordinary singularity types}, and the
quotient
\[
 \QPSH(X)/{\sim}
\]
is the finest of the three quotients used in this book.  The two coarser
quotients \(\QPSH(X)/{\sim_P}\) and \(\QPSH(X)/{\sim_{\mathcal I}}\) are
constructed in \Cref{sec:ch51-p-envelope,sec:ch51-i-envelope}, and the
canonical surjections between them are used in
\eqref{eq:verified-trace-factorization}.

Third, the elementary calculus of \(\preceq\) is exactly what one expects.
If \(\varphi\preceq\psi\) and \(\varphi'\preceq\psi'\), then
\(\varphi+\varphi'\preceq\psi+\psi'\),
\(\max\{\varphi,\varphi'\}\preceq\max\{\psi,\psi'\}\), and
\(\lambda\varphi\preceq\lambda\psi\) for \(\lambda>0\).  Also
\(\varphi\wedge_\theta\psi\preceq\varphi\) and
\(\varphi\wedge_\theta\psi\preceq\psi\), where the rooftop is
\eqref{eq:complex-rooftop}, while
\(\max\{\varphi,\psi\}\) is the least upper bound for \(\preceq\) among
\(\theta\)-psh functions.  All of these are immediate from the
definitions; they are recorded because
\Cref{prop:operations-trace-calculus} asserts the corresponding
statements for \(\preceq_{\mathcal I}\) and one should know which of them
are formal and which are not.

\begin{example}[The three types on a disk]
\label{ex:ch51-disk-three-types}
On the unit disk \(\mathbf D\subset\mathbf C\), consider
\[
 \begin{gathered}
 u_1=\log|z|^2,
 \qquad
 u_2=\log|z|^2+\log(-\log|z|^2)^{-1},\\
 u_3=-\log(-\log|z|^2),
 \qquad
 u_4=0 .
 \end{gathered}
\]
The second and third are psh near \(0\) after shrinking the disk.  One has
\(u_1\preceq u_2\preceq u_3\preceq u_4\) but no reverse relation, and
\[
 \nu(u_1,0)=1,\qquad
 \nu(u_2,0)=1,\qquad
 \nu(u_3,0)=0,\qquad
 \nu(u_4,0)=0 .
\]
Thus \(u_1\sim_{\mathcal I}u_2\) and \(u_3\sim_{\mathcal I}u_4\), whereas
\(u_1\not\sim u_2\) and \(u_3\not\sim u_4\).  This exhibits the first
strictness in the hierarchy, and it does so already locally: no
compactness, no cohomology class, and no mass is needed.  What it does
\emph{not} exhibit is any separation of \(\sim_P\) from
\(\sim_{\mathcal I}\); that separation is global and is the content of
\Cref{ex:ch51-cantor}.
\end{example}

\subsection{Compactness}
\label{sec:ch51-compactness}

The following compactness statement is the analytic engine behind every
existence argument in \cref{part:complex}.  It is stated for nets because
the transcendental construction of \Cref{chap:transcendental-metrics}
is indexed by the K\"ahler forms \(\omega\) on \(X\), directed by
\(\omega\geq\omega'\), and this index set has no countable cofinal subset
in general.

\begin{proposition}\label{prop:complex-l1-compactness}
For fixed \(\theta\), the set
\[
     \{\varphi\in\PSH(X,\theta)\mid \sup_X\varphi=0\}
\]
is compact in \(L^1(X)\).  If a decreasing net
\((\varphi_i)\subset\PSH(X,\theta)\) has
\(\inf_i\sup_X\varphi_i>-\infty\), then its pointwise infimum is
\(\theta\)-psh and the net converges to it in \(L^1\).
\end{proposition}

\begin{proof}
Local psh compactness gives \(L^1_{\mathrm{loc}}\)-compactness after
normalization.  A finite cover and the compatibility of local potentials
globalize it.  The map
\(\eta\mapsto\sup_X\eta\) is continuous on
\(\PSH(X,\theta)\) for the \(L^1\)-topology.

For the second assertion, fix \(i_0\) and pass to the cofinal tail
\(i\geq i_0\).  On this tail the numbers \(\sup_X\varphi_i\) are bounded
both above and below.  Compactness after normalization, together with
the continuity of the supremum, therefore puts the tail in an
\(L^1\)-compact set \(K\).
Set \(g:=\inf_i\varphi_i\), pointwise.  Let a subnet converge in
\(L^1\) to \(\psi\in\PSH(X,\theta)\).  For each fixed \(i\), the subnet
is eventually contained in the tail indexed by \(i\), so Hartogs'
lemma gives \(\psi\leq\varphi_i\), hence \(\psi\leq g\).
Choose recursively terms of this convergent subnet whose original
indices increase and whose \(L^1\)-distance to \(\psi\) is less than
\(2^{-m}\).  The resulting sequence decreases pointwise.  The classical
decreasing-sequence theorem for psh functions shows that its pointwise
limit \(h\) is psh or identically \(-\infty\); its \(L^1\)-convergence
to \(\psi\) rules out the latter and gives \(h=\psi\).  Since \(g\) is
the infimum over the whole net, \(g\leq h=\psi\).  Combined with the
opposite inequality above, this gives \(g=\psi\) pointwise and, in
particular, proves that \(g\) is \(\theta\)-psh.  We have proved that every
\(L^1\)-convergent subnet has the same limit \(g\).  A net in a compact
Hausdorff space whose every cluster point is \(g\) converges to \(g\):
otherwise a subnet outside some neighborhood of \(g\) would have a
different cluster point.  Therefore \(\varphi_i\to g\) in \(L^1\).
This genuine subnet argument is needed because the inverse systems in
\Cref{chap:transcendental-metrics} need not have countable
cofinal subsets.
\end{proof}

Two corollaries are used repeatedly and are worth isolating.

\begin{corollary}\label{cor:ch51-hartogs}
Let \((\varphi_i)\subset\PSH(X,\theta)\) be a net converging in
\(L^1(X)\) to \(\varphi\in\PSH(X,\theta)\).  Then:
\begin{enumerate}
\item\label{itm:ch51-hartogs:sup}
      \(\sup_X\varphi_i\to\sup_X\varphi\);
\item\label{itm:ch51-hartogs:hartogs}
      for every compact \(K\subset X\) and every continuous
      \(h:K\to\mathbf R\) with \(\varphi<h\) on \(K\), one has
      \(\varphi_i<h\) on \(K\) for \(i\) large;
\item\label{itm:ch51-hartogs:limsup}
      \(\limsup_i\varphi_i\leq\varphi\) pointwise, with equality outside a
      pluripolar set whenever the net is monotone.
\end{enumerate}
\end{corollary}

\begin{proof}
Item \cref{itm:ch51-hartogs:sup} is the continuity of the supremum
functional noted in the proof of \Cref{prop:complex-l1-compactness}: on
\(\PSH(X,\theta)\) the \(L^1\)-topology and the topology of local uniform
convergence of the sup-functional agree, because a normalized family is
\(L^1\)-compact and \(\sup_X\) is upper semicontinuous and, by the
sub-mean-value inequality applied on a fixed finite cover by coordinate
balls, also lower semicontinuous along \(L^1\)-convergent nets.

Item \cref{itm:ch51-hartogs:hartogs} is Hartogs' lemma in its net form.
Fix a finite cover of \(X\) by coordinate balls \(B(x_a,3r_a)\) on which
\(g_a+\varphi_i\) is psh for every \(i\), with \(g_a\) smooth and
independent of \(i\).  For \(x\in B(x_a,r_a)\) and \(0<s<r_a\) the
sub-mean-value inequality gives
\[
 (g_a+\varphi_i)(x)
 \leq\frac{1}{|B(x,s)|}\int_{B(x,s)}(g_a+\varphi_i)\,\rd V
 =:A_{i,a}(x,s).
\]
Since \(\varphi_i\to\varphi\) in \(L^1(X)\), we have
\(A_{i,a}(x,s)\to A_{\infty,a}(x,s)\) uniformly in \(x\) for each fixed
\(s\), the error being at most
\(|B(x,s)|^{-1}\|\varphi_i-\varphi\|_{L^1}\).  Moreover
\(A_{\infty,a}(\cdot,s)\) is continuous, and it decreases as
\(s\downarrow0\) to \((g_a+\varphi)(\cdot)\).

Let now \(h\) be continuous on the compact \(K\) with \(\varphi<h\) there,
and put \(3\epsilon:=\min_K(h-\varphi)\), a positive number because
\(h-\varphi\) is lower semicontinuous and positive on \(K\).  Applying
Dini's theorem to the decreasing family of continuous functions
\(x\mapsto A_{\infty,a}(x,s)-g_a(x)\) on the compact set
\(K\cap\overline{B(x_a,r_a)}\), whose limit \(\varphi\) is bounded above
there by \(h-3\epsilon\), we obtain \(s_a>0\) with
\[
 A_{\infty,a}(x,s_a)-g_a(x)<h(x)-2\epsilon
 \qquad\text{on }K\cap\overline{B(x_a,r_a)} .
\]
Choosing \(i\) so large that
\(\|\varphi_i-\varphi\|_{L^1}<\epsilon\min_a\inf_x|B(x,s_a)|\) gives
\(\varphi_i(x)\leq A_{i,a}(x,s_a)-g_a(x)<h(x)-\epsilon\) on
\(K\cap\overline{B(x_a,r_a)}\) for every \(a\).  The balls
\(B(x_a,r_a)\) cover \(X\), so this proves the claim.

Item \cref{itm:ch51-hartogs:limsup} is the pointwise form of the same
inequality, together with \Cref{prop:ch51-psh-operations}, item
\cref{itm:ch51-op:sup}, in the increasing case and
\Cref{prop:complex-l1-compactness} in the decreasing case.
\end{proof}

\begin{corollary}\label{cor:ch51-choquet}
Let \(\{\varphi_i\}_{i\in I}\subset\PSH(X,\theta)\) be a family that is
uniformly bounded above, and set \(\varphi=\uscsup_i\varphi_i\).  Then
there is a countable subset \(\{i_1,i_2,\ldots\}\subset I\) such that
\[
 \varphi=\uscsup_a\max\{\varphi_{i_1},\ldots,\varphi_{i_a}\},
\]
and the finite maxima on the right increase.  If \(I\) is directed and
\(i\mapsto\varphi_i\) is increasing, the \(i_a\) may be chosen increasing
and the finite maxima replaced by the \(\varphi_{i_a}\) themselves.
\end{corollary}

\begin{proof}
This is Choquet's lemma \cite{DemBook}: the upper
semicontinuous regularization of the supremum of an arbitrary family
coincides with that of a suitable countable subfamily, because the
topology of \(X\) has a countable basis and the regularization is
determined by the values of the supremum on that basis.  Replacing the
countable subfamily by its finite maxima makes the sequence increasing
without changing the regularized supremum, by
\Cref{prop:ch51-psh-operations}, item \cref{itm:ch51-op:max}.  If the
family is a directed increasing net, choose \(i_a\) inductively above
\(i_{a-1}\) and above \(i_1,\ldots,i_a\); then
\(\max\{\varphi_{i_1},\ldots,\varphi_{i_a}\}\leq\varphi_{i_a}\leq\varphi\).
\end{proof}

\begin{warning}{Choquet does not make a net a sequence}
\label{warn:ch51-choquet}
\Cref{cor:ch51-choquet} extracts a countable subfamily with the same
regularized supremum.  It does not produce a countable cofinal subset of
\(I\), and it says nothing about decreasing families.  Every use of a
Choquet reduction below is a use for \emph{increasing} families and for a
quantity --- a regularized supremum, or a mass --- that is determined by
that supremum.  Statements about decreasing nets, such as the second half
of \Cref{prop:complex-l1-compactness}, are proved by a genuine subnet
argument.  The distinction is the reason
\Cref{cor:operations-increasing-net} and
\Cref{rem:operations-nets} are stated separately in
\Cref{chap:operations-trace}.
\end{warning}

\subsection{Minimal singularities and positivity of the class}
\label{sec:ch51-minimal}

\begin{definition}[Minimal singularities]
\label{def:complex-minimal-potential}
If \(\{\theta\}\) is pseudoeffective, set
\[
 V_\theta
 :=\uscsup\{\psi\in\PSH(X,\theta)\mid \psi\leq0\}.
\tag{51.5}\label{eq:complex-vtheta}
\]
Then \(V_\theta\) has minimal singularities: every
\(\psi\in\PSH(X,\theta)\) satisfies \(\psi\preceq V_\theta\).
The class is big exactly when it contains a K\"ahler current.
\end{definition}

\begin{remark}[Convention checkpoint]
The symbols \(\theta_\varphi\) and \(V_\theta\) depend on a smooth
representative, while the underlying current and all equivalence
relations below are invariant under
\eqref{eq:complex-change-representative}.
\end{remark}

The two assertions of \Cref{def:complex-minimal-potential} are proved as
follows.  That \(V_\theta\in\PSH(X,\theta)\) is
\Cref{prop:ch51-psh-operations}, item \cref{itm:ch51-op:sup}, applied to a
family that is nonempty (pseudoeffectivity) and bounded above by \(0\); one
has \(V_\theta\leq0\) because the regularization of a family of
nonpositive upper semicontinuous functions is nonpositive.  If
\(\psi\in\PSH(X,\theta)\), then \(\psi-\sup_X\psi\leq0\) is a competitor,
so \(\psi\leq V_\theta+\sup_X\psi\), which is
\(\psi\preceq V_\theta\).  Thus \(V_\theta\) is a maximal element for
\(\preceq\), and any two such differ by a bounded function.

\begin{definition}[Positivity of a class]
\label{def:ch51-positivity}
The class \(\{\theta\}\) is
\begin{enumerate}
\item \emph{pseudoeffective} if \(\PSH(X,\theta)\neq\emptyset\);
\item \emph{big} if there are \(\varphi\in\PSH(X,\theta)\), a K\"ahler
      form \(\omega\), and \(\epsilon>0\) with
      \(\theta_\varphi\geq\epsilon\omega\); such a
      \(\theta_\varphi\) is a \emph{K\"ahler current};
\item \emph{nef} if for every \(\epsilon>0\) there is a smooth
      \(\varphi_\epsilon\) with
      \(\theta+\ddc\varphi_\epsilon\geq-\epsilon\omega\);
\item \emph{K\"ahler} if \(\theta+\ddc\varphi\) is a smooth positive
      definite form for some smooth \(\varphi\).
\end{enumerate}
\end{definition}

These four conditions are strictly decreasing in strength.  A K\"ahler
class is nef and big; a nef class is pseudoeffective but need not be big;
a big class need not be nef; a pseudoeffective class need not be big.
\Cref{ex:ch51-blowup-mass} realizes a big non-nef class on a blown-up
surface and computes its nonpluripolar volume, and
\Cref{ex:ch51-nef-zero-volume} realizes a nef class of vanishing volume on
the same surface.  Neither example is exotic, and both occur in
\Cref{chap:transcendental-metrics}.

\begin{proposition}\label{prop:ch51-vtheta-bounded-iff}
Let \(\{\theta\}\) be pseudoeffective.
\begin{enumerate}
\item\label{itm:ch51-vtheta:kahler}
      \(V_\theta\) is bounded if and only if for some, equivalently every,
      K\"ahler form \(\omega\) and every \(\epsilon>0\) the class
      \(\{\theta\}+\epsilon\{\omega\}\) contains a smooth positive form;
      that is, if and only if \(\{\theta\}\) is nef.
\item\label{itm:ch51-vtheta:continuous}
      If \(\{\theta\}\) is K\"ahler, then \(V_\theta\) may be taken to be
      \(0\) after changing the smooth representative.
\item\label{itm:ch51-vtheta:big}
      If \(\{\theta\}\) is big, \(V_\theta\) is locally bounded on a
      nonempty Zariski-open subset, the \emph{ample locus} of
      \(\{\theta\}\), but need not be bounded on \(X\).
\end{enumerate}
\end{proposition}

\begin{proof}
For item \cref{itm:ch51-vtheta:kahler}, suppose \(V_\theta\) is bounded.
Then \(\theta+\ddc V_\theta\geq0\) with bounded potential; regularizing
\(V_\theta\) by Demailly's approximation with loss \(\epsilon\omega\)
produces smooth \(\varphi_\epsilon\) with
\(\theta+\ddc\varphi_\epsilon\geq-\epsilon\omega\), so \(\{\theta\}\) is
nef.  Conversely, if \(\{\theta\}\) is nef, the smooth potentials
\(\varphi_\epsilon\) satisfy
\(\varphi_\epsilon\in\PSH(X,\theta+\epsilon\omega)\); normalizing
\(\sup_X\varphi_\epsilon=0\) and using
\Cref{prop:complex-l1-compactness} to extract an \(L^1\)-limit as
\(\epsilon\downarrow0\) shows \(V_\theta\geq\lim\varphi_\epsilon>-\infty\)
almost everywhere; the uniform bound on \(\sup_X\) and the sub-mean-value
inequality then give a uniform lower bound.  Item
\cref{itm:ch51-vtheta:continuous} is immediate: if
\(\theta+\ddc f\) is a K\"ahler form, replace \(\theta\) by
\(\theta+\ddc f\) and use \eqref{eq:complex-change-representative}; then
\(0\in\PSH(X,\theta)\) and \(V_\theta=0\).  Item
\cref{itm:ch51-vtheta:big} combines two imported facts: a big class
contains a K\"ahler current with analytic singularities
\cite{Dem92reg,DPS01}, and the ample locus, defined as the complement of
the intersection of the singular loci of all such currents, is a nonempty
Zariski-open set \cite{Bou04}.  On the ample locus, \(V_\theta\) is
locally bounded, because it dominates the locally bounded potential of any
one such current up to a constant.
\end{proof}

\begin{remark}[The zero-volume trap]
\label{rem:ch51-zero-volume}
Even for nef \(\{\theta\}\) the number \(\int_X\langle\theta_{V_\theta}^n
\rangle\) may vanish, in which case there is no probability
normalization of the Monge--Amp\`ere measure and no relative finite-energy
theory over \(\theta\); see \eqref{eq:appJ-npp} and
\Cref{ex:ch51-nef-zero-volume}.  This is the single most common source of
vacuous statements in transcendental non-Archimedean geometry, and it is
the reason \Cref{chap:transcendental-metrics} builds its inverse system out
of the strictly positive perturbations \(\theta+\omega\) rather than out of
\(\theta\) itself.
\end{remark}

\section{Lelong numbers}
\label{sec:ch51-lelong}

The normalization \eqref{eq:complex-ddc-normalization} fixes our Lelong
number.  If \(F:U\to\mathbf C^n\) is a coordinate chart at \(x\), define
\[
 \nu(\varphi,x)
 =\sup\left\{\gamma\geq0\ \middle|\ 
 \varphi\circ F^{-1}(z)
 \leq\gamma\log|z-F(x)|^2+O(1)\right\}.
\tag{51.6}\label{eq:complex-lelong-number}
\]
In particular,
\[
                         \nu(\log|z|^2,0)=1.
\tag{51.7}\label{eq:complex-lelong-model}
\]
This is twice the value arising from conventions based on \(\log|z|\).

\subsection{Three descriptions of the same number}
\label{sec:ch51-lelong-three}

Definition \eqref{eq:complex-lelong-number} is the one used in this book
because it is manifestly a statement about singularity types: it depends
only on \(\varphi\) modulo bounded functions, and it is monotone,
\[
 \varphi\preceq\psi
 \quad\Longrightarrow\quad
 \nu(\varphi,x)\geq\nu(\psi,x).
\]
Two other descriptions are used in proofs, and their agreement with
\eqref{eq:complex-lelong-number} is the imported content of this
subsection.

Fix a coordinate chart identifying a neighbourhood of \(x\) with a ball in
\(\mathbf C^n\) centred at the origin, and let \(u\) be psh there.  Write
\[
 \lambda_u(t)
 :=\frac{1}{\sigma(S_r)}\int_{S_r}u\,\rd\sigma,
 \qquad
 t=\log r^2,
\]
for the mean of \(u\) over the sphere \(S_r=\{|z|=r\}\), and
\[
 M_u(t):=\sup_{|z|\leq r}u,\qquad t=\log r^2 .
\]

\begin{proposition}\label{prop:ch51-lelong-three}
Let \(u\) be psh on a ball around \(0\) in \(\mathbf C^n\).
\begin{enumerate}
\item\label{itm:ch51-lelong:convex}
      The functions \(\lambda_u\) and \(M_u\) are convex and increasing on
      \((-\infty,t_0)\), and \(\lambda_u\leq M_u\).
\item\label{itm:ch51-lelong:limits}
      The three limits
      \[
       \lim_{t\to-\infty}\frac{\lambda_u(t)}{t},
       \qquad
       \lim_{t\to-\infty}\frac{M_u(t)}{t},
       \qquad
       \sup\{\gamma\geq0\mid u\leq\gamma t+O(1)\ \text{as}\ t\to-\infty\}
      \]
      exist and coincide; their common value is \(\nu(u,0)\).
\item\label{itm:ch51-lelong:mass}
      One also has
      \[
       \nu(u,0)
       =\lim_{r\to0}
         \frac{\displaystyle\int_{|z|<r}\ddc u\wedge(\ddc|z|^2)^{n-1}}
              {\displaystyle\int_{|z|<r}
                 \ddc\log|z|^2\wedge(\ddc|z|^2)^{n-1}},
      \]
      the mass of \(\ddc u\) on small balls normalized against the model
      pole; the ratio makes the normalizing constant explicit and
      reproduces \eqref{eq:complex-lelong-model} by inspection.
\item\label{itm:ch51-lelong:slope}
      \(\nu(u,0)=\inf_{t}\lambda_u'(t)=\lim_{t\to-\infty}\lambda_u'(t)\),
      where \(\lambda_u'\) denotes the right derivative.
\end{enumerate}
\end{proposition}

\begin{proof}[Imported, with the normalization checked]
Items \cref{itm:ch51-lelong:convex}--\cref{itm:ch51-lelong:mass} are
Demailly's theory of Lelong numbers \cite[Chapter~III]{DemBook}; see also
\cite{Lel68,Siu74,Dem93}.  What must be checked is only the
normalization.  In our convention \(\ddc=\rmi\partial\bar\partial/2\pi\),
so \(\ddc\log|z|^2=\delta_0\) in dimension one and, in dimension \(n\),
\((\ddc\log|z|^2)^n=\delta_0\).  Testing item
\cref{itm:ch51-lelong:mass} on \(u=\log|z|^2\) gives the value \(1\), which
matches \eqref{eq:complex-lelong-model}; a source normalizing by
\(\log|z|\) obtains \(2\) here.  For item \cref{itm:ch51-lelong:slope},
convexity of \(\lambda_u\) in \(t\) makes \(\lambda_u'\) increasing, so
the limit at \(-\infty\) is the infimum, and it equals
\(\lim_{t\to-\infty}\lambda_u(t)/t\) because for \(t<t_1\)
\[
 \frac{\lambda_u(t)}{t}
 =\frac{\lambda_u(t_1)}{t}
   -\frac{1}{t}\int_t^{t_1}\lambda_u'(s)\,\rd s
 \longrightarrow\lim_{s\to-\infty}\lambda_u'(s).
\]
Finally, the agreement of the sup-description
\eqref{eq:complex-lelong-number} with the two mean-value descriptions is
the substance of item \cref{itm:ch51-lelong:limits}.  One inequality is the
sub-mean-value property, which gives
\(\lambda_u\leq M_u\) and hence
\(\lim M_u(t)/t\leq\lim\lambda_u(t)/t\); the other follows from the
Harnack-type comparison \(M_u(t)\leq\lambda_u(t+c_n)+C\), valid for psh
functions on balls with \(c_n\) depending only on the dimension, which
forces the two asymptotic slopes to coincide.
\end{proof}

\begin{warning}{Do not mix the two Lelong normalizations}
\label{warn:ch51-lelong-normalization}
With \eqref{eq:complex-lelong-number}, the Lelong number of the current
\([D]\) of a prime divisor \(D\) along \(D\) is \(1\), and the Lelong
number of \(c\log|f|^2\) at a point of \(\{f=0\}\) is
\(c\cdot\ord_x(f)\).  Sources normalizing by \(\log|z|\) report twice
these numbers; sources normalizing multiplier ideals by
\(\rme^{-2\varphi}\) compensate by halving the weight.  Both rescalings
must be applied together, or formula
\eqref{eq:complex-lelong-from-ideals} acquires a spurious factor \(2\).
\end{warning}

\subsection{Elementary calculus}
\label{sec:ch51-lelong-calculus}

\begin{proposition}\label{prop:ch51-lelong-calculus}
Let \(\varphi,\psi\) be quasi-psh on \(X\), let \(x\in X\), and let
\(\lambda>0\).
\begin{enumerate}
\item\label{itm:ch51-lc:sum}
      \(\nu(\varphi+\psi,x)=\nu(\varphi,x)+\nu(\psi,x)\).
\item\label{itm:ch51-lc:scale}
      \(\nu(\lambda\varphi,x)=\lambda\nu(\varphi,x)\).
\item\label{itm:ch51-lc:max}
      \(\nu(\max\{\varphi,\psi\},x)=\min\{\nu(\varphi,x),\nu(\psi,x)\}\).
\item\label{itm:ch51-lc:convex}
      \(\nu((1-t)\varphi+t\psi,x)=(1-t)\nu(\varphi,x)+t\nu(\psi,x)\) for
      \(t\in[0,1]\).
\item\label{itm:ch51-lc:bounded}
      If \(\varphi-\psi\) is bounded near \(x\), then
      \(\nu(\varphi,x)=\nu(\psi,x)\); in particular \(\nu(\cdot,x)\)
      descends to ordinary singularity types.
\item\label{itm:ch51-lc:smooth}
      \(\nu(\varphi,x)=0\) whenever \(\varphi\) is bounded near \(x\);
      \(\nu(\varphi,x)<1\) implies that \(\rme^{-\varphi}\) is integrable
      near \(x\); and \(\nu(\varphi,x)>n\) implies that it is not.
\end{enumerate}
\end{proposition}

\begin{proof}
Items \cref{itm:ch51-lc:sum}, \cref{itm:ch51-lc:scale} and
\cref{itm:ch51-lc:convex} follow from
\Cref{prop:ch51-lelong-three}, item \cref{itm:ch51-lelong:slope}, because
\(\lambda_{u+v}=\lambda_u+\lambda_v\) and
\(\lambda_{cu}=c\lambda_u\), and asymptotic slopes are additive and
positively homogeneous.  Item \cref{itm:ch51-lc:bounded} is immediate from
the definition \eqref{eq:complex-lelong-number}.

For item \cref{itm:ch51-lc:max}, write \(\gamma_1=\nu(\varphi,x)\),
\(\gamma_2=\nu(\psi,x)\), and \(\gamma=\nu(\max\{\varphi,\psi\},x)\).
Since \(\max\{\varphi,\psi\}\geq\varphi\), monotonicity gives
\(\gamma\leq\gamma_1\), and likewise \(\gamma\leq\gamma_2\).  Conversely,
for \(\gamma'<\min\{\gamma_1,\gamma_2\}\) both \(\varphi\) and \(\psi\)
are bounded above by \(\gamma'\log|z-x|^2+O(1)\), hence so is their
maximum, whence \(\gamma\geq\gamma'\).

Item \cref{itm:ch51-lc:smooth} is Skoda's theorem \cite{Sko72}: if
\(\nu(\varphi,x)<1\), then \(\rme^{-\varphi}\) is integrable near \(x\);
if \(\nu(\varphi,x)>n\), it is not.  The bounds are sharp at both ends:
\(\varphi=\log|z_1|^2\) has \(\nu=1\) and \(\rme^{-\varphi}=|z_1|^{-2}\)
is not integrable, while
\(\varphi=\log(|z_1|^2+\cdots+|z_n|^2)^n=n\log|z|^2\) has \(\nu=n\) and
\(\rme^{-\varphi}=|z|^{-2n}\) is not integrable, whereas
\(\varphi=\gamma\log|z|^2\) with \(\gamma<n\) gives an integrable weight.
\end{proof}

\begin{remark}[Lelong numbers do not detect \(\preceq\)]
\label{rem:ch51-lelong-not-faithful}
Item \cref{itm:ch51-lc:bounded} says that \(\nu(\cdot,x)\) factors through
\(\sim\).  It does \emph{not} say that the family of all Lelong numbers
determines the ordinary singularity type, and it does not.  In
\Cref{ex:ch51-disk-three-types} the potentials \(u_3\) and \(u_4\) have
all Lelong numbers zero while \(u_3\not\sim u_4\); globally,
\Cref{ex:ch51-cantor} does the same with a full set of vanishing Lelong
numbers and different masses.  The exact assertion that the Lelong numbers
\emph{do} determine is \Cref{thm:complex-strong-openness}, item
\cref{itm:ch51-so:equivalence}: they determine all multiplier ideals, and
nothing more.
\end{remark}

\subsection{Monotone limits}
\label{sec:ch51-lelong-monotone}

The next lemma is the one place in this chapter where a small piece of
convexity does genuine work.  It is used in the proof of
\eqref{eq:complex-p-i-comparison}, in the construction of the
\(\mathcal I\)-envelope, in the levelwise argument of
\Cref{prop:rays-kahler-transition}, and in \Cref{ex:ch51-cantor}.

\begin{lemma}
\label{lem:ch51-lelong-increasing}
Let \(\{\varphi_i\}_{i\in I}\subset\PSH(X,\theta)\) be a family that is
uniformly bounded above and \emph{directed}, in the sense that for
\(i,j\in I\) there is \(k\) with \(\varphi_k\geq\max\{\varphi_i,
\varphi_j\}\).  Put \(\varphi=\uscsup_i\varphi_i\).  Then for every
\(x\in X\)
\[
 \nu(\varphi,x)=\inf_{i\in I}\nu(\varphi_i,x),
\]
and for every prime divisor \(E\) over \(X\)
\[
 \nu(\varphi,E)=\inf_{i\in I}\nu(\varphi_i,E).
\]
\end{lemma}

\begin{proof}
Work in a coordinate ball around \(x\) and replace \(\varphi_i\) by
\(g+\varphi_i\), which is psh; adding the fixed smooth \(g\) changes no
Lelong number.  Write \(u_i\) for the resulting psh functions and
\(u\) for \(g+\varphi\), so that \(u=\uscsup_iu_i\).

\emph{Step 1: the easy inequality.}  Since \(u\geq u_i\), monotonicity of
\eqref{eq:complex-lelong-number} gives
\(\nu(u,x)\leq\nu(u_i,x)\) for each \(i\), hence
\(\nu(u,x)\leq\gamma:=\inf_i\nu(u_i,x)\).

\emph{Step 2: spherical means converge.}  Let \(\lambda_i:=\lambda_{u_i}\)
and \(\lambda:=\lambda_u\) in the notation of
\Cref{sec:ch51-lelong-three}.  By
\Cref{prop:ch51-psh-operations}, item \cref{itm:ch51-op:sup}, the set
\(\{u<\sup_iu_i\}\) is pluripolar, hence Lebesgue-null, so
\(u=\sup_iu_i\) almost everywhere.  Because the family is directed and
uniformly bounded above, monotone convergence for the increasing net of
integrals gives
\[
 \lambda(t)=\sup_i\lambda_i(t)
 \qquad\text{for every }t<t_0 .
\]

\emph{Step 3: the convexity estimate.}  By
\Cref{prop:ch51-lelong-three}, item \cref{itm:ch51-lelong:slope},
\(\lambda_i'\) is increasing with
\(\lim_{t\to-\infty}\lambda_i'(t)=\nu(u_i,x)\geq\gamma\); hence
\(\lambda_i'\geq\gamma\) everywhere.  Fix \(t_1<t_0\) with
\(\lambda(t_1)<+\infty\).  For \(t\leq t_1\),
\[
 \lambda_i(t)
 =\lambda_i(t_1)-\int_t^{t_1}\lambda_i'(s)\,\rd s
 \leq\lambda_i(t_1)-\gamma(t_1-t)
 \leq\lambda(t_1)+\gamma(t-t_1).
\]
Taking the supremum over \(i\) and using Step 2,
\[
 \lambda(t)\leq\lambda(t_1)+\gamma(t-t_1)
 \qquad(t\leq t_1).
\]
Dividing by \(t<0\) and letting \(t\to-\infty\) gives
\(\nu(u,x)=\lim_{t\to-\infty}\lambda(t)/t\geq\gamma\).  With Step 1 this
proves the first identity.

\emph{Step 4: divisors.}  Let \(\pi:Y\to X\) be a modification and
\(E\subset Y\) a prime divisor.  By Siu's theorem the generic Lelong
number satisfies \(\nu(\psi,E)=\inf_{y\in E}\nu(\pi^*\psi,y)\) for every
quasi-psh \(\psi\) on \(X\) whose pullback is not identically \(-\infty\).
The family \(\{\pi^*\varphi_i\}\) is again directed and uniformly bounded
above, and \(\uscsup_i\pi^*\varphi_i=\pi^*\varphi\) because the two sides
are upper semicontinuous, agree outside a pluripolar set, and each
dominates the other's unregularized supremum.  Applying the first identity
at each \(y\in E\) and exchanging the two infima,
\[
 \nu(\varphi,E)
 =\inf_{y\in E}\inf_i\nu(\pi^*\varphi_i,y)
 =\inf_i\inf_{y\in E}\nu(\pi^*\varphi_i,y)
 =\inf_i\nu(\varphi_i,E).
\]
\end{proof}

\begin{warning}{No decreasing analogue}
\label{warn:ch51-lelong-decreasing}
For a decreasing sequence \(\varphi_j\searrow\varphi\) one has
\(\nu(\varphi_j,x)\leq\nu(\varphi,x)\), but the inequality can be strict
in the limit: on the disk, \(\varphi_j=\max\{\log|z|^2,-j\}\) decreases to
\(\log|z|^2\) with \(\nu(\varphi_j,0)=0\) for every \(j\) and
\(\nu(\varphi,0)=1\).  Every use of a decreasing family below therefore
carries its own hypothesis, and none of them is a formal consequence of
\Cref{lem:ch51-lelong-increasing}.  The same asymmetry reappears for
masses in \Cref{prop:complex-mass-continuity} and for the trace in
\Cref{rem:operations-trace-not-increasing}, where it goes in the opposite
direction.
\end{warning}

\subsection{Divisorial Lelong numbers}
\label{sec:ch51-divisorial}

Let \(\pi:Y\to X\) be a modification and let \(E\subset Y\) be a prime
divisor.  The generic value of
\(\nu(\pi^*\varphi,\cdot)\) on \(E\) is denoted
\[
                         \nu(\varphi,E).
\tag{51.8}\label{eq:complex-divisorial-lelong}
\]
It depends only on the divisor over \(X\), not on the chosen realization.

Three points require proof or at least an argument.

\emph{The generic value exists.}  By Siu's semicontinuity theorem
\cite{Siu74}, for each \(c>0\) the set \(\{y\mid\nu(\pi^*\varphi,y)\geq
c\}\) is an analytic subset of \(Y\).  Hence
\(y\mapsto\nu(\pi^*\varphi,y)\) is constant on the complement in \(E\) of
a countable union of proper analytic subsets, and that constant value is
the minimum \(\inf_{y\in E}\nu(\pi^*\varphi,y)\).  This is the identity
used in Step 4 of the proof of \Cref{lem:ch51-lelong-increasing}.

\emph{Independence of the realization.}  If \(\pi':Y'\to X\) is another
modification on which the same divisorial valuation \(\ord_E\) of the
function field is realized by a prime divisor \(E'\), dominate both by a
third modification \(Y''\to Y,Y'\).  The strict transforms of \(E\) and
\(E'\) in \(Y''\) coincide, because they define the same valuation; and
the generic Lelong number of \(\pi''^*\varphi\) along that common divisor
computes both numbers, since Lelong numbers are preserved by pullback
along a biholomorphism between the generic points.  Thus
\eqref{eq:complex-divisorial-lelong} depends only on \(\ord_E\).

\emph{The relation to the Siu decomposition.}  If
\(T=\theta_\varphi\) is a closed positive \((1,1)\)-current on \(X\) and
\(D\subset X\) is a prime divisor, then
\[
 T\geq\nu(\varphi,D)\,[D],
\]
and the difference is again a closed positive current.  Iterating over the
countably many \(D\) with \(\nu(\varphi,D)>0\) gives the Siu
decomposition
\[
 T=\sum_{k}\nu_k[D_k]+R,
\]
with \(R\) closed positive and \(\nu(R,D)=0\) for every prime divisor
\(D\subset X\).  This is \cite{Siu74}, and it is used in
\Cref{ex:ch51-blowup-mass} to compute a minimal singularity type from an
intersection number.

\begin{proposition}\label{prop:ch51-divisorial-calculus}
The assignment \(\varphi\mapsto(\nu(\varphi,E))_E\), indexed by the prime
divisors over \(X\), satisfies:
\begin{enumerate}
\item\label{itm:ch51-dc:calculus}
      additivity, positive homogeneity, and the minimum rule for maxima,
      as in \Cref{prop:ch51-lelong-calculus};
\item\label{itm:ch51-dc:monotone}
      \(\varphi\preceq\psi\Rightarrow\nu(\varphi,E)\geq\nu(\psi,E)\);
\item\label{itm:ch51-dc:increasing}
      the increasing-limit formula of \Cref{lem:ch51-lelong-increasing};
\item\label{itm:ch51-dc:pullback}
      \(\nu(\pi^*\varphi,E)=\nu(\varphi,E)\) for every modification
      \(\pi\) and every prime divisor \(E\) over \(Y\), viewed as a
      divisor over \(X\);
\item\label{itm:ch51-dc:analytic}
      if \(\varphi=c\log|\mathfrak a|^2+O(1)\) locally for a coherent
      ideal \(\mathfrak a\), then
      \(\nu(\varphi,E)=c\,\ord_E(\mathfrak a)\).
\end{enumerate}
\end{proposition}

\begin{proof}
Items \cref{itm:ch51-dc:calculus}--\cref{itm:ch51-dc:increasing} follow from
the corresponding pointwise statements applied at a generic point of
\(E\), using the identity \(\nu(\cdot,E)=\inf_{y\in E}\nu(\cdot,y)\)
established above and the fact that an infimum of sums is not in general a
sum of infima --- which is why the additivity in item
\cref{itm:ch51-dc:calculus} needs the extra remark that the exceptional
loci for \(\varphi\) and \(\psi\) are countable unions of proper analytic
subsets of \(E\), so that a common generic point may be chosen.  Item
\cref{itm:ch51-dc:pullback} is the definition together with the
independence of realization proved above.

For item \cref{itm:ch51-dc:analytic}, let \(\rho:Z\to X\) be a log
resolution on which \(\mathfrak a\cdot\cO_Z=\cO_Z(-F)\) is invertible.
Then \(\rho^*\varphi=c\log|s_F|^2+O(1)\) with \(s_F\) a local equation of
\(F\), so \(\nu(\rho^*\varphi,\cdot)\) equals \(c\) times the coefficient
of \(F\) at the point in question; taking generic points and using item
\cref{itm:ch51-dc:pullback} gives
\(\nu(\varphi,E)=c\,\ord_E(F)=c\,\ord_E(\mathfrak a)\).
\end{proof}

\subsection{A worked example on \texorpdfstring{\(\mathbf P^2\)}{P2}}
\label{sec:ch51-p2-example}

The following potential is used twice: here for its Lelong numbers, and in
\Cref{ex:ch51-p2-ideals} for its multiplier ideals.  It is the simplest
example in which the ideal is not a power of a maximal ideal and the two
singular points are genuinely different.

\begin{example}[A Fubini--Study potential with two base points]
\label{ex:ch51-p2-lelong}
On \(X=\mathbf P^2\) with homogeneous coordinates \([Z_0:Z_1:Z_2]\), let
\(h\) be the Fubini--Study metric on \(\cO(1)\), so that
\(\omega:=c_1(\cO(1),h)\) satisfies \(\int_X\omega^2=1\).  Put
\[
 s_1:=Z_1^3,\qquad s_2:=Z_0Z_2^2\ \in H^0(\mathbf P^2,\cO(3)),
\]
and define
\[
 \varphi:=\log\bigl(|s_1|^2_{h^3}+|s_2|^2_{h^3}\bigr)
 =\log\frac{|Z_1|^6+|Z_0|^2|Z_2|^4}{\|Z\|^6},
 \qquad
 \|Z\|^2=|Z_0|^2+|Z_1|^2+|Z_2|^2 .
\]
Then \(\varphi\in\PSH(X,3\omega)\), because
\(3\omega+\ddc\varphi=\ddc\log(\sum_j|s_j|^2_{h^3})+3\omega\) is the
curvature current of the singular metric \(h^3\rme^{-\varphi}\), and the
latter is the pullback of a positive current under a rational map defined
by holomorphic sections.  Explicitly, \(\varphi\) has analytic
singularities in the sense of \Cref{def:complex-analytic-singularities},
with \(c=1\) and the ideal \(\mathfrak a=(s_1,s_2)\).

The base locus is
\[
 \{s_1=s_2=0\}
 =\{Z_1=0\}\cap\bigl(\{Z_0=0\}\cup\{Z_2=0\}\bigr)
 =\{p_1,p_2\},
\]
where
\[
 p_1=[0:0:1],
 \qquad
 p_2=[1:0:0],
\]
and \(\varphi\) is smooth on \(X\setminus\{p_1,p_2\}\).  In the affine
chart \(Z_0=1\) with \(z_i=Z_i/Z_0\), one has near \(p_2=(0,0)\)
\[
 \varphi=\log\bigl(|z_1|^6+|z_2|^4\bigr)+O(1),
 \qquad
 \mathfrak a_{p_2}=(z_1^3,z_2^2);
\]
in the chart \(Z_2=1\) with \(w_0=Z_0/Z_2,\ w_1=Z_1/Z_2\), one has near
\(p_1=(0,0)\)
\[
 \varphi=\log\bigl(|w_1|^6+|w_0|^2\bigr)+O(1),
 \qquad
 \mathfrak a_{p_1}=(w_1^3,w_0).
\]

\emph{Lelong numbers at points.}  By
\Cref{prop:ch51-divisorial-calculus}, item \cref{itm:ch51-dc:analytic}, or
directly from \eqref{eq:complex-lelong-number},
\[
 \begin{gathered}
 \nu(\varphi,p_2)=\min\{3,2\}=2,
 \qquad
 \nu(\varphi,p_1)=\min\{3,1\}=1,\\
 \nu(\varphi,x)=0\quad\text{for }x\neq p_1,p_2 .
 \end{gathered}
\]
For instance at \(p_2\), along \(z_2=0\) one has
\(\varphi=3\log|z_1|^2+O(1)\) and along \(z_1=0\) one has
\(\varphi=2\log|z_2|^2+O(1)\), while along the diagonal
\(z_1=z_2=\zeta\) one has \(\varphi=\log(|\zeta|^6+|\zeta|^4)+O(1)
=2\log|\zeta|^2+O(1)\) and \(\log|z|^2=\log|\zeta|^2+\log2\); the ratio
tends to \(2\), which is the minimum, hence the Lelong number.

\emph{Divisorial Lelong numbers.}  Let \(v_{a,b}\) with \(a,b>0\) be the
monomial valuation on \(\cO_{\mathbf P^2,p_2}\) with
\(v_{a,b}(z_1)=a\), \(v_{a,b}(z_2)=b\); it is divisorial when \(a/b\) is
rational, realized on a suitable weighted blow-up.  Then
\[
 \nu(\varphi,v_{a,b})=\min\{3a,2b\}.
\]
In particular the exceptional divisor \(E_2\) of the ordinary blow-up of
\(p_2\) is \(v_{1,1}\), and \(\nu(\varphi,E_2)=2\).  Similarly at
\(p_1\), with \(v'_{a,b}(w_0)=a\) and \(v'_{a,b}(w_1)=b\),
\[
 \nu(\varphi,v'_{a,b})=\min\{a,3b\},
 \qquad
 \nu(\varphi,E_1)=1
\]
for the exceptional divisor \(E_1\) of the blow-up of \(p_1\).  Every
other prime divisor over \(\mathbf P^2\) has \(\nu(\varphi,\cdot)=0\)
unless its centre is \(p_1\) or \(p_2\).

\emph{A mass computation, for contrast.}  The two sections \(s_1,s_2\)
define a rational map \(\Phi=[s_1:s_2]:\mathbf P^2\dashrightarrow
\mathbf P^1\), and \(3\omega+\ddc\varphi=\Phi^*\omega_{\mathbf P^1}\) on
\(X\setminus\{p_1,p_2\}\), where \(\omega_{\mathbf P^1}\) is the
Fubini--Study form of total mass \(1\).  A pullback from a curve has
vanishing square, so
\[
 \langle(3\omega+\ddc\varphi)^2\rangle=0,
 \qquad
 M_{3\omega}(\varphi)=0,
 \qquad\text{while}\qquad
 \int_X(3\omega)^2=9 .
\]
Thus a potential with analytic singularities in a K\"ahler class can carry
no nonpluripolar mass at all.  By
\eqref{eq:complex-analytic-envelope-agreement} it is nevertheless a fixed
point of both envelope operators; it is not a \emph{model potential} in
the sense of \Cref{def:complex-p-envelope}, because that term is reserved
for positive mass.  The distinction is exactly the one that
\Cref{chap:transcendental-metrics} is built to respect.
\end{example}

\begin{remark}[Why the mass vanishes, in ideal-theoretic terms]
\label{rem:ch51-p2-multiplicity}
The vanishing above can be read off numerically.  For a potential
\(\log|\mathfrak a|^2\) with \(\mathfrak a\) a base-point ideal of a
linear system in \(|L|\), the nonpluripolar mass equals
\((L^n)\) minus the sum of the local Hilbert--Samuel multiplicities of
\(\mathfrak a\) at its base points, when those are isolated.  Here
\((L^2)=9\), while
\(e(\mathfrak a_{p_2})=3\cdot2=6\) and
\(e(\mathfrak a_{p_1})=3\cdot1=3\), for a total of \(9\).  The two
computations agree.  In \Cref{ex:ch51-blowup-mass} the same bookkeeping is
carried out on a blown-up surface, where the answer is not zero.
\end{remark}

\section{Multiplier ideals}
\label{sec:ch51-multiplier}

A Lelong number is a single real number attached to a point or a divisor.
A multiplier ideal is a coherent sheaf, and it records the singularity of
\(\varphi\) at every point at once, in the way that \(L^2\) methods can
use.  The two are equivalent packages of information, but the equivalence
--- \Cref{thm:complex-strong-openness}, item
\cref{itm:ch51-so:equivalence} --- is a theorem, and it is the reason the
\(\mathcal I\)-preorder of \Cref{def:complex-i-preorder} can be defined by
ideals and computed by valuations.

\subsection{The definition and its normalization}
\label{sec:ch51-multiplier-def}

\begin{definition}[Multiplier ideal]
\label{def:complex-multiplier-ideal}
For a quasi-psh function \(\varphi\), its multiplier ideal is
\[
 \mathcal I(\varphi)(U)
 :=\left\{f\in\mathcal O_X(U)\ \middle|\
 |f|^2\rme^{-\varphi}\in L^1_{\mathrm{loc}}(U)\right\}.
\tag{51.9}\label{eq:complex-multiplier-ideal}
\]
\end{definition}

The exponent is \(\rme^{-\varphi}\).  For example, on a disk,
\[
 \mathcal I\!\left(\lambda\log|z|^2\right)
       =\mathcal O_{\mathbf C}\bigl(-\lfloor\lambda\rfloor\{0\}\bigr).
\tag{51.10}\label{eq:complex-multiplier-example}
\]
Thus the endpoint \(\lambda\in\mathbf Z\) is included in the jump.

Formula \eqref{eq:complex-multiplier-example} is worth verifying by hand,
because it is the single most effective test of the two conventions in
force.  For \(f=z^k\) one has
\[
 \int_{|z|<1}|z|^{2k}|z|^{-2\lambda}\,\rd V
 =2\pi\int_0^1r^{2k-2\lambda+1}\,\rd r,
\]
which is finite exactly when \(2k-2\lambda+1>-1\), that is
\(k>\lambda-1\).  For \(\lambda\notin\mathbf Z\) this means
\(k\geq\lceil\lambda\rceil-1+1=\lfloor\lambda\rfloor\); for
\(\lambda\in\mathbf Z\) it means \(k\geq\lambda=\lfloor\lambda\rfloor\).
In both cases the answer is \((z^{\lfloor\lambda\rfloor})\), which is
\eqref{eq:complex-multiplier-example}.  A source that reports
\((z^{\lceil\lambda\rceil-1})\) is using the strict convention
\(k>\lambda-1\) with a different treatment of the integral endpoint; a
source that reports \((z^{\lfloor2\lambda\rfloor})\) is using the weight
\(\rme^{-2\varphi}\) of \Cref{warn:ch51-half-log}.

\begin{proposition}\label{prop:ch51-ideal-formal}
Let \(\varphi,\psi\) be quasi-psh on \(X\), and \(\lambda,\mu>0\).
\begin{enumerate}
\item\label{itm:ch51-if:coherent}
      \(\mathcal I(\varphi)\) is a coherent ideal sheaf, and
      \(\cO_X/\mathcal I(\varphi)\) is supported on
      \(\{\varphi=-\infty\}\).
\item\label{itm:ch51-if:monotone}
      \(\varphi\leq\psi\Rightarrow
      \mathcal I(\varphi)\subseteq\mathcal I(\psi)\); more generally
      \(\varphi\preceq\psi\Rightarrow
      \mathcal I(\lambda\varphi)\subseteq\mathcal I(\lambda\psi)\) for
      every \(\lambda>0\).
\item\label{itm:ch51-if:scale}
      \(\lambda\leq\mu\Rightarrow
      \mathcal I(\mu\varphi)\subseteq\mathcal I(\lambda\varphi)\), and
      \(\mathcal I(\lambda\varphi)=\cO_X\) for \(\lambda\) small.
\item\label{itm:ch51-if:product}
      Subadditivity holds:
      \(\mathcal I(\varphi+\psi)\subseteq
      \mathcal I(\varphi)\cdot\mathcal I(\psi)\).  The reverse inclusion
      is false.
\item\label{itm:ch51-if:skoda-bounds}
      Skoda's bounds hold:
      \(\nu(\varphi,x)<1\) implies
      \(\mathcal I(\varphi)_x=\cO_{X,x}\), and
      \(\nu(\varphi,x)>n\) implies
      \(\mathcal I(\varphi)_x\subseteq\mathfrak m_x\).
\end{enumerate}
\end{proposition}

\begin{proof}
Item \cref{itm:ch51-if:coherent} is the coherence theorem of Nadel
\cite{Nad89,Nad90}, proved from the H\"ormander--Bombieri--Skoda
\(L^2\)-extension technique; the support statement holds because
\(\rme^{-\varphi}\) is bounded, hence integrable, wherever \(\varphi\) is
bounded below.  Items \cref{itm:ch51-if:monotone} and
\cref{itm:ch51-if:scale} are immediate from the definition: if
\(\varphi\leq\psi\) then \(\rme^{-\psi}\leq\rme^{-\varphi}\); adding a
constant to \(\varphi\) multiplies the weight by a constant and does not
change integrability; and \(\mathcal I(\lambda\varphi)=\cO_X\) for
\(\lambda\) small, because \(\sup_x\nu(\lambda\varphi,x)
=\lambda\sup_x\nu(\varphi,x)<1\) once \(\lambda\) is small, the supremum
being finite by compactness of \(X\) and upper semicontinuity of
\(\nu(\varphi,\cdot)\).

Item \cref{itm:ch51-if:product} is the subadditivity theorem of
Demailly--Ein--Lazarsfeld \cite{DEL00}, proved by restricting to the
diagonal of \(X\times X\) using the Ohsawa--Takegoshi extension theorem
\cite{OT87}.  The reverse inclusion fails already on a disk: with
\(\varphi=\psi=0.6\log|z|^2\) one has
\(\mathcal I(\varphi)=\mathcal I(\psi)=\cO\) by
\eqref{eq:complex-multiplier-example}, while
\(\mathcal I(\varphi+\psi)=\mathcal I(1.2\log|z|^2)=(z)\); thus
\(\mathcal I(\varphi)\cdot\mathcal I(\psi)=\cO\not\subseteq(z)\).

Item \cref{itm:ch51-if:skoda-bounds} is Skoda's theorem \cite{Sko72} in the
normalization \eqref{eq:complex-lelong-number}.  Neither bound can be
improved without further hypotheses: \(\varphi=\log|z_1|^2\) on
\(\mathbf C^2\) has \(\nu(\varphi,0)=1\) and
\(\mathcal I(\varphi)=(z_1)\subsetneq\cO\), while
\(\varphi=\gamma\log|z|^2\) with \(\gamma<n\) has
\(\mathcal I(\varphi)_0=\cO_0\).
\end{proof}

\begin{example}[Log canonical thresholds]
\label{ex:ch51-lct}
For a quasi-psh \(\varphi\) and \(x\in X\), the \emph{complex singularity
exponent} or log canonical threshold at \(x\) is
\[
 \lct_x(\varphi)
 :=\sup\{\lambda>0\mid \mathcal I(\lambda\varphi)_x=\cO_{X,x}\}.
\]
By \Cref{prop:ch51-lelong-calculus}, item \cref{itm:ch51-lc:smooth},
\[
 \frac{1}{\nu(\varphi,x)}\leq\lct_x(\varphi)\leq
 \frac{n}{\nu(\varphi,x)},
\]
with both bounds attained: the left by \(\varphi=\log|z_1|^2\) and the
right by \(\varphi=\log|z|^2\).  For the potential of
\Cref{ex:ch51-p2-lelong},
\[
 \lct_{p_2}(\varphi)=\frac13+\frac12=\frac56,
 \qquad
 \lct_{p_1}(\varphi)=\frac11+\frac13=\frac43 ,
\]
by the monomial formula recalled in \Cref{prop:ch51-howald}.  Both lie in
the predicted intervals \([1/2,1]\) and \([1,2]\).  The
semicontinuity of \(\lct\) in families is \cite{DK01}; its
effective refinement is \cite{Hiep14}.  Neither is used below, but both
are the natural strengthenings of \Cref{thm:complex-strong-openness}.
\end{example}

\subsection{Strong openness and the valuative criterion}
\label{sec:ch51-openness}

\refstepcounter{theorem}\label{thm:complex-strong-openness}
\begin{quote}
\textbf{Foundational input \thetheorem\ (strong openness and its
valuative package; source boundary).}
Let \(\varphi\) be quasi-psh.
\begin{enumerate}
\item\label{itm:ch51-so:openness} The strong openness identity holds:
      \[
      \mathcal I(\varphi)
       =\bigcup_{\epsilon>0}\mathcal I((1+\epsilon)\varphi).
      \tag{51.11}\label{eq:complex-strong-openness}
      \]
\item\label{itm:ch51-so:lelong} For every prime divisor \(E\) over \(X\),
      \[
      \nu(\varphi,E)
       =\lim_{m\to\infty}\frac1m
                    \ord_E\mathcal I(m\varphi).
      \tag{51.12}\label{eq:complex-lelong-from-ideals}
      \]
\item\label{itm:ch51-so:equivalence} For quasi-psh \(\varphi,\psi\), equality
      \(\mathcal I(m\varphi)=\mathcal I(m\psi)\) for all
      \(m\geq1\) is equivalent to equality
      \(\mathcal I(\lambda\varphi)=\mathcal I(\lambda\psi)\) for
      every \(\lambda>0\), and to
      \(\nu(\varphi,E)=\nu(\psi,E)\) for every prime divisor over \(X\).
\end{enumerate}
\par
Item \cref{itm:ch51-so:openness} is imported from
\cite[Theorem~1.1]{GZ15}; item \cref{itm:ch51-so:lelong} uses Demailly's
approximation theorem \cite{Dem92reg}; and item
\cref{itm:ch51-so:equivalence} uses Boucksom's valuative criterion
\cite[Corollary~10.18]{Bou17}.  These analytic inputs are not reproved
here.  The deductions and normalization checks are recorded below.
\end{quote}

Before the proof we isolate the imported valuative statement, in the exact
form in which it is used, and check its normalization.

\refstepcounter{proposition}\label{prop:ch51-valuative-criterion}
\begin{quote}
\textbf{Foundational input \theproposition\ (valuative criterion for
multiplier ideals; source boundary).}
Let \(\varphi\) be quasi-psh on \(X\), let \(\lambda>0\), let \(U\subseteq
X\) be open, and let \(f\in\cO_X(U)\).  For a prime divisor \(E\) over
\(X\), write \(A(E)\) for its log discrepancy, that is,
\(A(E)=1+\ord_E(K_{Y/X})\) for any modification \(\pi:Y\to X\) on which
\(E\) appears.  Then
\[
 f\in\mathcal I(\lambda\varphi)(U)
 \quad\Longleftrightarrow\quad
 \begin{aligned}[t]
 &\text{there is }\eta>0\text{ such that}\\
 &\ord_E(f)\geq(1+\eta)\lambda\,\nu(\varphi,E)-A(E)\\
 &\text{for every prime divisor }E\text{ over }U .
 \end{aligned}
\]
In particular \(\mathcal I(\lambda\varphi)\) is determined by the family
\((\nu(\varphi,E))_E\).
This equivalence is imported from \cite[Corollary~10.18]{Bou17}; no proof
is claimed here.  We only check its normalization.

This is \cite[Corollary~10.18]{Bou17}.  That source normalizes Lelong
numbers by \(\log|z|\) and multiplier ideals by the weight
\(\rme^{-2\varphi}\); the substitution
\(\varphi\rightsquigarrow\lambda\varphi/2\) converts its statement into the
displayed one, both rescalings being applied simultaneously as required by
\Cref{warn:ch51-half-log}.  Two unit tests confirm the result.  First, with
\(\varphi=\log|z|^2\) on the unit disk, the only relevant divisor is
\(E=\{0\}\) with \(A(E)=1\) and \(\nu(\varphi,E)=1\); the criterion reads
\(\ord_0(f)\geq(1+\eta)\lambda-1\) for some \(\eta>0\), that is
\(\ord_0(f)>\lambda-1\), which reproduces
\eqref{eq:complex-multiplier-example} including the endpoint behaviour at
\(\lambda\in\mathbf Z\).  Second, for \(f=1\) the criterion says
\(\mathcal I(\lambda\varphi)=\cO_X\) exactly when
\(\lambda<\inf_EA(E)/\nu(\varphi,E)\), which is the standard valuative
formula for the log canonical threshold of \Cref{ex:ch51-lct}.  The single
\(\eta\), uniform over all \(E\), is essential; replacing it by pointwise
strict inequalities gives a strictly larger set of \(f\).
\end{quote}

\paragraph{Deductions and normalization checks.}
\emph{Item \cref{itm:ch51-so:openness}.}  The inclusion \(\supseteq\) is
\Cref{prop:ch51-ideal-formal}, item \cref{itm:ch51-if:scale}.  The
inclusion \(\subseteq\) is the strong openness theorem of Guan--Zhou
\cite[Theorem~1.1]{GZ15}.  In the normalization
\eqref{eq:complex-multiplier-ideal} the statement is exactly as displayed;
sources using the weight \(\rme^{-2\varphi}\) state it for
\(\mathcal I(2\varphi)\) and the two forms are equivalent under
\(\varphi\rightsquigarrow\varphi/2\).

\emph{Item \cref{itm:ch51-so:lelong}.}  Fix \(m\geq1\) and let
\(\{\sigma_{m,j}\}_j\) be a family of local generators of
\(\mathcal I(m\varphi)\) obtained from an orthonormal basis of the
Bergman space of \(L^2\) holomorphic functions with weight
\(\rme^{-m\varphi}\); the generation statement is the
H\"ormander--Bombieri--Skoda argument used in \cite{Dem92reg}.  Set
\[
 \varphi_m:=\frac1m\log\sum_j|\sigma_{m,j}|^2 ,
\]
which has analytic singularities with ideal \(\mathcal I(m\varphi)\) and
coefficient \(1/m\).  By \Cref{prop:ch51-divisorial-calculus}, item
\cref{itm:ch51-dc:analytic},
\[
 \nu(\varphi_m,E)=\frac1m\ord_E\mathcal I(m\varphi)
\]
for every prime divisor \(E\) over \(X\).  Demailly's approximation
theorem \cite{Dem92reg} gives, in our normalization,
\[
 \nu(\varphi,E)-\frac nm\leq\nu(\varphi_m,E)\leq\nu(\varphi,E),
\]
the right inequality because \(\varphi_m\geq\varphi-C_m/m\) up to a
constant with \(C_m/m\to0\), hence \(\varphi\preceq\varphi_m\) up to that
error, and the left inequality by the \(L^2\) estimate producing a section
of controlled vanishing.  Letting \(m\to\infty\) gives
\eqref{eq:complex-lelong-from-ideals}.

An alternative derivation of the inequality
\(\liminf_m m^{-1}\ord_E\mathcal I(m\varphi)\geq\nu(\varphi,E)\), which
uses no approximation theorem, is to apply
\Cref{prop:ch51-valuative-criterion} with \(\lambda=m\): every
\(f\in\mathcal I(m\varphi)\) satisfies
\(\ord_E(f)\geq m\nu(\varphi,E)-A(E)\), and \(A(E)\) does not depend on
\(m\).

\emph{Item \cref{itm:ch51-so:equivalence}.}  We prove a cycle of
implications.  Suppose \(\mathcal I(m\varphi)=\mathcal I(m\psi)\) for
every integer \(m\geq1\).  Applying
\eqref{eq:complex-lelong-from-ideals} to both sides gives
\(\nu(\varphi,E)=\nu(\psi,E)\) for every prime divisor \(E\) over \(X\).
Suppose next that all these divisorial numbers agree.  Then
\Cref{prop:ch51-valuative-criterion} produces the same membership
criterion for \(\mathcal I(\lambda\varphi)\) and
\(\mathcal I(\lambda\psi)\) at every real \(\lambda>0\), so the two
sheaves coincide.  Finally, equality at every real \(\lambda>0\)
trivially implies equality at every integer \(m\geq1\).  This closes the
cycle.

\begin{warning}{What item \cref{itm:ch51-so:equivalence} does not say}
\label{warn:ch51-so-scope}
The theorem identifies two packages of information --- all multiplier
ideals, and all divisorial Lelong numbers --- with each other.  It does
\emph{not} identify either of them with the ordinary singularity type.  A
potential can have all divisorial Lelong numbers equal to zero and still
be unbounded, as in \Cref{ex:ch51-disk-three-types} and, globally, in
\Cref{ex:ch51-cantor}; log--log poles are the standard local source of
such behaviour.  The whole difficulty of \Cref{sec:ch51-p-vs-i} is that
the missing information is not only ``\(O(1)\)-sized'': it can change a
nonpluripolar mass.
\end{warning}

\subsection{Monomial weights}
\label{sec:ch51-howald}

The multiplier ideals of the toric weights that occur in
\cref{part:complex} can be written down.  Let
\(\mathfrak a\subset\cO_{\mathbf C^n,0}\) be a monomial ideal, and let
\[
 \operatorname{Newt}(\mathfrak a)
 \subset\mathbf R^n_{\geq0}
\]
be the convex hull of the union of \(k+\mathbf R^n_{\geq0}\) over the
exponents \(k\) of monomials in \(\mathfrak a\).

\begin{proposition}\label{prop:ch51-howald}
Let \(\mathfrak a\) be a monomial ideal and let \(c>0\).  With
\(\varphi=c\log|\mathfrak a|^2\), the multiplier ideal
\(\mathcal I(\varphi)\) is the monomial ideal generated by those \(z^k\)
with
\[
 k+\mathbf 1\in\Int\bigl(c\operatorname{Newt}(\mathfrak a)\bigr),
 \qquad \mathbf 1=(1,\ldots,1).
\]
For \(\mathfrak a=(z_1^{a_1},\ldots,z_n^{a_n})\) this reads
\[
 z^k\in\mathcal I(c\log|\mathfrak a|^2)
 \quad\Longleftrightarrow\quad
 \sum_{i=1}^n\frac{k_i+1}{a_i}>c ,
\]
and \(\lct_0(\log|\mathfrak a|^2)=\sum_i1/a_i\).  Moreover the divisorial
Lelong numbers are given by the monomial valuations: for \(v_a\) with
\(v_a(z_i)=a_i>0\),
\[
 \nu\bigl(c\log|\mathfrak a|^2,v_a\bigr)
 =c\min\{\,v_a(z^k)\mid z^k\in\mathfrak a\,\}.
\]
\end{proposition}

\begin{proof}[Imported, with the normalization checked]
Howald's theorem computes the multiplier ideal of a monomial ideal; the
algebraic statement is in \cite{PAG2}, and the plurisubharmonic version
used here, valid for arbitrary toric weights and not only for monomial
ideals, is \cite{Guenancia12}.  In the normalization
\eqref{eq:complex-multiplier-ideal} the weight is
\(\rme^{-\varphi}=(\sum|z^{k}|^2)^{-c}\), so the exponent \(c\) is the
same on both sides and no factor \(2\) appears; a source normalizing by
\(\rme^{-2\varphi}\) states the criterion with \(2c\).

The one-variable case is a direct check: with
\(\mathfrak a=(z^{a})\) the condition reads \((k+1)/a>c\), that is
\(k>ca-1\), which agrees with
\eqref{eq:complex-multiplier-example} for \(\lambda=ca\).

The last formula is \Cref{prop:ch51-divisorial-calculus}, item
\cref{itm:ch51-dc:analytic}, applied to the monomial valuation \(v_a\),
whose value on a monomial ideal is the minimum of its values on the
generators.
\end{proof}

\begin{example}[The \(\mathbf P^2\) potential, continued]
\label{ex:ch51-p2-ideals}
Return to \(\varphi\in\PSH(\mathbf P^2,3\omega)\) of
\Cref{ex:ch51-p2-lelong}, with local ideals
\(\mathfrak a_{p_2}=(z_1^3,z_2^2)\) and
\(\mathfrak a_{p_1}=(w_1^3,w_0)\).  Since \(\varphi\) is smooth away from
\(p_1,p_2\), the sheaf \(\mathcal I(\lambda\varphi)\) equals \(\cO_X\)
outside these two points, and \Cref{prop:ch51-howald} computes its stalks.

\emph{At \(p_2\).}  The condition is
\(\frac{k_1+1}{3}+\frac{k_2+1}{2}>\lambda\), equivalently
\(2k_1+3k_2>6\lambda-5\).
\begin{itemize}
\item \(\lambda=\tfrac56\): the condition becomes \(2k_1+3k_2>0\), so
      \(\mathcal I=(z_1,z_2)=\mathfrak m_{p_2}\); at any smaller
      \(\lambda\) it is \(\cO\).  Hence
      \(\lct_{p_2}(\varphi)=\tfrac13+\tfrac12=\tfrac56\).
\item \(\lambda=1\): \(2k_1+3k_2>1\), so \(k=(0,0)\) is excluded and
      \(k=(1,0),(0,1)\) are allowed:
      \(\mathcal I(\varphi)_{p_2}=(z_1,z_2)\).
\item \(\lambda=2\): \(2k_1+3k_2>7\), that is \(2k_1+3k_2\geq8\).  The
      minimal exponents are \((4,0),(3,1),(1,2),(0,3)\), so
      \[
       \mathcal I(2\varphi)_{p_2}
       =(z_1^4,\;z_1^3z_2,\;z_1z_2^2,\;z_2^3).
      \]
\item \(\lambda=3\): \(2k_1+3k_2>13\), that is \(2k_1+3k_2\geq14\); the
      minimal exponents are
      \[
       (7,0),\ (6,1),\ (4,2),\ (3,3),\ (1,4),\ (0,5).
      \]
\end{itemize}

\emph{At \(p_1\).}  The condition is
\(\frac{k_1+1}{1}+\frac{k_2+1}{3}>\lambda\) in the coordinates
\((w_0,w_1)\), equivalently \(3k_1+k_2>3\lambda-4\).
\begin{itemize}
\item \(\lambda=\tfrac43\): the threshold, so
      \(\lct_{p_1}(\varphi)=1+\tfrac13=\tfrac43\).
\item \(\lambda=1\): \(3k_1+k_2>-1\), so
      \(\mathcal I(\varphi)_{p_1}=\cO\).  The point \(p_1\) is invisible
      to \(\mathcal I(\varphi)\) and becomes visible only at
      \(\lambda>4/3\).
\item \(\lambda=2\): \(3k_1+k_2>2\), that is \(\geq3\), with minimal
      exponents \((1,0),(0,3)\):
      \(\mathcal I(2\varphi)_{p_1}=(w_0,w_1^3)\).
\end{itemize}

\emph{Consistency with \eqref{eq:complex-lelong-from-ideals}.}  Let
\(E_2\) be the exceptional divisor of the blow-up of \(p_2\), so that
\(\ord_{E_2}(z^k)=k_1+k_2\).  Then
\(\ord_{E_2}\mathcal I(\lambda\varphi)\) is the minimum of \(k_1+k_2\)
over the monomials of the ideal, that is subject to
\(2k_1+3k_2>6\lambda-5\).  Since \(k_2\) has the larger coefficient, the
minimum is attained on the \(k_2\)-axis, and
\[
 \ord_{E_2}\mathcal I(\lambda\varphi)
 =\Bigl\lfloor2\lambda-\tfrac53\Bigr\rfloor+1 .
\]
For \(\lambda=2\) this gives \(3\), realized by \(z_2^3\) and also by
\(z_1z_2^2\), in agreement with the generators listed above.  Dividing by
\(\lambda\) and letting \(\lambda\to\infty\) gives \(2=\nu(\varphi,E_2)\),
as computed in \Cref{ex:ch51-p2-lelong}.  At \(p_1\) the constraint is
\(3k_1+k_2>3\lambda-4\), the minimum is attained on the \(k_1\)-axis, and
\[
 \ord_{E_1}\mathcal I(\lambda\varphi)
 =\Bigl\lfloor\lambda-\tfrac43\Bigr\rfloor+1 ,
\]
which is \(1\) at \(\lambda=2\) and has normalized limit
\(\nu(\varphi,E_1)=1\).
\end{example}

\begin{warning}{Which variable carries the minimum}
\label{warn:ch51-ord-monomial}
The order of vanishing of a monomial ideal along a divisor is the minimum
over \emph{all} its monomials, not over a chosen generating set, and the
minimizing monomial is the one loading the variable with the largest
coefficient in the defining inequality.  Loading the wrong variable in the
computation above returns \(3\) instead of \(2\) for
\(\nu(\varphi,E_2)\), contradicting \Cref{ex:ch51-p2-lelong}.  This is the
fastest available check on a multiplier-ideal computation, and it is one
reason \eqref{eq:complex-lelong-from-ideals} is stated as a limit rather
than as an equality at a single scale.
\end{warning}

\begin{remark}[Reading the whole \(\mathcal I\)-type off the Newton
polygon]
\label{rem:ch51-newton}
For an analytic singularity with monomial ideal the entire
\(\mathcal I\)-equivalence class is encoded by
\(\operatorname{Newt}(\mathfrak a)\), and the divisorial Lelong numbers
are the support function of that Newton polyhedron evaluated at monomial
valuations.  Two monomial ideals with the same integral closure --- hence
the same Newton polyhedron --- give \(\mathcal I\)-equivalent, indeed
\(\preceq\)-equivalent, potentials.  This is the local shadow of the
tropical picture of \Cref{chap:tropical-spectra-complex} and of the toric
computations in \Cref{chap:transcendental-okounkov-dh}.
\end{remark}

\subsection{Analytic singularities and quasi-equisingular approximation}
\label{sec:ch51-analytic-sing}

Analytic singularities are the singularities for which all four
comparisons of \Cref{sec:ch51-four} collapse into one.  They are also the
only singularities that a resolution of singularities can see, and they
are the target of every regularization theorem in the subject.

\begin{definition}[Analytic singularities]
\label{def:complex-analytic-singularities}
A quasi-psh function has analytic singularities if locally
\[
 \varphi=c\log\!\left(\sum_{j=1}^N|f_j|^2\right)+O(1),
 \qquad c>0,
\tag{51.13}\label{eq:complex-analytic-singularities}
\]
for holomorphic \(f_j\).  A log resolution converts this to
\(c\log|s_D|^2+O(1)\) for an effective divisor \(D\).
\end{definition}

\begin{proposition}
\label{prop:ch51-analytic-rigid}
Let \(\varphi,\psi\) be quasi-psh on \(X\) with analytic singularities.
Then
\[
 \varphi\preceq\psi
 \quad\Longleftrightarrow\quad
 \varphi\preceq_{\mathcal I}\psi ,
\]
and consequently \(\varphi\sim\psi\iff\varphi\sim_{\mathcal I}\psi\).
\end{proposition}

\begin{proof}
The implication \(\Rightarrow\) is
\Cref{prop:ch51-ideal-formal}, item \cref{itm:ch51-if:monotone}.  For the
converse, choose by Hironaka a single modification
\(\rho:Z\to X\) that is simultaneously a log resolution for the ideals
defining both singularities, so that
\[
 \rho^*\varphi=c\log|s_D|^2+O(1),
 \qquad
 \rho^*\psi=c'\log|s_{D'}|^2+O(1)
\]
with \(D,D'\) effective divisors with simple normal crossing support.
Assume \(\varphi\preceq_{\mathcal I}\psi\); by
\Cref{thm:complex-strong-openness}, item \cref{itm:ch51-so:equivalence},
this gives \(\nu(\varphi,E)\geq\nu(\psi,E)\) for every prime divisor
\(E\) over \(X\), hence in particular
\(c\,\ord_E(D)\geq c'\,\ord_E(D')\) for every prime component \(E\) of
\(\Supp D\cup\Supp D'\) on \(Z\), by
\Cref{prop:ch51-divisorial-calculus}, item \cref{itm:ch51-dc:analytic}.
Therefore \(cD\geq c'D'\) as divisors, so
\(c\log|s_D|^2\leq c'\log|s_{D'}|^2+O(1)\) on \(Z\), that is
\(\rho^*\varphi\preceq\rho^*\psi\).  Pushing forward, and using that
\(\rho\) is an isomorphism over a dense open set while both potentials
are bounded above on the compact \(X\), gives \(\varphi\preceq\psi\).
\end{proof}

\begin{remark}[The two rigidity statements are different]
\label{rem:ch51-rigidity}
\Cref{prop:ch51-analytic-rigid} compares two potentials that both have
analytic singularities.  The last assertion of
\Cref{thm:complex-i-projection}, formula
\eqref{eq:complex-analytic-envelope-agreement}, compares one such
potential with its two envelopes, which are \emph{not} known in advance to
have analytic singularities; the extra input there is envelope maximality,
not resolution of singularities.  Confusing the two produces the false
principle that equality of all divisorial Lelong numbers implies bounded
difference, which \Cref{ex:ch51-disk-three-types} refutes.
\end{remark}

\begin{proposition}\label{prop:complex-quasi-equisingular}
Given \(\varphi\in\PSH(X,\theta)\) and a K\"ahler form \(\omega\), there
are \(\epsilon_j\searrow0\) and decreasing potentials
\[
 \varphi_j\in\PSH(X,\theta+\epsilon_j\omega)
\]
with analytic singularities such that \(\varphi_j\searrow\varphi\) and,
for every \(\lambda'>\lambda>0\),
\[
 \mathcal I(\lambda'\varphi_j)
          \subseteq\mathcal I(\lambda\varphi)
 \quad\text{for }j\gg1.
\tag{51.14}\label{eq:complex-quasi-equisingular}
\]
They may be chosen of gentle type
\((2^{-j},\mathcal I(2^j\varphi))\).
\end{proposition}

\begin{proof}[Imported regularization theorem]
This is the quasi-equisingular form of Demailly regularization.  The
Bergman-kernel approximant at scale \(2^j\) has analytic singularities
of type \((2^{-j},\mathcal I(2^j\varphi))\); the standard diagonal
extraction makes the approximants decreasing while their curvature
losses \(\epsilon_j\omega\) tend to zero.  Strong openness gives the
displayed multiplier-ideal inclusion.  See
\cite{Dem92reg,DPS01} for the regularization theorem and its
quasi-equisingular refinement.
\end{proof}

\begin{definition}[Quasi-equisingular approximation]
\label{def:ch51-quasi-equisingular}
A sequence \((\varphi_j)\) as in
\Cref{prop:complex-quasi-equisingular} --- decreasing, with analytic
singularities, with curvature loss \(\epsilon_j\omega\downarrow0\), and
satisfying \eqref{eq:complex-quasi-equisingular} --- is a
\emph{quasi-equisingular approximation} of \(\varphi\).
\end{definition}

\begin{proposition}\label{prop:ch51-qes-preserves-i}
Fix \(\varphi\in\PSH(X,\theta)\), and let \((\varphi_j)\) be one of its
quasi-equisingular approximations.  For every prime divisor \(E\) over \(X\),
\[
 \nu(\varphi,E)=\lim_{j\to\infty}\nu(\varphi_j,E),
\]
the limit being increasing in \(j\); and for every \(\lambda>0\),
\[
 \mathcal I(\lambda\varphi)
 =\bigcup_{\lambda'<\lambda}\ \bigcap_{j}\mathcal I(\lambda'\varphi_j).
\]
In particular a quasi-equisingular approximation determines, and is
determined up to \(\mathcal I\)-equivalence by, the
\(\mathcal I\)-equivalence class of \(\varphi\); it does \emph{not}
determine the ordinary singularity type.
\end{proposition}

\begin{proof}
Since \(\varphi_j\searrow\varphi\), we have \(\varphi\leq\varphi_j\),
hence \(\nu(\varphi,E)\geq\nu(\varphi_j,E)\), and the sequence
\(\nu(\varphi_j,E)\) is increasing.  For the reverse inequality apply
\eqref{eq:complex-quasi-equisingular}: for \(\lambda'>\lambda>0\) and
\(j\) large, \(\mathcal I(\lambda'\varphi_j)\subseteq
\mathcal I(\lambda\varphi)\).  Taking \(\ord_E\) and using
\Cref{prop:ch51-divisorial-calculus}, item \cref{itm:ch51-dc:analytic},
together with \eqref{eq:complex-lelong-from-ideals} applied to
\(\varphi\), gives
\(\lambda'\nu(\varphi_j,E)\geq\lambda\nu(\varphi,E)-A(E)\) up to a bounded
error independent of \(\lambda\); dividing by \(\lambda\) and letting
\(\lambda\to\infty\) then \(\lambda'/\lambda\to1\) gives
\(\liminf_j\nu(\varphi_j,E)\geq\nu(\varphi,E)\).

The ideal formula follows from
\Cref{prop:ch51-valuative-criterion} applied to \(\varphi\) and to each
\(\varphi_j\), because both sides are described by the same divisorial
inequalities in the limit; strong openness
\eqref{eq:complex-strong-openness} converts the union over
\(\lambda'<\lambda\) into an equality rather than an inclusion.

The final sentence is the point of the definition: the approximants have
analytic singularities, so by
\Cref{prop:ch51-analytic-rigid} their ordinary types are determined by
their divisorial data, and the limit retains only the
\(\mathcal I\)-class.  \Cref{ex:ch51-cantor} shows that this loss is
genuine: the potential there admits the constant approximation
\(\varphi_j\equiv0\).
\end{proof}

\begin{warning}{Quasi-equisingular approximants are not
\texorpdfstring{\(\preceq\)}{order}-faithful}
\label{warn:ch51-qes}
Two potentials with the same quasi-equisingular approximation need not
have the same ordinary singularity type, and need not have the same
nonpluripolar mass.  The trace operator of
\Cref{chap:operations-trace} is also defined through quasi-equisingular
approximants, but its restriction theorem supplies extra information:
restricted approximants have a common \(P\)-class, and their limits are
\(\mathcal I\)-good.  Thus trace depends only on the input
\(\mathcal I\)-class and takes values in output \(P\)-classes; see
\Cref{thm:operations-trace-well-defined,thm:operations-trace-p-descent}
and the factorization \eqref{eq:verified-trace-factorization}.
\end{warning}

\section{Nonpluripolar products and masses}
\label{sec:ch51-npp}

The Monge--Amp\`ere operator is not continuous, and it is not defined at
all on arbitrary psh functions: the product \((\ddc u)^n\) of a psh
function with a strong enough pole is not a measure.  The nonpluripolar
product is the maximal extension that remains a positive closed current
and that ignores exactly the sets a psh function is allowed to blow up on.
Everything downstream --- masses, volumes, energies, test-curve
Duistermaat--Heckman measures --- is built from it, and every failure of
continuity in \cref{part:complex} traces back to the truncation in its
construction.

\subsection{The local construction}
\label{sec:ch51-npp-local}

For psh functions \(u_1,\ldots,u_p\) on a coordinate open set, put
\[
 O_k=\bigcap_{j=1}^p\{u_j>-k\}.
\]
On \(O_k\), Bedford--Taylor theory defines
\[
 \one_{O_k}\ddc u_1\wedge\cdots\wedge\ddc u_p
 =
 \one_{O_k}\bigwedge_{j=1}^p\ddc\max\{u_j,-k\}.
\tag{51.15}\label{eq:complex-npp-truncation}
\]
The right side increases with \(k\).  Extending it by zero across the
pluripolar complement and taking the limit gives
\[
              \left\langle\ddc u_1\wedge\cdots\wedge\ddc u_p\right\rangle .
\]
For global closed positive currents \(T_j=\theta_j+\ddc\varphi_j\) on
a compact K\"ahler manifold, the resulting product
\[
       \langle T_1\wedge\cdots\wedge T_p\rangle
\tag{51.16}\label{eq:complex-global-npp}
\]
is always well-defined.

Three points in this construction deserve to be spelled out, because each
is a place where a later proof either succeeds or fails.

\emph{Why the truncations are compatible.}  Identity
\eqref{eq:complex-npp-truncation} is plurifine locality for the
Bedford--Taylor product: on the plurifine open set \(O_k\), the functions
\(u_j\) and \(\max\{u_j,-k\}\) coincide, and the Bedford--Taylor product
of locally bounded psh functions depends only on their restriction to a
plurifine open set \cite{BT87}.  Consequently the measures
\(\one_{O_k}\bigwedge_j\ddc\max\{u_j,-k\}\) form an increasing family in
\(k\), and their increasing limit exists as a positive current.

\emph{Why the limit is locally finite.}  On a compact K\"ahler manifold
with a K\"ahler form \(\omega\), Stokes' theorem applied to the truncated
potentials gives, for every \(k\),
\[
 \begin{aligned}
 \int_X\one_{O_k}\bigwedge_{j=1}^p\ddc\max\{u_j,-k\}\wedge\omega^{n-p}
 &\leq
 \int_X\bigwedge_{j=1}^p
     \bigl(\theta_j+\ddc\max\{\varphi_j,-k\}\bigr)\wedge\omega^{n-p}\\
 &=\{\theta_1\}\cdots\{\theta_p\}\cdot\{\omega\}^{n-p},
 \end{aligned}
\]
a bound independent of \(k\).  Hence the increasing limit has finite
total mass, and the local statement follows by restriction.  This is the
only place where compactness is used; on a noncompact space the
nonpluripolar product may fail to be locally finite, and \emph{is} defined
only when it is, which is why \eqref{eq:complex-global-npp} is asserted
for compact \(X\).

\emph{Why the answer does not depend on the exhaustion.}  Any increasing
sequence of plurifine open sets whose union has pluripolar complement in
\(\bigcap_j\{u_j>-\infty\}\) gives the same limit, again by plurifine
locality.  In particular the definition does not depend on the choice of
local potentials \(u_j\) in \eqref{eq:complex-npp-truncation}.

\subsection{Bidegree one}
\label{sec:ch51-npp-bidegree-one}

In bidegree \((1,1)\) the nonpluripolar product has a completely explicit
description, and every mass computation in this chapter uses it.

\begin{lemma}\label{lem:ch51-npp-bidegree-one}
Let \(\varphi\in\PSH(X,\theta)\).  Then
\[
 \langle\theta_\varphi\rangle
 =\one_{\{\varphi>-\infty\}}\,\theta_\varphi
\]
as measures on \(X\) when \(n=1\), and as currents in general.  In
particular \(\langle\theta_\varphi\rangle\leq\theta_\varphi\), the two
agree if and only if \(\theta_\varphi\) puts no mass on the pluripolar
set \(\{\varphi=-\infty\}\), and, for \(n=1\),
\[
 M_\theta(\varphi)
 =\int_X\theta_\varphi-\theta_\varphi(\{\varphi=-\infty\})
 =\int_X\{\theta\}-\theta_\varphi(\{\varphi=-\infty\}).
\]
\end{lemma}

\begin{proof}
Take \(p=1\) in \eqref{eq:complex-npp-truncation}.  With
\(O_k=\{\varphi>-k\}\) and \(\varphi_k:=\max\{\varphi,-k\}\), plurifine
locality gives
\[
 \one_{O_k}\theta_\varphi=\one_{O_k}\theta_{\varphi_k}
\]
for every \(k\), since \(\varphi=\varphi_k\) on \(O_k\).  The sets
\(O_k\) increase to \(\{\varphi>-\infty\}\), so the increasing limit of
the left-hand sides is \(\one_{\{\varphi>-\infty\}}\theta_\varphi\), while
the increasing limit of the right-hand sides is by definition
\(\langle\theta_\varphi\rangle\).  The inequality
\(\langle\theta_\varphi\rangle\leq\theta_\varphi\) and the criterion for
equality are immediate, and the mass formula for \(n=1\) follows because
then \(\theta_\varphi\) is a positive measure of total mass
\(\int_X\{\theta\}\), by Stokes.
\end{proof}

\begin{remark}[The higher-degree analogue fails]
\label{rem:ch51-npp-no-restriction}
For \(p\geq2\) there is no formula
\(\langle\theta_{\varphi_1}\wedge\cdots\wedge\theta_{\varphi_p}\rangle
=\one_{\{\text{all }\varphi_j>-\infty\}}
\theta_{\varphi_1}\wedge\cdots\wedge\theta_{\varphi_p}\), because the
right-hand side has no meaning: the unrestricted product of unbounded
potentials is not defined.  The nonpluripolar product is precisely the
substitute, and \Cref{lem:ch51-npp-bidegree-one} is the one degree in
which the substitute can be described by a restriction.  This is why the
computations in \Cref{ex:ch51-cantor} take place on a curve and the
computations in \Cref{ex:ch51-blowup-mass} on a surface use a Siu
decomposition instead.
\end{remark}

\subsection{The ledger}
\label{sec:ch51-npp-ledger}

\begin{theorem}\label{thm:complex-npp-ledger}
The product \eqref{eq:complex-global-npp} is a closed positive
\((p,p)\)-current.  It is symmetric and multilinear in the currents,
local for the plurifine topology, agrees with the Bedford--Taylor
product for locally bounded potentials, and charges no pluripolar set.
If \(\varphi_j\) is replaced by a less singular potential in the same
class, every mixed total mass can only increase.  Proper bimeromorphic
pullback preserves total masses.
\end{theorem}

\begin{proof}[Why truncation is sufficient]
Plurifine locality makes the truncations in
\eqref{eq:complex-npp-truncation} compatible.  Their masses on relatively
compact sets are uniformly controlled after wedging with a K\"ahler
form, so the increasing limit is locally finite.  Bedford--Taylor
integration by parts yields closedness and symmetry before passage to
the limit.  Monotonicity is the global theorem of
\cite{BEGZ10}; bimeromorphic invariance follows by pulling the
truncations to the common isomorphism locus, since the omitted sets are
pluripolar.
\end{proof}

The argument deserves to be run through once more, slowly, because five
distinct assertions are packed into the statement and they have five
different proofs.

\emph{Closedness, symmetry, multilinearity.}  Each truncated product
\(\bigwedge_j\ddc\max\{u_j,-k\}\) is a closed positive current, symmetric
and multilinear in its arguments, by Bedford--Taylor theory
\cite{BT76,BT82}.  These properties are stable under the increasing limit
of the restrictions \(\one_{O_k}(\cdot)\): positivity and symmetry are
pointwise, multilinearity holds at each level \(k\), and closedness
passes to the limit because a limit of closed currents in the weak
topology is closed.

\emph{Plurifine locality.}  If \(u_j=v_j\) on a plurifine open set \(O\)
for every \(j\), then the truncations coincide on
\(O\cap O_k\), hence so do the limits on \(O\).  This is the property that
makes the whole construction work, and it is
\cite{BT87}.

\emph{Agreement with Bedford--Taylor.}  If every \(u_j\) is locally
bounded, then \(O_k\) is the whole set for \(k\) large, and the
truncations stabilize.

\emph{No mass on pluripolar sets.}  Let \(A\) be pluripolar; by
\cite{Vu19}, after a local reduction it is contained in
\(\{v=-\infty\}\) for some psh \(v\), which we may take among the
truncated potentials.  On \(A\), every \(O_k\) misses \(A\), so
\(\one_{O_k}(\cdot)(A)=0\) for every \(k\), and the increasing limit
gives \(\langle\cdots\rangle(A)=0\).

\emph{Monotonicity of total masses.}  If
\(\varphi_j\preceq\psi_j\) for every \(j\), with \(\psi_j\) in the same
class, then
\[
 \int_X\langle\theta_{\varphi_1}\wedge\cdots\wedge\theta_{\varphi_p}
 \wedge\omega^{n-p}\rangle
 \leq
 \int_X\langle\theta_{\psi_1}\wedge\cdots\wedge\theta_{\psi_p}
 \wedge\omega^{n-p}\rangle .
\]
This is \emph{not} a pointwise inequality between the two currents, and it
is not proved by integration by parts on the truncations, because the
error terms supported on the truncation boundaries do not obviously
vanish.  The unrestricted statement is
\cite[Theorem~1.1]{DDNL18mono}; its predecessor, under the hypothesis that
the potentials have small unbounded locus, is
\cite[Theorem~1.16]{BEGZ10}.  We import the former.

\emph{Bimeromorphic invariance.}  Let \(\pi:Y\to X\) be a modification
with exceptional locus \(F\) and indeterminacy locus of \(\pi^{-1}\)
equal to \(Z=\pi(F)\).  Both \(F\) and \(Z\) are pluripolar --- they are
proper analytic subsets --- so no mass sits on them.  On the complement,
\(\pi\) is a biholomorphism and the change-of-variables formula applies
term by term to the truncations.  Hence each mixed total mass is
preserved.  This is the statement used in
\Cref{thm:tna-bimeromorphic}.

\begin{warning}{Monotonicity is about masses, not measures}
\label{warn:ch51-monotonicity}
\Cref{thm:complex-npp-ledger} asserts that replacing a potential by a
less singular one in the same class can only increase the \emph{total}
mixed masses.  It does not assert any inequality between the two
Monge--Amp\`ere measures.  On \(\mathbf P^1\), take \(\varphi=\max\{G_0,
-1\}\) and \(\psi=0\): both have mass \(1\), yet
\(\MA_\omega(\varphi)\) charges a circle and \(\MA_\omega(\psi)=\omega\)
does not.  This is computed in \Cref{ex:ch51-p1-envelope}.
\end{warning}

\subsection{Masses}
\label{sec:ch51-masses}

We abbreviate
\[
 \theta_{\varphi_1}\wedge\cdots\wedge\theta_{\varphi_p}
 :=\left\langle
   (\theta_1+\ddc\varphi_1)\wedge\cdots\wedge
   (\theta_p+\ddc\varphi_p)\right\rangle .
\]
For \(\varphi\in\PSH(X,\theta)\), its nonpluripolar mass is
\[
                         M_\theta(\varphi)
 :=\int_X\theta_\varphi^n.
\tag{51.17}\label{eq:complex-npp-mass}
\]
Write
\[
 \PSH(X,\theta)_{>0}
 :=\{\varphi\in\PSH(X,\theta)\mid M_\theta(\varphi)>0\}.
\tag{51.18}\label{eq:complex-positive-mass}
\]
If this set is nonempty, then \(\{\theta\}\) is big.

We also write \(\MA_\theta(\varphi):=\langle\theta_\varphi^n\rangle\) for
the (unnormalized) Monge--Amp\`ere measure, so that
\(M_\theta(\varphi)=\int_X\MA_\theta(\varphi)\).  When
\(M_\theta(\varphi)>0\) the probability measure
\(M_\theta(\varphi)^{-1}\MA_\theta(\varphi)\) is available; when the mass
vanishes it is not, and no statement of this book divides by a mass
without a positivity hypothesis in the same sentence.  This is
\eqref{eq:appJ-npp}.

The exact inequalities are the following, and they are used constantly.

\begin{proposition}\label{prop:ch51-mass-inequalities}
Let \(\{\theta\}\) be pseudoeffective and \(\varphi,\psi\in
\PSH(X,\theta)\).
\begin{enumerate}
\item\label{itm:ch51-mi:monotone}
      \(\varphi\preceq\psi\Rightarrow
      M_\theta(\varphi)\leq M_\theta(\psi)\).
\item\label{itm:ch51-mi:minimal}
      \(M_\theta(\varphi)\leq M_\theta(V_\theta)=:\vol(\{\theta\})\) for
      every \(\varphi\), and \(\vol(\{\theta\})>0\) if and only if
      \(\{\theta\}\) is big.
\item\label{itm:ch51-mi:nef}
      If \(\{\theta\}\) is nef, then
      \(\vol(\{\theta\})=\int_X\{\theta\}^n\).  If \(\{\theta\}\) is
      merely pseudoeffective, the two numbers can differ; on a surface one
      always has \(\vol(\{\theta\})\geq\int_X\{\theta\}^2\), with strict
      inequality as soon as the negative part of the divisorial Zariski
      decomposition is nonzero.
\item\label{itm:ch51-mi:translation}
      \(M_\theta(\varphi+c)=M_\theta(\varphi)\) for \(c\in\mathbf R\), and
      \(M_{\lambda\theta}(\lambda\varphi)=\lambda^nM_\theta(\varphi)\) for
      \(\lambda>0\).
\item\label{itm:ch51-mi:perturb}
      If \(\omega\geq0\) is smooth, then
      \(M_{\theta+\omega}(\varphi)\geq M_\theta(\varphi)\); in particular
      a potential of vanishing mass in \(\{\theta\}\) may acquire positive
      mass after an arbitrarily small K\"ahler perturbation.
\end{enumerate}
\end{proposition}

\begin{proof}
Item \cref{itm:ch51-mi:monotone} is
\Cref{thm:complex-npp-ledger} with all entries equal.  Item
\cref{itm:ch51-mi:minimal} follows from it and
\Cref{def:complex-minimal-potential}, since \(\varphi\preceq V_\theta\)
for every \(\varphi\); the characterization of positivity of
\(\vol(\{\theta\})\) is the theorem of Boucksom
\cite{Bou02,Bou02b} that the volume of a pseudoeffective class is
positive exactly when the class is big.  Item \cref{itm:ch51-mi:nef}: for
nef \(\{\theta\}\), the smooth potentials \(\varphi_\epsilon\) of
\Cref{def:ch51-positivity} make \(V_{\theta+\epsilon\omega}\) bounded, so
\(\MA_{\theta+\epsilon\omega}(V_{\theta+\epsilon\omega})\) is the
Bedford--Taylor measure of a bounded potential and has total mass
\(\int_X(\{\theta\}+\epsilon\{\omega\})^n\); letting \(\epsilon\downarrow0\)
and using item \cref{itm:ch51-mi:perturb} together with continuity of the
cohomological expression gives the claim.  The surface statement is the
divisorial Zariski decomposition \(\{\theta\}=P+N\) of \cite{Bou04}: one
has \(\vol(\{\theta\})=P^2\), \(P\cdot N=0\), and \(N^2\leq0\) with
equality only for \(N=0\), whence
\(\int_X\{\theta\}^2=P^2+N^2\leq\vol(\{\theta\})\).
\Cref{ex:ch51-blowup-mass} realizes the strict case with
\(\vol(\{\theta\})=1\) and \(\int_X\{\theta\}^2=0\).

Item \cref{itm:ch51-mi:translation} is immediate, because
\(\theta_{\varphi+c}=\theta_\varphi\) and the nonpluripolar construction
is unchanged.  Item \cref{itm:ch51-mi:perturb} follows from
\(\theta_\varphi+\omega\geq\theta_\varphi\) and expanding
\(\langle(\theta_\varphi+\omega)^n\rangle\) by multilinearity: all mixed
terms are positive.  The final assertion is realized by
\Cref{ex:ch51-p2-lelong}, where \(M_{3\omega}(\varphi)=0\) while
\(M_{3\omega+\epsilon\omega}(\varphi)>0\) for every \(\epsilon>0\),
because the added positivity is not absorbed by the base locus.
\end{proof}

\begin{warning}{Mass is not a cohomological invariant}
\label{warn:ch51-mass-not-cohomological}
For big non-nef classes, \(\vol(\{\theta\})\neq\int_X\{\theta\}^n\) in
general, and the difference is not a sign error but the whole content of
divisorial Zariski decomposition.  \Cref{ex:ch51-blowup-mass} exhibits a
class with \(\int_X\{\theta\}^2=0\) and \(\vol(\{\theta\})=1\).  Any
argument that computes a nonpluripolar mass by expanding a cohomology
class must first verify nefness.
\end{warning}

\subsection{Continuity along monotone families}
\label{sec:ch51-mass-continuity}

\begin{proposition}\label{prop:complex-mass-continuity}
Let \((\varphi_i)\) be an increasing net in \(\PSH(X,\theta)\), uniformly
bounded above, and put
\(\varphi=\uscsup_i\varphi_i\).  Then
\[
                   M_\theta(\varphi)=\lim_iM_\theta(\varphi_i).
\tag{51.19}\label{eq:complex-mass-increasing}
\]
For decreasing families, equality requires the appropriate
full-mass or strong-singularity convergence hypothesis; it is false as
an unrestricted formal rule.
\end{proposition}

\begin{proof}
Mass monotonicity shows that \(M_\theta(\varphi_i)\) is an increasing
net bounded above by \(M_\theta(\varphi)\); let its limit be \(L\).
Choquet's lemma provides a countable set of indices \(j_1,j_2,\ldots\)
whose potentials have usc supremum \(\varphi\).  Recursively use
directedness to choose
\[
 i_a\geq i_{a-1},j_1,\ldots,j_a .
\]
Then \(\varphi_{i_a}\) increases almost everywhere to \(\varphi\).
The increasing-sequence semicontinuity theorem for nonpluripolar
products gives
\[
 M_\theta(\varphi)
 \leq\lim_{a\to\infty}M_\theta(\varphi_{i_a})
 \leq L.
\]
The opposite inequality \(L\leq M_\theta(\varphi)\) is the monotonicity
bound already noted.  This proves \eqref{eq:complex-mass-increasing}
for the original net, without assuming a countable cofinal subset.
\end{proof}

\begin{warning}{The three continuity rules, and only those}
\label{warn:ch51-continuity-rules}
The rules that may be used without further qualification are exactly
those listed under ``safe continuity rules'' in \Cref{sec:appH-npp}:
\begin{enumerate}
\item increasing nets, uniformly bounded above: masses converge, by
      \Cref{prop:complex-mass-continuity};
\item decreasing sequences with a common singularity type and uniformly
      controlled energy: the nonpluripolar products converge weakly;
\item nothing else.
\end{enumerate}
For an arbitrary decreasing sequence, neither statement holds.  On
\(\mathbf P^1\) with \(\int_X\omega=1\), let
\(\varphi_j:=\max\{G_0,-j\}\), where \(G_0\) is the Green function of
\Cref{sec:ch51-unit-tests}.  Then \(\varphi_j\searrow G_0\) and
\(M_\omega(\varphi_j)=1\) for every \(j\), each \(\varphi_j\) being
bounded.  On the other hand \(\omega+\ddc G_0=\delta_0\) is carried by
the polar set \(\{G_0=-\infty\}=\{0\}\), so
\Cref{lem:ch51-npp-bidegree-one} gives
\[
 \langle\omega+\ddc G_0\rangle
 =\one_{\{G_0>-\infty\}}\delta_0=0,
 \qquad
 M_\omega(G_0)=0 .
\]
Thus masses drop discontinuously in the limit, from \(1\) to \(0\).  The Monge--Amp\`ere measures
\(\MA_\omega(\varphi_j)\) do converge weakly, to \(\delta_0\); but their
nonpluripolar limits do not, because \(\delta_0\) is carried by a polar
set.  Both phenomena occur again, in higher codimension and with less
warning, in \Cref{chap:operations-trace}.
\end{warning}

\subsection{Two masses computed}
\label{sec:ch51-masses-computed}

The following two examples are on the same surface and are used
throughout \cref{part:complex} whenever a statement has to be tested
against a class that is not K\"ahler.

\begin{example}[A big non-nef class: masses on a blown-up plane]
\label{ex:ch51-blowup-mass}
Let \(\pi:X\to\mathbf P^2\) be the blow-up of a point \(p\), with
exceptional divisor \(E\).  In \(H^{1,1}(X,\mathbf R)\) write
\(H:=\pi^*c_1(\cO_{\mathbf P^2}(1))\) and \(e:=\{E\}\), so that
\[
 H^2=1,\qquad H\cdot e=0,\qquad e^2=-1 .
\]
Set \(\ell:=H-e\), the class of the strict transform of a line through
\(p\).  The nef cone of \(X\) is spanned by \(H\) and \(\ell\), and the
pseudoeffective cone by \(\ell\) and \(e\).

Take
\[
 \alpha:=H+e=\ell+2e .
\]
It is pseudoeffective, being a nonnegative combination of \(\ell\) and
\(e\), and it is \emph{not} nef, because
\[
 \alpha\cdot e=H\cdot e+e^2=-1<0 .
\]
Its self-intersection is
\[
 \int_X\alpha^2=H^2+2H\cdot e+e^2=1+0-1=0 .
\]

Fix a smooth Hermitian metric \(h\) on \(\cO_X(E)\) with curvature form
\(\Theta\in c_1(\cO_X(E))\), let \(s_E\) be the canonical section, and
put
\[
 \theta:=\pi^*\omega_{\mathrm{FS}}+\Theta\in\alpha,
 \qquad
 V:=\log|s_E|^2_h ,
\]
normalized so that \(\int_{\mathbf P^2}\omega_{\mathrm{FS}}^2=1\).  By
Poincar\'e--Lelong in the form \eqref{eq:appJ-complex-section-pl},
\(\ddc V=[E]-\Theta\), so
\[
 \theta_V=\pi^*\omega_{\mathrm{FS}}+[E]\geq0,
\]
and \(V\in\PSH(X,\theta)\).

\emph{The minimal singularity type.}  We claim
\(V\sim V_\theta\), so that the potential of minimal singularities in
\(\alpha\) is unbounded, with
\[
 \nu(V_\theta,E)=1 .
\]
Indeed, let \(\varphi\in\PSH(X,\theta)\) be arbitrary and let
\(\nu:=\nu(\varphi,E)\).  The Siu decomposition of
\Cref{sec:ch51-divisorial} gives \(\theta_\varphi=\nu[E]+R\) with \(R\)
closed positive and \(\nu(R,E)=0\); a closed positive current with
vanishing generic Lelong number along an irreducible curve has
nonnegative intersection with that curve \cite{Bou04}, so
\[
 0\leq\{R\}\cdot e=(\alpha-\nu e)\cdot e=-1+\nu ,
\]
that is \(\nu\geq1\).  Consequently \(\theta_\varphi\geq[E]\), and
\[
 \pi^*\omega_{\mathrm{FS}}+\ddc(\varphi-V)
 =\theta_\varphi-[E]\geq0 ,
\]
so \(\varphi-V\in\PSH(X,\pi^*\omega_{\mathrm{FS}})\) is upper
semicontinuous on the compact \(X\), hence bounded above.  Thus
\(\varphi\preceq V\) for every \(\varphi\), which is the claim.

\emph{The volume.}  By \Cref{thm:complex-npp-ledger} the nonpluripolar
product charges no pluripolar set, and \(E\) is pluripolar; on
\(X\setminus E\) the current \(\theta_V\) is the smooth semipositive form
\(\pi^*\omega_{\mathrm{FS}}\).  Therefore
\[
 \langle\theta_V^2\rangle
 =\one_{X\setminus E}\bigl(\pi^*\omega_{\mathrm{FS}}\bigr)^2,
 \qquad
 \vol(\alpha)=M_\theta(V_\theta)
 =\int_X\bigl(\pi^*\omega_{\mathrm{FS}}\bigr)^2
 =\int_{\mathbf P^2}\omega_{\mathrm{FS}}^2=1 .
\]
So
\[
 \vol(\alpha)=1
 \qquad\text{while}\qquad
 \int_X\alpha^2=0 .
\]
This is \Cref{warn:ch51-mass-not-cohomological} realized with the
smallest possible numbers.  It is consistent with the divisorial Zariski
decomposition \(\alpha=P+N\) with \(P=H\), \(N=e\): indeed
\(P^2=1=\vol(\alpha)\), \(P\cdot N=0\), and \(N^2=-1\), so
\(\int_X\alpha^2=P^2+N^2=0\).

\emph{A family of masses.}  For \(\tau\in[0,1]\) set
\[
 \eta_\tau:=\pi^*\omega_{\mathrm{FS}}-\tau\Theta
 \ \in\ H-\tau e ,
\]
a smooth form in a nef class, since \(H-\tau e=\tau\ell+(1-\tau)H\) lies
in the nef cone.  By \Cref{prop:ch51-vtheta-bounded-iff}, item
\cref{itm:ch51-vtheta:kahler}, the potential \(V_{\eta_\tau}\) is bounded.
Define
\[
 \varphi_\tau:=(1+\tau)V+V_{\eta_\tau}\ \in\PSH(X,\theta),
 \qquad
 \theta_{\varphi_\tau}
 =\bigl(\eta_\tau+\ddc V_{\eta_\tau}\bigr)+(1+\tau)[E] .
\]
Its divisorial Lelong number is \(\nu(\varphi_\tau,E)=1+\tau\), and
exactly as above the nonpluripolar mass is computed off \(E\), where the
potential is bounded, so that the Bedford--Taylor measure of a bounded
potential in a nef class gives
\[
 M_\theta(\varphi_\tau)
 =\int_X(H-\tau e)^2
 =1-\tau^2 .
\]
Thus, as the imposed multiplicity along \(E\) increases from \(1\) to
\(2\), the nonpluripolar mass decreases from \(1\) to \(0\) along the
concave profile \(\tau\mapsto1-\tau^2\).  For \(\tau>1\) the class
\(H-\tau e\) is no longer pseudoeffective and there is no such potential
at all.  This profile is the elementary prototype of the
Duistermaat--Heckman measures of
\Cref{chap:transcendental-okounkov-dh}: the associated measure on the
parameter is \(-\rd(1-\tau^2)=2\tau\,\rd\tau\) on \([0,1]\), of total mass
\(1=\vol(\alpha)\).
\end{example}

\begin{example}[A nef class of vanishing volume]
\label{ex:ch51-nef-zero-volume}
On the same surface take \(\beta:=\ell=H-e\), the class of the strict
transform of a line through \(p\).  It is nef, since it spans an extremal
ray of the nef cone, and
\[
 \int_X\beta^2=H^2-2H\cdot e+e^2=1-0-1=0 .
\]
Choose a smooth \(\theta'\in\beta\).  By
\Cref{prop:ch51-vtheta-bounded-iff}, item \cref{itm:ch51-vtheta:kahler},
\(V_{\theta'}\) is bounded, so \(\MA_{\theta'}(V_{\theta'})\) is a
Bedford--Taylor measure of total mass \(\int_X\beta^2=0\); that is,
\[
 \vol(\beta)=0,
 \qquad
 M_{\theta'}(\varphi)=0
 \quad\text{for every }\varphi\in\PSH(X,\theta').
\]
Consequently: there is no probability normalization of
\(\MA_{\theta'}\); there is no positive-mass model potential in
\(\beta\); \Cref{thm:complex-p-projection} has no content in \(\beta\);
and the relative finite-energy classes of
\Cref{def:appH-relative-full-mass} are not defined.  Every one of these
expressions becomes available after replacing \(\beta\) by
\(\beta+\epsilon\{\omega\}\) with \(\omega\) K\"ahler, for which
\(\vol=\epsilon^2\int_X\{\omega\}^2+2\epsilon\,\beta\cdot\{\omega\}>0\).
The perturbative inverse system of
\Cref{chap:transcendental-metrics} is precisely the systematic form of
this replacement, and \Cref{rem:ch51-zero-volume} is its motivation.
\end{example}

\section{The \texorpdfstring{\(P\)}{P}-envelope}
\label{sec:ch51-p-envelope}

An envelope replaces a potential by the largest potential of a prescribed
kind lying below it.  Two such operators are needed in
\cref{part:complex}, and they are genuinely different.  The
\(P\)-envelope keeps the ordinary singularity type as a constraint and
maximizes; the \(\mathcal I\)-envelope of \Cref{sec:ch51-i-envelope}
keeps only the multiplier ideals and maximizes.  Since the second
constraint is weaker, the second envelope is larger,
\eqref{eq:complex-p-i-comparison}, and the two can differ.

\subsection{Rooftops}
\label{sec:ch51-rooftop}

For \(\varphi,\psi\in\PSH(X,\theta)\), their rooftop is
\[
 \varphi\wedge_\theta\psi
 :=\uscsup\{\eta\in\PSH(X,\theta)\mid
             \eta\leq\varphi,\ \eta\leq\psi\},
\tag{51.20}\label{eq:complex-rooftop}
\]
with value \(-\infty\) if the candidate set is empty.

The rooftop is the greatest lower bound of \(\varphi\) and \(\psi\) in
\(\PSH(X,\theta)\) for the pointwise order, when it exists; it is
\emph{not} \(\min\{\varphi,\psi\}\), which is rarely \(\theta\)-psh.
Elementary properties, all immediate from
\Cref{prop:ch51-psh-operations}, item \cref{itm:ch51-op:sup}:
\[
 \begin{gathered}
 \varphi\wedge_\theta\psi\leq\min\{\varphi,\psi\},
 \qquad
 (\varphi+c)\wedge_\theta(\psi+c)=(\varphi\wedge_\theta\psi)+c,\\
 \varphi\leq\varphi'\ \Longrightarrow\
 \varphi\wedge_\theta\psi\leq\varphi'\wedge_\theta\psi .
 \end{gathered}
\]
The rooftop is used here only through
\eqref{eq:complex-p-envelope-rooftop}; its metric role, in the formula
\(d_{1,\phi}(u,v)=E_\phi(u)+E_\phi(v)-2E_\phi(P_\theta(u,v))\) of
\Cref{def:appH-i-d1}, belongs to \Cref{chap:rays-test-curves}.

\begin{lemma}\label{lem:ch51-rooftop-contact}
Let \(u,v\in\PSH(X,\theta)\) with \(P:=u\wedge_\theta v\not\equiv-\infty\)
and suppose \(u,v\) have bounded difference from a fixed positive-mass
model potential \(\phi\).  Then
\[
 \MA_\theta(P)\ \text{is carried by}\ \{P=u\}\cup\{P=v\}.
\]
\end{lemma}

\begin{proof}
This is the standard orthogonality property of obstacle problems, and it
is proved by the balayage argument: if \(x\) lies in the open set
\(\{P<\min(u,v)\}\), then \(P\) is maximal there, because any local
perturbation upwards that stays below \(\min(u,v)\) would contradict the
definition of \(P\) as the largest competitor.  A maximal
\(\theta\)-psh function on an open set has vanishing Monge--Amp\`ere
measure there, by the Bedford--Taylor characterization of maximality for
bounded potentials \cite{BT82}, applied after the relative truncation
\(u_k=\max\{u,\phi-k\}\) that reduces the statement to bounded
potentials and using plurifine locality to pass to the limit.  The
hypothesis of bounded difference from \(\phi\) is what makes the
truncations converge with no loss of mass, by
\Cref{thm:appH-comparison}.
\end{proof}

\subsection{The envelope and model potentials}
\label{sec:ch51-p-envelope-def}

\begin{definition}[\(P\)-envelope and model potential]
\label{def:complex-p-envelope}
For \(\varphi\in\PSH(X,\theta)\), set
\[
 P_\theta[\varphi]
 :=\uscsup\{\psi\in\PSH(X,\theta)\mid
              \psi\leq0,\ \psi\preceq\varphi\}.
\tag{51.21}\label{eq:complex-p-envelope}
\]
Equivalently,
\[
 P_\theta[\varphi]
 =\uscsup_{C\in\mathbf Z_{>0}}
       \bigl((\varphi+C)\wedge_\theta V_\theta\bigr).
\tag{51.22}\label{eq:complex-p-envelope-rooftop}
\]
A positive-mass potential \(\phi\) is a model potential if
\(\phi=P_\theta[\phi]\).
\end{definition}

This book reserves model potential for positive mass.  The
positive-mass hypothesis is essential in the projection theorem below;
the zero-mass case is not silently imported.

Two remarks on the definition.  First, the equivalence of
\eqref{eq:complex-p-envelope} and
\eqref{eq:complex-p-envelope-rooftop} is elementary: a competitor
\(\psi\leq0\) with \(\psi\preceq\varphi\) satisfies \(\psi\leq\varphi+C\)
for some \(C\), and also \(\psi\leq V_\theta\) since
\(V_\theta\) is the largest nonpositive potential; hence
\(\psi\leq(\varphi+C)\wedge_\theta V_\theta\).  Conversely each
\((\varphi+C)\wedge_\theta V_\theta\) is itself a competitor, being
\(\leq V_\theta\leq0\) and \(\leq\varphi+C\).  The family indexed by
\(C\in\mathbf Z_{>0}\) is increasing, so
\Cref{cor:ch51-choquet} and
\Cref{lem:ch51-lelong-increasing} apply to it directly; this is the only
form of the envelope that is ever used in a proof.

Second, the operator depends on \(\theta\) and on the normalization
\(\psi\leq0\).  Replacing \(\theta\) by \(\theta+\ddc f\) replaces
\(P_\theta[\varphi]\) by \(P_{\theta+\ddc f}[\varphi-f]\), which differs
from \(P_\theta[\varphi]-f\) by a bounded function but not in general by
zero.  The \emph{relation} \(\preceq_P\) of
\Cref{def:complex-p-preorder} is unaffected, which is
\Cref{rem:ch51-p-preorder-background}.

\begin{proposition}
\label{prop:ch51-p-elementary}
Let \(\varphi,\psi\in\PSH(X,\theta)\).
\begin{enumerate}
\item\label{itm:ch51-pe:monotone}
      \(\varphi\preceq\psi\Rightarrow
      P_\theta[\varphi]\leq P_\theta[\psi]\); in particular
      \(P_\theta[\varphi]\) depends only on the ordinary singularity type
      of \(\varphi\).
\item\label{itm:ch51-pe:normalize}
      \(P_\theta[\varphi]\leq V_\theta\), and
      \(P_\theta[V_\theta]=V_\theta\).  If \(\varphi\) is bounded, then
      \(P_\theta[\varphi]=V_\theta\).
\item\label{itm:ch51-pe:above}
      If \(\varphi\leq0\), then \(\varphi\leq P_\theta[\varphi]\).  In
      general \(\varphi\preceq P_\theta[\varphi]\).
\item\label{itm:ch51-pe:concave}
      \(P_\theta[(1-t)\varphi+t\psi]\geq
      (1-t)P_\theta[\varphi]+tP_\theta[\psi]\) for \(t\in[0,1]\), the
      two backgrounds being combined accordingly.
\item\label{itm:ch51-pe:max}
      \(P_\theta[\max\{\varphi,\psi\}]
      \geq\max\{P_\theta[\varphi],P_\theta[\psi]\}\).
\end{enumerate}
\end{proposition}

\begin{proof}
Item \cref{itm:ch51-pe:monotone}: if \(\varphi\leq\psi+C\) then every
competitor for \(P_\theta[\varphi]\) is a competitor for
\(P_\theta[\psi]\), because \(\preceq\) is transitive; take the
regularized supremum.  Applying this in both directions when
\(\varphi\sim\psi\) gives the dependence statement.

Item \cref{itm:ch51-pe:normalize}: every competitor is \(\leq0\), hence
\(\leq V_\theta\) by \eqref{eq:complex-vtheta}, so the regularized
supremum is \(\leq V_\theta\); and \(V_\theta\) itself is a competitor for
\(P_\theta[V_\theta]\).  If \(\varphi\) is bounded then
\(\varphi\sim0\sim V_\theta\), so item \cref{itm:ch51-pe:monotone} gives
\(P_\theta[\varphi]=P_\theta[V_\theta]=V_\theta\).

Item \cref{itm:ch51-pe:above}: if \(\varphi\leq0\) then \(\varphi\) is
itself a competitor.  In general, \(\varphi-\sup_X\varphi\) is, whence
\(\varphi\preceq P_\theta[\varphi]\).

Item \cref{itm:ch51-pe:concave}: if \(\eta_1\) and \(\eta_2\) are
competitors for \(\varphi\) and \(\psi\) respectively, then
\((1-t)\eta_1+t\eta_2\leq0\) and
\((1-t)\eta_1+t\eta_2\preceq(1-t)\varphi+t\psi\), so it is a competitor
for the convex combination.  Item \cref{itm:ch51-pe:max}: a competitor for
\(\varphi\) is a competitor for \(\max\{\varphi,\psi\}\), since
\(\varphi\leq\max\{\varphi,\psi\}\).
\end{proof}

\begin{warning}{The inequalities in items 4 and 5 can be strict}
\label{warn:ch51-p-not-linear}
The \(P\)-envelope is concave and superadditive under maxima but is not
additive, and neither inequality reverses.  This is why the operations
of \Cref{chap:operations-trace} are defined by applying an envelope
\emph{after} the naive operation, as in
\eqref{eq:rays-test-sum}, rather than by asserting that the operation
commutes with the envelope.
\end{warning}

\subsection{The projection theorem}
\label{sec:ch51-p-projection}

\begin{theorem}\label{thm:complex-p-projection}
If \(\varphi\in\PSH(X,\theta)_{>0}\), then
\[
 P_\theta[P_\theta[\varphi]]=P_\theta[\varphi],
\qquad
 M_\theta(P_\theta[\varphi])=M_\theta(\varphi).
\tag{51.23}\label{eq:complex-p-projection}
\]
Moreover,
\[
 P_\theta[\varphi]
 =\sup\left\{\psi\in\PSH(X,\theta)\ \middle|\
 \begin{array}{c}
 \psi\leq0,\ \varphi\preceq\psi,\\
 M_\theta(\psi)=M_\theta(\varphi)
 \end{array}\right\}.
\tag{51.24}\label{eq:complex-p-characterization}
\]
\end{theorem}

\begin{proof}[Proof architecture]
Formula \eqref{eq:complex-p-envelope-rooftop} gives an increasing
approximation by potentials of the singularity type of \(\varphi\).
Mass monotonicity and increasing continuity prove mass preservation.
The rooftop contact inequality concentrates the Monge--Amp\`ere measure
of \(P_\theta[\varphi]\) on its contact set with the normalized obstacle.
The domination principle then shows that every competitor in
\eqref{eq:complex-p-characterization} lies below the envelope.
Using the envelope itself as a competitor proves equality and
idempotence.
\end{proof}

\begin{proof}[The four steps in full]
We now run the argument in detail, isolating the two imported inputs.

\emph{Step 1: the rooftops have the singularity type of \(\varphi\).}
Write \(u_C:=(\varphi+C)\wedge_\theta V_\theta\) for
\(C\in\mathbf Z_{>0}\).  By definition \(u_C\leq\varphi+C\), so
\(u_C\preceq\varphi\).  Conversely, \(\varphi-\sup_X\varphi\leq V_\theta\)
by \eqref{eq:complex-vtheta}, and \(\varphi-\sup_X\varphi\leq\varphi+C\)
as soon as \(C\geq-\sup_X\varphi\); hence \(\varphi-\sup_X\varphi\) is a
competitor in the rooftop \eqref{eq:complex-rooftop} and
\[
 \varphi-\sup_X\varphi\ \leq\ u_C\ \leq\ \varphi+C
 \qquad(C\gg0).
\]
Therefore \(u_C\sim\varphi\), and
\Cref{thm:complex-npp-ledger} applied in both directions gives
\(M_\theta(u_C)=M_\theta(\varphi)\) for all large \(C\).

\emph{Step 2: the envelope preserves mass.}  The family \((u_C)\) is
increasing in \(C\) and bounded above by \(V_\theta\leq0\), with
\(P_\theta[\varphi]=\uscsup_Cu_C\) by
\eqref{eq:complex-p-envelope-rooftop}.  \Cref{prop:complex-mass-continuity}
gives
\[
 M_\theta(P_\theta[\varphi])=\lim_CM_\theta(u_C)=M_\theta(\varphi),
\]
which is the second half of \eqref{eq:complex-p-projection}.  Note that
this step does not use \(M_\theta(\varphi)>0\).  Step 1 also shows
\(\varphi\preceq P_\theta[\varphi]\) and
\(P_\theta[\varphi]\leq V_\theta\leq0\), so \(P_\theta[\varphi]\) belongs
to the competitor set on the right-hand side of
\eqref{eq:complex-p-characterization}; hence \(P_\theta[\varphi]\) is at
most the supremum there.

\emph{Step 3: the contact property.}  This is the first imported input.
For \(\varphi\) of positive mass,
\[
 \MA_\theta\bigl(P_\theta[\varphi]\bigr)
 \ \text{is carried by}\
 \{P_\theta[\varphi]=V_\theta\},
\]
and in fact
\(\MA_\theta(P_\theta[\varphi])\leq
\one_{\{P_\theta[\varphi]=V_\theta\}}\MA_\theta(V_\theta)\).  For each
fixed \(C\) this is \Cref{lem:ch51-rooftop-contact} applied to the pair
\((\varphi+C,V_\theta)\); the passage to the increasing limit, which is
where positivity of the mass is used, is \cite{DDNL18mono}; see also the
survey \cite{DDNLsurv}.  On a curve the statement can be read off from
\Cref{lem:ch51-npp-bidegree-one}, and
\Cref{ex:ch51-p1-envelope} exhibits it with explicit measures.

\emph{Step 4: the characterization and idempotence.}  Let \(\psi\leq0\)
satisfy \(\varphi\preceq\psi\) and
\(M_\theta(\psi)=M_\theta(\varphi)\).  Applying Steps 1--2 to \(\psi\)
gives a model potential \(Q:=P_\theta[\psi]\) with
\(\psi\leq Q\leq0\) and \(M_\theta(Q)=M_\theta(\psi)\).  Also
\(\varphi\preceq\psi\) gives, by
\Cref{prop:ch51-p-elementary}, item \cref{itm:ch51-pe:monotone},
\[
 P:=P_\theta[\varphi]\ \leq\ Q,
 \qquad
 M_\theta(P)=M_\theta(Q)>0 .
\]
By Step 3, \(\MA_\theta(P)\) is carried by \(\{P=V_\theta\}\); on that
set \(P=V_\theta\geq Q\geq P\), hence \(P=Q\) there, so
\(\MA_\theta(P)\) is carried by \(\{P\geq Q\}\).  Since \(P\preceq Q\)
and \(M_\theta(P)=M_\theta(Q)\), we have \(P\in\mathcal E(X,\theta,Q)\)
in the notation of \Cref{def:appH-relative-full-mass}, and \(Q\) lies
there as well.  The domination principle in the relative full-mass class,
\Cref{thm:appH-comparison}, therefore yields \(P\geq Q\), whence
\(P=Q\).  Consequently \(\psi\leq Q=P\), which is the inequality
\(\geq\) missing from \eqref{eq:complex-p-characterization}; combined
with Step 2 this proves the characterization.

Idempotence follows at once: \(P_\theta[\varphi]\) is itself
\(\leq0\), less singular than \(\varphi\), and of the same mass, so
applying the characterization to it gives
\(P_\theta[P_\theta[\varphi]]=P_\theta[\varphi]\).
\end{proof}

\begin{warning}{Positive mass is not decorative}
\label{warn:ch51-p-positive-mass}
Steps 1 and 2 hold for every \(\varphi\); Steps 3 and 4 do not.  The
domination principle of \Cref{thm:appH-comparison} is a statement about
relative full-mass classes over a positive-mass model, and there is no
substitute when the mass vanishes.  \Cref{ex:ch51-nef-zero-volume}
exhibits an entire pseudoeffective class in which every potential has
zero mass, so that \Cref{thm:complex-p-projection} is vacuous there;
\Cref{ex:ch51-p2-lelong} exhibits a single zero-mass potential inside a
K\"ahler class.  In both cases \cref{part:complex} proceeds by K\"ahler
perturbation, never by extending the theorem.
\end{warning}

\subsection{An envelope computed on
\texorpdfstring{\(\mathbf P^1\)}{P1}}
\label{sec:ch51-p1-envelope}

The following computation is completely explicit and settles, by
inspection, several questions that are otherwise easy to get wrong: that
\(P_\theta[\varphi]\) is generally \emph{not} equal to \(\varphi\), that
it is nevertheless of the same ordinary singularity type when the
singularity is analytic, that the two Monge--Amp\`ere measures differ
while the two nonpluripolar masses agree, and that the contact property of
Step 3 above is visible as a support statement.

\begin{example}[The envelope of a half-pole on \(\mathbf P^1\)]
\label{ex:ch51-p1-envelope}
Let \(X=\mathbf P^1\) with \(\omega=\ddc\log(1+|z|^2)\), so
\(\int_X\omega=1\), and let
\[
 G_0(z)=\log\frac{|z|^2}{1+|z|^2},
 \qquad
 \omega+\ddc G_0=\delta_0,\qquad \sup_XG_0=0,
\]
as in \Cref{sec:ch51-unit-tests}.  Put
\[
 \varphi:=\tfrac12G_0\in\PSH(X,\omega).
\]

\emph{Data of \(\varphi\).}  One has
\[
 \omega+\ddc\varphi=\tfrac12\delta_0+\tfrac12\omega,
 \qquad
 \nu(\varphi,0)=\tfrac12,
 \qquad
 \MA_\omega(\varphi)=\tfrac12\omega,
 \qquad
 M_\omega(\varphi)=\tfrac12 ,
\]
the third equality by \Cref{lem:ch51-npp-bidegree-one}, since the atom
sits on the polar set \(\{0\}\).  The class is K\"ahler with
\(\vol(\{\omega\})=1\), so \(\varphi\) loses exactly half the available
mass.

\emph{The envelope.}  Write \(t=\log|z|^2\in\mathbf R\).  A rotationally
invariant function \(\eta\) is \(\omega\)-psh precisely when
\[
 K(t):=\eta(t)+\log(1+\rme^{t})
\]
is convex and nondecreasing with \(K'\leq1\); the total mass condition is
\(K'(+\infty)-K'(-\infty)\leq1\), and
\(\nu(\eta,0)=K'(-\infty)\).  The three constraints defining the
competitors in \eqref{eq:complex-p-envelope} become
\[
 K\ \text{convex},\qquad
 K'\in[\tfrac12,1],\qquad
 K(t)\leq L(t):=\log(1+\rme^{t}),
\]
the middle one because \(\eta\preceq\varphi\) is equivalent, for
rotationally invariant \(\eta\), to \(\nu(\eta,0)\geq\frac12\), and the
last one because \(\eta\leq0\).  Since \(L'(t)=\rme^t/(1+\rme^t)\geq
\frac12\) exactly for \(t\geq0\), the largest admissible \(K\) is
\[
 K(t)=
 \begin{cases}
 \log2+\tfrac t2, & t\leq0,\\[2pt]
 \log(1+\rme^t), & t\geq0,
 \end{cases}
\]
which is convex, has \(K'\in[\frac12,1)\), and satisfies \(K\leq L\)
because \(L(t)-\log2-\frac t2\) vanishes at \(t=0\) and has derivative
\(L'(t)-\frac12\), negative for \(t<0\) and positive for \(t>0\).  By
symmetrization --- the rotational average of a competitor is again a
competitor, by \Cref{prop:ch51-psh-operations}, item
\cref{itm:ch51-op:convex}, applied to the rotations --- this
rotationally invariant maximum is the envelope.  Explicitly,
\[
 P_\omega[\varphi](z)=
 \begin{cases}
 \log2+\tfrac12\log|z|^2-\log(1+|z|^2), & |z|\leq1,\\[3pt]
 0, & |z|\geq1 .
 \end{cases}
\]

\emph{Comparison.}  On \(|z|\leq1\),
\[
 P_\omega[\varphi]-\varphi
 =\log2-\tfrac12\log(1+|z|^2)
 \ \in\ \bigl[\tfrac12\log2,\ \log2\bigr],
\]
and on \(|z|\geq1\) the difference is
\(-\varphi\geq0\), tending to \(0\) at infinity.  So
\[
 \varphi\ \neq\ P_\omega[\varphi],
 \qquad
 \varphi\ \sim\ P_\omega[\varphi],
\]
in accordance with \eqref{eq:complex-analytic-envelope-agreement}: the
singularity of \(\varphi\) is analytic, so the envelope changes the
function but not its ordinary type.  In particular \(\varphi\) is not a
model potential, while \(P_\omega[\varphi]\) is.

\emph{Measures.}  From \(K\) one reads off
\[
 \omega+\ddc P_\omega[\varphi]
 =\tfrac12\delta_0+\one_{\{|z|>1\}}\omega,
 \qquad
 \MA_\omega\bigl(P_\omega[\varphi]\bigr)
 =\one_{\{|z|>1\}}\omega ,
\]
because \(K\) is affine on \(t<0\), agrees with \(L\) on \(t>0\), and has
\(K'(-\infty)=\frac12\), giving the atom.  There is no atom on the circle
\(|z|=1\), since \(K'\) is continuous there with value
\(\frac12=L'(0)\).  Hence
\[
 M_\omega\bigl(P_\omega[\varphi]\bigr)
 =\int_{\{|z|>1\}}\omega
 =1-\frac{1}{1+1}=\frac12
 =M_\omega(\varphi),
\]
which is \eqref{eq:complex-p-projection} verified numerically.  The two
Monge--Amp\`ere measures are different --- \(\frac12\omega\) versus
\(\one_{\{|z|>1\}}\omega\) --- although their masses agree; this is
\Cref{warn:ch51-monotonicity}.  Finally,
\(\MA_\omega(P_\omega[\varphi])\) is carried by
\(\{|z|\geq1\}=\{P_\omega[\varphi]=V_\omega\}\), since \(V_\omega=0\).
That is Step 3 of the proof of \Cref{thm:complex-p-projection}, made
visible.
\end{example}

\begin{remark}[A frequent slip]
\label{rem:ch51-model-vs-type}
It is tempting to write ``\(\frac12G_0\) is the model potential of its
type''.  It is not: it is a potential whose type is model, and the model
representative is the strictly larger function computed above.  The
distinction matters in \Cref{chap:rays-test-curves}, where the levels of
a test curve are required to satisfy
\(P_\theta[\Gamma_\tau]=\Gamma_\tau\) on the nose, not merely up to
bounded error, so that the Legendre transform of
\eqref{eq:rays-forward-legendre} produces a genuine geodesic rather than
a subgeodesic.
\end{remark}

\subsection{The \texorpdfstring{\(P\)}{P}-preorder}
\label{sec:ch51-p-preorder}

\begin{definition}[\(P\)-preorder]
\label{def:complex-p-preorder}
For quasi-psh \(\varphi,\psi\), choose a smooth \(\theta\) for which both
lie in \(\PSH(X,\theta)_{>0}\), after adding a K\"ahler form if needed,
and write
\[
 \varphi\preceq_P\psi
 \quad\Longleftrightarrow\quad
 P_\theta[\varphi]\leq P_\theta[\psi].
\tag{51.25}\label{eq:complex-p-preorder}
\]
This condition is independent of the common background.  Write
\(\varphi\sim_P\psi\) when both comparisons hold.
\end{definition}

The next proposition supplies the reformulation of
\eqref{eq:complex-p-preorder} that is actually used in proofs: the
\(P\)-preorder is a comparison of \emph{ordinary} singularity types,
applied to the two model representatives.  This is what makes the
independence assertion in \Cref{def:complex-p-preorder} plausible and
what makes the relation usable.

\begin{proposition}\label{prop:ch51-p-preorder-reformulation}
Let \(\varphi,\psi\in\PSH(X,\theta)_{>0}\).  Then
\[
 P_\theta[\varphi]\leq P_\theta[\psi]
 \quad\Longleftrightarrow\quad
 P_\theta[\varphi]\preceq P_\theta[\psi].
\]
Consequently \(\preceq_P\) is a preorder on
\(\PSH(X,\theta)_{>0}\), it depends only on the ordinary singularity
types of the two arguments, and
\[
 \varphi\preceq\psi\ \Longrightarrow\ \varphi\preceq_P\psi .
\]
Moreover \(\varphi\sim_P\psi\) if and only if
\(P_\theta[\varphi]=P_\theta[\psi]\), and each \(\sim_P\)-class contains
exactly one model potential.
\end{proposition}

\begin{proof}
The implication \(\Rightarrow\) is trivial.  For \(\Leftarrow\), suppose
\(P_\theta[\varphi]\preceq P_\theta[\psi]\).  By
\Cref{prop:ch51-p-elementary}, item \cref{itm:ch51-pe:monotone},
\[
 P_\theta\bigl[P_\theta[\varphi]\bigr]
 \leq
 P_\theta\bigl[P_\theta[\psi]\bigr],
\]
and by the idempotence in \eqref{eq:complex-p-projection} the two sides
are \(P_\theta[\varphi]\) and \(P_\theta[\psi]\).

Reflexivity and transitivity are inherited from \(\leq\).  Dependence
only on ordinary types is
\Cref{prop:ch51-p-elementary}, item \cref{itm:ch51-pe:monotone}, applied
in both directions, and that item also gives the displayed implication.
If \(\varphi\sim_P\psi\), the two inequalities give
\(P_\theta[\varphi]=P_\theta[\psi]\).  Finally, a model potential
\(\phi\) satisfies \(P_\theta[\phi]=\phi\), so it is the unique model
potential in its \(\sim_P\)-class.
\end{proof}

\begin{remark}[On the independence of the background]
\label{rem:ch51-p-preorder-background}
\Cref{def:complex-p-preorder} asserts that
\eqref{eq:complex-p-preorder} does not depend on the choice of a common
sufficiently positive background \(\theta\).  This is the convention of
the supplied complex book and of \cite{DDNLmetric,DDNLsurv}, where the
relation is introduced on singularity types rather than on potentials in
a fixed class, and it is used in that form in
\Cref{chap:operations-trace,chap:functoriality-trace}.  We do not
reprove it here; what is used below is only the reformulation of
\Cref{prop:ch51-p-preorder-reformulation}, which is internal to a single
background, together with the background-free implication
\(\preceq_P\Rightarrow\preceq_{\mathcal I}\) of
\Cref{cor:complex-comparison-hierarchy} --- and the target of that
implication, the \(\mathcal I\)-preorder of
\Cref{def:complex-i-preorder}, involves no background at all.
\end{remark}

\begin{warning}{Where the background does matter}
\label{warn:ch51-background}
The \emph{relation} \(\preceq_P\) is background-independent; the
\emph{function} \(P_\theta[\varphi]\) is not, and neither is the
number \(M_\theta(\varphi)\).  In
\Cref{chap:transcendental-metrics} the whole point of the inverse system
is that \(P_{\theta+\omega}[\varphi]\) varies with \(\omega\), and that
its variation is what carries the information lost when
\(\vol(\{\theta\})=0\).  Compare
\Cref{prop:rays-kahler-transition}, where the transition maps of the
system are exactly the envelope operators for varying backgrounds.
\end{warning}

\section{The \texorpdfstring{\(\mathcal I\)}{I}-envelope}
\label{sec:ch51-i-envelope}

The second envelope replaces the constraint ``more singular than
\(\varphi\)'' by the weaker constraint ``the same multiplier ideals as
\(\varphi\)''.  Because the constraint is weaker, the envelope is larger,
and because the constraint refers to no background form, the resulting
preorder is manifestly background-free.  That last feature is what makes
the \(\mathcal I\)-envelope the right operator for the K\"ahler
perturbation systems of \Cref{chap:transcendental-metrics} --- see
\Cref{prop:rays-kahler-transition} --- and, at the same time, is what
makes it blind to nonpluripolar mass.

\subsection{The \texorpdfstring{\(\mathcal I\)}{I}-preorder}
\label{sec:ch51-i-preorder}

\begin{definition}[\(\mathcal I\)-preorder]
\label{def:complex-i-preorder}
For quasi-psh \(\varphi,\psi\), write
\[
 \varphi\preceq_{\mathcal I}\psi
 \quad\Longleftrightarrow\quad
 \mathcal I(m\varphi)\subseteq\mathcal I(m\psi)
 \quad\text{for every }m\geq1.
\tag{51.26}\label{eq:complex-i-preorder}
\]
Equivalently, \(\nu(\varphi,E)\geq\nu(\psi,E)\) for every prime divisor
over \(X\).  Write \(\varphi\sim_{\mathcal I}\psi\) for equality in both
directions.
\end{definition}

The asserted equivalence of the two descriptions is
\Cref{thm:complex-strong-openness}, item
\cref{itm:ch51-so:equivalence}, applied to the pair of one-sided
statements; note that the ideal inclusion and the Lelong inequality point
in opposite directions, because a more singular potential has larger
Lelong numbers and smaller multiplier ideals.  Reading the direction
backwards is the commonest error in this circle of ideas, and it is
detected by the one-pole test of \Cref{sec:ch51-unit-tests}: with
\(\varphi=\log|z|^2\) and \(\psi=0\) one has
\(\varphi\preceq\psi\), \(\nu(\varphi,0)=1>0=\nu(\psi,0)\), and
\(\mathcal I(\varphi)=(z)\subsetneq\cO=\mathcal I(\psi)\).

\begin{proposition}
\label{prop:ch51-i-calculus}
Let \(\varphi,\psi,\varphi',\psi'\) be quasi-psh on \(X\) and
\(\lambda>0\).
\begin{enumerate}
\item\label{itm:ch51-ic:preorder}
      \(\preceq_{\mathcal I}\) is a preorder on quasi-psh functions,
      independent of every choice of background form.
\item\label{itm:ch51-ic:sum}
      \(\varphi\preceq_{\mathcal I}\psi\) and
      \(\varphi'\preceq_{\mathcal I}\psi'\) imply
      \(\varphi+\varphi'\preceq_{\mathcal I}\psi+\psi'\),
      \(\max\{\varphi,\varphi'\}\preceq_{\mathcal I}
      \max\{\psi,\psi'\}\), and
      \(\lambda\varphi\preceq_{\mathcal I}\lambda\psi\).
\item\label{itm:ch51-ic:pullback}
      For a modification \(\pi:Y\to X\),
      \(\varphi\preceq_{\mathcal I}\psi\) if and only if
      \(\pi^*\varphi\preceq_{\mathcal I}\pi^*\psi\).
\item\label{itm:ch51-ic:increasing}
      If \(\{\varphi_i\}\) and \(\{\psi_i\}\) are directed families,
      uniformly bounded above, with
      \(\varphi_i\preceq_{\mathcal I}\psi_i\) for every \(i\), then
      \(\uscsup_i\varphi_i\preceq_{\mathcal I}\uscsup_i\psi_i\).
\item\label{itm:ch51-ic:bounded}
      Every bounded quasi-psh function is \(\sim_{\mathcal I}0\).
\end{enumerate}
\end{proposition}

\begin{proof}
Item \cref{itm:ch51-ic:preorder} is immediate from
\eqref{eq:complex-i-preorder}, which mentions no background.  Item
\cref{itm:ch51-ic:sum} follows from the divisorial description together
with \Cref{prop:ch51-divisorial-calculus}, item
\cref{itm:ch51-dc:calculus}: Lelong numbers are additive, positively
homogeneous, and satisfy the minimum rule for maxima, and each of these
operations is monotone.  Item \cref{itm:ch51-ic:pullback} is
\Cref{prop:ch51-divisorial-calculus}, item \cref{itm:ch51-dc:pullback},
together with the observation that the prime divisors over \(Y\) are
cofinal among those over \(X\).  Item \cref{itm:ch51-ic:increasing} is
\Cref{lem:ch51-lelong-increasing} applied to both families:
\(\nu(\uscsup_i\varphi_i,E)=\inf_i\nu(\varphi_i,E)
\geq\inf_i\nu(\psi_i,E)=\nu(\uscsup_i\psi_i,E)\).  Item
\cref{itm:ch51-ic:bounded} is \Cref{prop:ch51-lelong-calculus}, item
\cref{itm:ch51-lc:smooth}.
\end{proof}

\subsection{The envelope}
\label{sec:ch51-i-envelope-def}

\begin{definition}[\(\mathcal I\)-envelope]
\label{def:complex-i-envelope}
For \(\varphi\in\PSH(X,\theta)\), set
\[
 P_\theta[\varphi]_{\mathcal I}
 :=\uscsup\{\psi\in\PSH(X,\theta)\mid
          \psi\leq0,\ \psi\sim_{\mathcal I}\varphi\}.
\tag{51.27}\label{eq:complex-i-envelope}
\]
It is \(\mathcal I\)-model if it is fixed by this operator.
\end{definition}

\begin{theorem}\label{thm:complex-i-projection}
For every \(\varphi\in\PSH(X,\theta)\),
\[
 \varphi\sim_{\mathcal I}P_\theta[\varphi]_{\mathcal I},
\qquad
 P_\theta[P_\theta[\varphi]_{\mathcal I}]_{\mathcal I}
      =P_\theta[\varphi]_{\mathcal I}.
\tag{51.28}\label{eq:complex-i-projection}
\]
Furthermore,
\[
 P_\theta[\varphi]\leq P_\theta[\varphi]_{\mathcal I},
\qquad
 \varphi\sim_{\mathcal I}P_\theta[\varphi].
\tag{51.29}\label{eq:complex-p-i-comparison}
\]
If \(\varphi\) has analytic singularities, then
\[
 \varphi\sim P_\theta[\varphi]
        \sim P_\theta[\varphi]_{\mathcal I}.
\tag{51.30}\label{eq:complex-analytic-envelope-agreement}
\]
\end{theorem}

\begin{proof}
Choquet reduction replaces the family in
\eqref{eq:complex-i-envelope} by an increasing sequence.  Finite maxima
preserve divisorial Lelong numbers, and strong openness shows that its
usc supremum has the same multiplier ideals as \(\varphi\).  This proves
\eqref{eq:complex-i-projection}.  Formula
\eqref{eq:complex-p-envelope-rooftop}, again with strong openness, gives
\(\varphi\sim_{\mathcal I}P_\theta[\varphi]\), hence
\eqref{eq:complex-p-i-comparison}.  For analytic singularities, pass to
a log resolution.  Divisorial orders then control the whole singularity
up to \(O(1)\), so the two envelopes cannot be strictly less singular
than \(\varphi\).
\end{proof}

\begin{proof}[The steps in full]
Write \(F:=\{\psi\in\PSH(X,\theta)\mid\psi\leq0,\
\psi\sim_{\mathcal I}\varphi\}\) for the competitor family in
\eqref{eq:complex-i-envelope}.

\emph{Step 1: \(F\) is nonempty and directed.}  It contains
\(\varphi-\sup_X\varphi\).  If \(\psi_1,\psi_2\in F\), then
\(\max\{\psi_1,\psi_2\}\in\PSH(X,\theta)\) by
\Cref{prop:ch51-psh-operations}, item \cref{itm:ch51-op:max}, it is
\(\leq0\), and by \Cref{prop:ch51-lelong-calculus}, item
\cref{itm:ch51-lc:max}, applied at a generic point of each prime divisor
\(E\) over \(X\),
\[
 \nu(\max\{\psi_1,\psi_2\},E)
 =\min\{\nu(\psi_1,E),\nu(\psi_2,E)\}
 =\nu(\varphi,E).
\]
Hence \(\max\{\psi_1,\psi_2\}\in F\), and \(F\) is directed.

\emph{Step 2: the envelope lies in \(F\), which gives
\eqref{eq:complex-i-projection}.}  By Step 1 and
\Cref{lem:ch51-lelong-increasing},
\[
 \nu\bigl(P_\theta[\varphi]_{\mathcal I},E\bigr)
 =\inf_{\psi\in F}\nu(\psi,E)
 =\nu(\varphi,E)
\]
for every prime divisor \(E\) over \(X\).  Since
\(P_\theta[\varphi]_{\mathcal I}\leq0\), this says
\(P_\theta[\varphi]_{\mathcal I}\in F\); it is therefore the largest
element of \(F\), and in particular
\(\varphi\sim_{\mathcal I}P_\theta[\varphi]_{\mathcal I}\).  As \(F\)
depends only on the \(\mathcal I\)-class of \(\varphi\), the family
attached to \(P_\theta[\varphi]_{\mathcal I}\) is the same \(F\), and
idempotence follows.

\emph{Step 3: \(\varphi\sim_{\mathcal I}P_\theta[\varphi]\).}  By Step 1
of the proof of \Cref{thm:complex-p-projection}, the rooftops
\(u_C=(\varphi+C)\wedge_\theta V_\theta\) satisfy \(u_C\sim\varphi\) for
\(C\) large, hence \(\nu(u_C,E)=\nu(\varphi,E)\) for every \(E\) by
\Cref{prop:ch51-lelong-calculus}, item \cref{itm:ch51-lc:bounded}.  The
family \((u_C)\) is increasing, and
\(P_\theta[\varphi]=\uscsup_Cu_C\); so
\Cref{lem:ch51-lelong-increasing} gives
\(\nu(P_\theta[\varphi],E)=\nu(\varphi,E)\), which is the second half of
\eqref{eq:complex-p-i-comparison}.  This step uses no positivity of
mass.

\emph{Step 4: \(P_\theta[\varphi]\leq P_\theta[\varphi]_{\mathcal I}\).}
By Step 3 and \(P_\theta[\varphi]\leq V_\theta\leq0\), the potential
\(P_\theta[\varphi]\) belongs to \(F\); by Step 2 it is therefore at most
the largest element of \(F\).

\emph{Step 5: analytic singularities.}  Assume \(\varphi\) has analytic
singularities, and let \(\rho:Z\to X\) be a log resolution with
\[
 \rho^*\varphi=c\log|s_D|_h^2+O(1),
\]
\(D\) an effective divisor with simple normal crossing support, \(h\) a
smooth metric on \(\cO_Z(D)\), and \(c>0\).  Put
\(\Psi:=P_\theta[\varphi]_{\mathcal I}\), which by Step 2 satisfies
\(\Psi\sim_{\mathcal I}\varphi\).  By
\Cref{prop:ch51-divisorial-calculus}, items
\cref{itm:ch51-dc:pullback} and \cref{itm:ch51-dc:analytic}, for every
prime divisor \(E\subset Z\),
\[
 \nu(\rho^*\Psi,E)=\nu(\Psi,E)=\nu(\varphi,E)=c\,\ord_E(D).
\]
The Siu decomposition of \Cref{sec:ch51-divisorial}, applied to the
closed positive current \(\rho^*\theta+\ddc\rho^*\Psi\), therefore gives
\[
 \rho^*\theta+\ddc\rho^*\Psi\ \geq\ c\,[D].
\]
Writing \(\Theta_D=c_1(\cO_Z(D),h)\) and using Poincar\'e--Lelong
\eqref{eq:appJ-complex-section-pl} in the form
\(\ddc\log|s_D|_h^2=[D]-\Theta_D\), we obtain
\[
 \bigl(\rho^*\theta-c\,\Theta_D\bigr)
 +\ddc\bigl(\rho^*\Psi-c\log|s_D|_h^2\bigr)\ \geq\ 0 ,
\]
so \(\rho^*\Psi-c\log|s_D|_h^2\) is quasi-psh on the compact manifold
\(Z\), hence bounded above.  Thus
\(\rho^*\Psi\leq c\log|s_D|_h^2+C=\rho^*\varphi+C'\), and since \(\rho\)
is an isomorphism over a dense open subset of \(X\) and both sides are
upper semicontinuous, \(\Psi\preceq\varphi\).  Combining with
\(\varphi\preceq P_\theta[\varphi]\leq\Psi\) from
\Cref{prop:ch51-p-elementary}, item \cref{itm:ch51-pe:above}, and Step 4,
\[
 \varphi\ \sim\ P_\theta[\varphi]\ \sim\ P_\theta[\varphi]_{\mathcal I},
\]
which is \eqref{eq:complex-analytic-envelope-agreement}.
\end{proof}

\begin{remark}[Where each hypothesis is used]
\label{rem:ch51-i-projection-hypotheses}
Steps 1--4 use no positivity of mass and no regularity of \(\varphi\);
they are consequences of \Cref{lem:ch51-lelong-increasing} and the
minimum rule for maxima.  Step 5 uses resolution of singularities and
Siu's decomposition, and it is the only step that upgrades an
\(\mathcal I\)-statement to an ordinary-type statement.  It is essential
that the upgrade is applied to the \emph{envelope} \(\Psi\) and not to an
arbitrary competitor: the argument produces a constant \(C'\) depending
on \(\Psi\), and no uniform constant over the whole family \(F\) is
claimed or needed.
\end{remark}

\begin{warning}{Do not read Step 5 as a general rigidity statement}
\label{warn:ch51-step5}
Step 5 does \emph{not} say that two potentials with the same divisorial
Lelong numbers have bounded difference.  It says that the
\(\mathcal I\)-envelope of a potential with analytic singularities has
bounded difference from it.  The general statement is false, by
\Cref{ex:ch51-disk-three-types}; and the difference is not merely
cosmetic, since \Cref{ex:ch51-cantor} produces an
\(\mathcal I\)-equivalent pair with different nonpluripolar masses.
\end{warning}

\begin{corollary}\label{cor:complex-comparison-hierarchy}
For quasi-psh functions,
\[
 \varphi\preceq\psi
 \Longrightarrow
 \varphi\preceq_P\psi
 \Longrightarrow
 \varphi\preceq_{\mathcal I}\psi.
\tag{51.31}\label{eq:complex-comparison-hierarchy}
\]
Neither reverse implication is asserted.  In particular,
\(\mathcal I\)-equivalence is coarser than \(P\)-equivalence.
\end{corollary}

\begin{proof}
If \(\varphi\preceq\psi\), then every normalized competitor defining
\(P_\theta[\varphi]\) is also a competitor for \(P_\theta[\psi]\), which
gives the first implication.  The second follows from
\eqref{eq:complex-p-i-comparison}: applying the \(\mathcal I\)-envelope
to the ordered \(P\)-envelopes gives the corresponding
\(\mathcal I\)-order.  Applying both implications in both directions
gives the assertion about equivalence relations.
\end{proof}

\begin{proof}[The second implication, in detail]
Assume \(\varphi\preceq_P\psi\), that is
\(P_\theta[\varphi]\leq P_\theta[\psi]\) for a common background in which
both have positive mass.  Monotonicity of Lelong numbers under the
pointwise order, which is \eqref{eq:complex-lelong-number}, gives
\[
 \nu\bigl(P_\theta[\varphi],E\bigr)
 \ \geq\
 \nu\bigl(P_\theta[\psi],E\bigr)
\]
for every prime divisor \(E\) over \(X\), the inequality reversing
because a smaller potential is more singular.  By Step 3 of the proof of
\Cref{thm:complex-i-projection}, the two sides equal
\(\nu(\varphi,E)\) and \(\nu(\psi,E)\) respectively.  Hence
\[
 \nu(\varphi,E)\ \geq\ \nu(\psi,E)
 \qquad\text{for every prime divisor }E\text{ over }X,
\]
which by \Cref{thm:complex-strong-openness}, item
\cref{itm:ch51-so:equivalence}, and
\Cref{def:complex-i-preorder} is exactly
\(\varphi\preceq_{\mathcal I}\psi\).  Note that the conclusion does not
mention \(\theta\), consistently with
\Cref{prop:ch51-i-calculus}, item \cref{itm:ch51-ic:preorder}.
\end{proof}

\begin{corollary}\label{cor:ch51-quotient-maps}
There are canonical surjections of preordered sets
\[
 \QPSH(X)/{\sim}
 \ \longrightarrow\
 \QPSH(X)/{\sim_P}
 \ \longrightarrow\
 \QPSH(X)/{\sim_{\mathcal I}} ,
\]
where the middle quotient is taken on the locus of potentials of
positive mass in a fixed sufficiently positive background.  Both maps are
strictly non-injective:
\Cref{ex:appH-loglog} separates the first two, and
\Cref{ex:ch51-cantor} separates the last two.
\end{corollary}

\begin{proof}
The maps exist and are order preserving by
\Cref{cor:complex-comparison-hierarchy}, and they are surjective because
each is induced by the identity on representatives.  The separations are
the cited examples.
\end{proof}

The chain of \Cref{cor:ch51-quotient-maps} is the source of the
factorization \eqref{eq:verified-trace-factorization}: the trace of
\Cref{chap:operations-trace} depends only on the input
\(\mathcal I\)-class but lands in a well-defined output \(P\)-class.
Hence the map obtained by starting with input \(P\)-classes factors through
the second surjection above.  \Cref{thm:verified-trace-factorization}
records this asymmetric, source-backed factorization.

\subsection{The volume detected by ideals}
\label{sec:ch51-current-volume}

\begin{definition}[Volume detected by ideals]
\label{def:complex-current-volume}
For \(\varphi\in\PSH(X,\theta)\), define
\[
 \vol(\theta,\varphi)
 :=\int_X\bigl(\theta+\ddc
             P_\theta[\varphi]_{\mathcal I}\bigr)^n.
\tag{51.32}\label{eq:complex-current-volume}
\]
It depends only on the current \(\theta_\varphi\).  Moreover,
\[
 M_\theta(\varphi)\leq\vol(\theta,\varphi),
\tag{51.33}\label{eq:complex-mass-volume}
\]
with equality for analytic singularities.
\end{definition}

Three comments.  The nonpluripolar product on the right of
\eqref{eq:complex-current-volume} is, as always, the one of
\eqref{eq:complex-global-npp}; the notation omits the angle brackets only
because the argument is now a specific model potential.  The inequality
\eqref{eq:complex-mass-volume} is
\Cref{thm:complex-npp-ledger} together with
\(P_\theta[\varphi]\leq P_\theta[\varphi]_{\mathcal I}\) from
\eqref{eq:complex-p-i-comparison} and the mass identity
\eqref{eq:complex-p-projection}; explicitly,
\[
 M_\theta(\varphi)
 =M_\theta\bigl(P_\theta[\varphi]\bigr)
 \leq M_\theta\bigl(P_\theta[\varphi]_{\mathcal I}\bigr)
 =\vol(\theta,\varphi),
\]
the middle step because \(P_\theta[\varphi]\preceq
P_\theta[\varphi]_{\mathcal I}\).  Equality for analytic singularities is
\eqref{eq:complex-analytic-envelope-agreement}.  Finally, the dependence
only on the current is the statement that
\(\vol(\theta,\varphi)\) is unchanged when \(\varphi\) is replaced by
\(\varphi+c\), which is
\Cref{prop:ch51-mass-inequalities}, item
\cref{itm:ch51-mi:translation}, together with the fact that
\eqref{eq:complex-i-envelope} depends only on the
\(\mathcal I\)-class, which is invariant under adding constants.

The quantity \(\vol(\theta,\varphi)\) is what
\Cref{prop:tna-homogeneous-volume} identifies with the volume of the
transcendental non-Archimedean metric attached to \(\varphi\), and
\eqref{eq:complex-mass-volume} is exactly the inequality that
\Cref{chap:transcendental-metrics} warns about: the non-Archimedean side
computes the \emph{ideal} volume, not the mass.  The gap can be
maximal, as the next section shows.

\begin{example}[Why the envelopes differ]
\label{ex:complex-loglog-envelope}
On \(\mathbf P^1\), a global \(\omega_{\mathrm{FS}}\)-psh potential with
local model
\[
                        -\log(-\log|z|^2)
\]
at \(0\) can have full nonpluripolar mass and
\(P_{\omega_{\mathrm{FS}}}[\varphi]=0\), although
\(\varphi\not\sim0\).  Its Lelong number is zero.  This example separates
ordinary singularity type from the coarser envelope types and explains
why traces are built from quasi-equisingular approximants rather than
naive restrictions.
\end{example}

The example can be written down and verified completely, which is worth
doing once because it is the standard witness for the strictness of the
first implication in \eqref{eq:complex-comparison-hierarchy}.

\begin{remark}[\Cref{ex:complex-loglog-envelope} made explicit]
\label{rem:ch51-loglog-explicit}
On \(\mathbf P^1\) with \(\omega=\ddc\log(1+|z|^2)\) and
\(t=\log|z|^2\), a rotationally invariant \(\eta\) is \(\omega\)-psh
exactly when \(K(t):=\eta(t)+\log(1+\rme^t)\) is convex and
nondecreasing with \(K'\leq1\); then \(\nu(\eta,0)=K'(-\infty)\), the
Monge--Amp\`ere measure corresponds to \(\rd K'\), and there are atoms of
mass \(K'(-\infty)\) at \(z=0\) and \(1-K'(+\infty)\) at \(z=\infty\).
Take
\[
 K(t):=
 \begin{cases}
 -\log(1-t), & t\leq0,\\
 t, & t\geq0 ,
 \end{cases}
 \qquad
 \eta:=K(t)-\log(1+\rme^t).
\]
Then \(K\) is convex, \(K'(t)=1/(1-t)\) for \(t<0\) and \(K'=1\) for
\(t>0\), continuous at \(0\), with \(K'\leq1\).  Hence
\(\eta\in\PSH(\mathbf P^1,\omega)\), and:
\begin{itemize}
\item near \(z=0\), \(\eta=-\log\bigl(1-\log|z|^2\bigr)+O(1)
      =-\log\bigl(-\log|z|^2\bigr)+O(1)\), so \(\eta\) is unbounded and
      \(\eta\not\sim0\);
\item \(\nu(\eta,0)=K'(-\infty)=0\), and every other Lelong number
      vanishes, so \(\eta\sim_{\mathcal I}0\);
\item \(K'(-\infty)=0\) and \(K'(+\infty)=1\), so
      \(\MA_\omega(\eta)=\rd K'\) carries no atom and
      \(M_\omega(\eta)=1=\vol(\{\omega\})\);
\item \(\eta\leq0\) with \(\sup_{\mathbf P^1}\eta=0\).
\end{itemize}
Since \(V_\omega=0\) and \(M_\omega(\eta)=M_\omega(V_\omega)\), the
characterization \eqref{eq:complex-p-characterization} applies with
\(\psi=V_\omega\) as a competitor and gives
\[
 P_\omega[\eta]=V_\omega=0 .
\]
Thus \(\eta\sim_P0\) and \(\eta\sim_{\mathcal I}0\), while
\(\eta\not\sim0\).  The first implication in
\eqref{eq:complex-comparison-hierarchy} is therefore strict.
\end{remark}

\begin{corollary}
\label{cor:ch51-full-mass-envelope}
Let \(\{\theta\}\) be pseudoeffective with \(\vol(\{\theta\})>0\), and let
\(\varphi\in\PSH(X,\theta)\).  Then
\[
 M_\theta(\varphi)=\vol(\{\theta\})
 \quad\Longleftrightarrow\quad
 P_\theta[\varphi]=V_\theta
 \quad\Longleftrightarrow\quad
 \varphi\sim_PV_\theta .
\]
\end{corollary}

\begin{proof}
If \(M_\theta(\varphi)=\vol(\{\theta\})=M_\theta(V_\theta)\), then
\(V_\theta\) is a competitor in
\eqref{eq:complex-p-characterization}: it is \(\leq0\), it satisfies
\(\varphi\preceq V_\theta\) by
\Cref{def:complex-minimal-potential}, and it has the required mass.
Hence \(P_\theta[\varphi]\geq V_\theta\), and the reverse inequality is
\Cref{prop:ch51-p-elementary}, item \cref{itm:ch51-pe:normalize}.
Conversely, if \(P_\theta[\varphi]=V_\theta\), then
\eqref{eq:complex-p-projection} gives
\(M_\theta(\varphi)=M_\theta(V_\theta)=\vol(\{\theta\})\).  The second
equivalence is \Cref{prop:ch51-p-preorder-reformulation}, since
\(V_\theta\) is model.
\end{proof}

\section{\texorpdfstring{\(P\)}{P}-equivalence versus
\texorpdfstring{\(\mathcal I\)}{I}-equivalence}
\label{sec:ch51-p-vs-i}

This section is the reason the chapter exists.  Everything above has been
preparation for two statements: that \(P\)-comparison implies
\(\mathcal I\)-comparison, and that the converse fails in a way that
destroys nonpluripolar mass.  The first is
\Cref{cor:complex-comparison-hierarchy}, proved twice above.  The second
is the example below.

\subsection{The implications, collected}
\label{sec:ch51-implications}

\begin{theorem}
\label{thm:ch51-hierarchy-strict}
Let \(X\) be a connected compact K\"ahler manifold and let
\(\varphi,\psi\) be quasi-psh functions on \(X\), placed in a common
background of positive mass where required.
\begin{enumerate}
\item\label{itm:ch51-hs:implications}
      \(\varphi\preceq\psi\Rightarrow\varphi\preceq_P\psi
      \Rightarrow\varphi\preceq_{\mathcal I}\psi\), and hence
      \(\varphi\sim\psi\Rightarrow\varphi\sim_P\psi
      \Rightarrow\varphi\sim_{\mathcal I}\psi\).
\item\label{itm:ch51-hs:first-strict}
      The first implication is strict already for \(X=\mathbf P^1\):
      there is \(\eta\in\PSH(\mathbf P^1,\omega)\) with
      \(\eta\sim_P0\) and \(\eta\not\sim0\).
\item\label{itm:ch51-hs:second-strict}
      The second implication is strict already for \(X=\mathbf P^1\):
      there is \(\varphi\in\PSH(\mathbf P^1,\omega)\) with
      \(\varphi\sim_{\mathcal I}0\) and \(\varphi\not\sim_P0\).
      Moreover \(\varphi\) may be chosen with
      \(M_\omega(\varphi)=\tfrac12\), while \(M_\omega(0)=1\).
\item\label{itm:ch51-hs:mass}
      \(\sim_P\) determines the nonpluripolar mass:
      \(\varphi\sim_P\psi\Rightarrow M_\theta(\varphi)=M_\theta(\psi)\).
      \(\sim_{\mathcal I}\) does not.
\item\label{itm:ch51-hs:ideal-volume}
      \(\sim_{\mathcal I}\) determines the ideal volume
      \(\vol(\theta,\cdot)\) of
      \eqref{eq:complex-current-volume}, and
      \(M_\theta(\varphi)\leq\vol(\theta,\varphi)\) with equality for
      analytic singularities.
\end{enumerate}
\end{theorem}

\begin{proof}
Item \cref{itm:ch51-hs:implications} is
\Cref{cor:complex-comparison-hierarchy}.  Item
\cref{itm:ch51-hs:first-strict} is
\Cref{rem:ch51-loglog-explicit}.  Item
\cref{itm:ch51-hs:second-strict} is \Cref{ex:ch51-cantor} below.

Item \cref{itm:ch51-hs:mass}: if \(\varphi\sim_P\psi\), then
\(P_\theta[\varphi]=P_\theta[\psi]\) by
\Cref{prop:ch51-p-preorder-reformulation}, and
\eqref{eq:complex-p-projection} gives
\(M_\theta(\varphi)=M_\theta(P_\theta[\varphi])
=M_\theta(P_\theta[\psi])=M_\theta(\psi)\).  That
\(\sim_{\mathcal I}\) does not determine the mass is again
\Cref{ex:ch51-cantor}, where the two masses are \(\tfrac12\) and \(1\).

Item \cref{itm:ch51-hs:ideal-volume}: the operator
\eqref{eq:complex-i-envelope} depends only on the
\(\mathcal I\)-class, so \(\vol(\theta,\cdot)\) does; the inequality and
the equality case are \eqref{eq:complex-mass-volume}.
\end{proof}

\subsection{A potential that is
\texorpdfstring{\(\mathcal I\)}{I}-equivalent, but not
\texorpdfstring{\(P\)}{P}-equivalent, to zero}
\label{sec:ch51-cantor}

The construction uses one classical input: a compact set of vanishing
logarithmic capacity carrying a probability measure without atoms.  We
recall it explicitly so that the numbers below can be checked.

\begin{lemma}\label{lem:ch51-polar-cantor}
Let \((\ell_m)_{m\geq0}\) be lengths with \(\ell_0=1\) and
\(\ell_{m+1}\leq\ell_m/2\), and let
\(F=\bigcap_mF_m\subset[0,1]\subset\mathbf C\) be the generalized Cantor
set in which \(F_m\) is a disjoint union of \(2^m\) closed intervals of
length \(\ell_m\), each containing two intervals of the next stage.  Let
\(\mu\) be the measure giving each of the \(2^m\) intervals of \(F_m\)
mass \(2^{-m}\).  Then \(\mu\) is a nonatomic Borel probability measure
carried by \(F\), and if
\[
 \sum_{m\geq0}2^{-m}\log\frac{1}{\ell_m}=+\infty ,
\]
then \(F\) has vanishing logarithmic capacity, hence is polar.  The
choice \(\ell_0=1\) and \(\ell_m=\rme^{-4^m}\) for \(m\geq1\)
satisfies both hypotheses, since the summands for \(m\geq1\) equal \(2^m\).
\end{lemma}

\begin{proof}[Classical, with the divergence check]
The measure \(\mu\) is the weak limit of the normalized Lebesgue measures
on \(F_m\); it is nonatomic because every point of \(F\) lies in nested
intervals of \(\mu\)-mass \(2^{-m}\to0\).  The capacity criterion for
generalized Cantor sets is \cite[p.~31]{Car83}: the set is polar exactly
when the displayed series diverges.  With \(\ell_0=1\) and
\(\ell_m=\rme^{-4^m}\) for \(m\geq1\), one has
\(\ell_1/\ell_0=\rme^{-4}\leq1/2\); for \(m\geq1\),
\(\ell_{m+1}/\ell_m=\rme^{-3\cdot4^m}\leq1/2\) and
\(2^{-m}\log(1/\ell_m)=2^m\), so the series diverges.
\end{proof}

\begin{example}[\(\mathcal I\)-equivalent, not \(P\)-equivalent]
\label{ex:ch51-cantor}
Let \(X=\mathbf P^1\) with \(\omega=\ddc\log(1+|z|^2)\), normalized by
\(\int_X\omega=1\), and let \(G_p\) be the Green function of
\Cref{sec:ch51-unit-tests}, so that
\[
 G_p(z)=\log\frac{|z-p|^2}{(1+|z|^2)(1+|p|^2)}\leq0,
 \qquad
 \omega+\ddc G_p=\delta_p .
\]
Let \(F\) and \(\mu\) be as in \Cref{lem:ch51-polar-cantor}, and set
\[
 G_\mu(z):=\int_XG_p(z)\,\rd\mu(p),
 \qquad
 \varphi:=\tfrac12G_\mu .
\]

\emph{\(\varphi\) is a genuine potential.}  Each \(G_p\leq0\), so
\(G_\mu\leq0\).  For \(z\) at chordal distance at least \(\delta\) from
\(F\) one has \(G_p(z)\geq2\log\delta\) for every \(p\in F\), so
\(G_\mu(z)\geq2\log\delta>-\infty\); in particular
\(G_\mu\not\equiv-\infty\).  Fubini gives
\[
 \omega+\ddc G_\mu=\mu ,
 \qquad
 \omega+\ddc\varphi=\tfrac12\omega+\tfrac12\mu\ \geq0,
\]
so \(\varphi\in\PSH(X,\omega)\) with \(\varphi\leq0\).

\emph{All Lelong numbers vanish.}  On a Riemann surface every
modification is an isomorphism, so the prime divisors over \(X\) are the
points of \(X\), and
\Cref{prop:ch51-lelong-three}, item \cref{itm:ch51-lelong:mass}, gives
\[
 \nu(\varphi,x)=(\omega+\ddc\varphi)(\{x\})
 =\tfrac12\omega(\{x\})+\tfrac12\mu(\{x\})=0
\]
for every \(x\), because \(\omega\) is smooth and \(\mu\) is nonatomic.
By \Cref{thm:complex-strong-openness}, item
\cref{itm:ch51-so:equivalence}, or directly by Skoda's bound
\Cref{prop:ch51-lelong-calculus}, item \cref{itm:ch51-lc:smooth}, applied
to \(m\varphi\) for every \(m\),
\[
 \mathcal I(m\varphi)=\cO_X\quad(m\geq1),
 \qquad\text{hence}\qquad
 \varphi\sim_{\mathcal I}0 .
\]

\emph{The nonpluripolar mass is \(1/2\).}  The set \(F\) is polar,
hence pluripolar, and \(\mu(F)=1\).  Also \(\omega(F)=0\), since \(F\)
has vanishing Lebesgue measure.  By
\Cref{lem:ch51-npp-bidegree-one} and
\Cref{thm:complex-npp-ledger}, the nonpluripolar product charges no
pluripolar set, so
\[
 \MA_\omega(\varphi)
 =\one_{X\setminus F}\bigl(\tfrac12\omega+\tfrac12\mu\bigr)
 =\tfrac12\omega,
 \qquad
 M_\omega(\varphi)=\tfrac12 .
\]
On the other hand \(M_\omega(0)=\int_X\omega=1\).

\emph{Conclusion: \(\varphi\not\sim_P0\).}  Both \(\varphi\) and \(0\)
have positive mass in the K\"ahler class \(\{\omega\}\), so
\Cref{thm:complex-p-projection} applies to both and
\Cref{thm:ch51-hierarchy-strict}, item \cref{itm:ch51-hs:mass}, gives:
if \(\varphi\sim_P0\), then
\(M_\omega(\varphi)=M_\omega(0)\), that is \(\tfrac12=1\).  Hence
\[
 \varphi\sim_{\mathcal I}0,
 \qquad
 \varphi\not\sim_P0,
 \qquad
 \varphi\not\sim0 .
\]
Explicitly, \(P_\omega[\varphi]\) is a model potential with
\(M_\omega(P_\omega[\varphi])=\tfrac12\) and
\(P_\omega[0]=V_\omega=0\) with mass \(1\), so the two envelopes differ.
Since \(\omega+\ddc P_\omega[\varphi]\) has total mass \(1\) and
nonpluripolar mass \(\tfrac12\), it carries mass \(\tfrac12\) on a polar
set; in particular \(P_\omega[\varphi]\) is unbounded, although all its
Lelong numbers vanish.

\emph{Direction of the failure.}  One implication does survive:
\(\varphi\leq0\) gives \(\varphi\preceq0\), hence
\(\varphi\preceq_P0\).  What fails is the reverse comparison
\(0\preceq_P\varphi\), even though \(0\preceq_{\mathcal I}\varphi\) holds.
So the second arrow in \eqref{eq:complex-comparison-hierarchy} is not
reversible, and it is not reversible even after restricting to potentials
of positive mass in a K\"ahler class on a curve.
\end{example}

\begin{remark}[Why the example needs a curve, and why that is enough]
\label{rem:ch51-cantor-dimension}
The construction is one-dimensional for two reasons.  First,
\Cref{lem:ch51-npp-bidegree-one} makes the mass computation a
restriction, with no truncation analysis.  Second, on a curve there are
no divisors over \(X\) other than points, so ``all divisorial Lelong
numbers vanish'' is a statement one can verify by inspecting a measure.
Neither restriction weakens the conclusion: multiplying \(\varphi\) by the
pullback construction \(p_1^*\varphi\) on \(\mathbf P^1\times Y\), and
adding a K\"ahler form in the second factor, produces the same separation
in every dimension \(\geq1\), by the smooth-pullback formula
\eqref{eq:verified-smooth-pullback-ideals}.
\end{remark}

\subsection{What the example implies downstream}
\label{sec:ch51-cantor-consequences}

\begin{corollary}
\label{cor:ch51-no-i-invariant-mass}
There is no function \(\Phi\) defined on
\(\PSH(\mathbf P^1,\omega)/{\sim_{\mathcal I}}\) with
\(\Phi([\varphi])=M_\omega(\varphi)\).
\end{corollary}

\begin{proof}
The classes of \(\varphi\) and \(0\) in
\Cref{ex:ch51-cantor} coincide, while the masses are \(\tfrac12\) and
\(1\).
\end{proof}

This is the elementary form of \Cref{thm:verified-mass-no-go}, which
upgrades it to a no-go statement for convex-body invariants of arbitrary
dimension.  It has three consequences that are used repeatedly.

\emph{The analytic trace is \(P\)-valued although its input dependence is
\(\mathcal I\)-theoretic.}  The trace of
\Cref{def:operations-analytic-trace} is constructed from
quasi-equisingular approximants, and by
\Cref{prop:ch51-qes-preserves-i} those approximants remember only the
\(\mathcal I\)-class.  However,
\Cref{thm:operations-trace-well-defined} proves that every output is
\(\mathcal I\)-good and that two choices are \(P\)-equivalent, while
\Cref{thm:operations-trace-p-descent} produces the map
\[
 D_Y/{\sim_{\mathcal I}}
 \longrightarrow
 \QPSH(\widetilde Y)/{\sim_P}.
\]
There is no contradiction with \Cref{cor:ch51-no-i-invariant-mass}:
that corollary excludes an arbitrary mass rule on the entire
\(\mathcal I\)-quotient, whereas trace is defined on the special domain
\(D_Y\), its values form \(\mathcal I\)-good classes, and the
\(P\)-uniqueness is obtained from the restriction theorem.

\emph{Volumes and masses are different invariants.}  By
\Cref{thm:ch51-hierarchy-strict}, items \cref{itm:ch51-hs:mass} and
\cref{itm:ch51-hs:ideal-volume}, the mass is a \(P\)-invariant and the
ideal volume is an \(\mathcal I\)-invariant, and
\eqref{eq:complex-mass-volume} is the only relation between them.  In
\Cref{ex:ch51-cantor}, \(\vol(\omega,\varphi)=1\) while
\(M_\omega(\varphi)=\tfrac12\), since
\(P_\omega[\varphi]_{\mathcal I}=P_\omega[0]_{\mathcal I}=0\).  This is
why \Cref{prop:tna-homogeneous-volume} computes
\(\vol(\theta,\varphi)\) and not \(M_\theta(\varphi)\), and why
\Cref{chap:transcendental-okounkov-dh} measures partial Okounkov bodies
against the former.

\emph{Zero mass is not an \(\mathcal I\)-condition.}  Combining the two
previous points with \Cref{ex:ch51-p2-lelong}, one sees that the
condition \(M_\theta(\varphi)=0\) cannot be read off from the multiplier
ideals: in \Cref{ex:ch51-cantor} a potential \(\mathcal I\)-equivalent to
\(0\) has mass \(\tfrac12<1\), and iterating the construction with the
measure \(\mu\) alone, that is with
\(\varphi'=G_\mu\in\PSH(X,\omega)\), gives
\[
 \omega+\ddc\varphi'=\mu,
 \qquad
 M_\omega(\varphi')=0,
 \qquad
 \varphi'\sim_{\mathcal I}0,
 \qquad
 M_\omega(0)=1 .
\]
So an \(\mathcal I\)-class contains potentials of every mass in
\([0,1]\); indeed \(\varphi_s:=sG_\mu\) has
\(M_\omega(\varphi_s)=1-s\) for \(s\in[0,1]\), and all of them are
\(\mathcal I\)-equivalent to \(0\).

\subsection{What is not claimed}
\label{sec:ch51-not-claimed}

Four statements that are easy to conflate with the above are not made
here and are not true in the stated generality.

\emph{Not claimed: \(\sim_{\mathcal I}\) implies \(\sim_P\) for the
ambient inputs on the domain of the trace.}  The domain \(D_Y\) of
\Cref{def:operations-trace-domain} is \(\mathcal I\)-saturated, by
\eqref{eq:operations-domain-saturated}, but it is not a locus on which
the two relations agree.  What the trace theorem proves is subtler:
\(\mathcal I\)-equivalent \emph{inputs} have \(P\)-equivalent
\emph{traces}.  It does not make the inputs themselves \(P\)-equivalent.
\Cref{ex:ch51-cantor} takes place on a curve where no trace is being taken.

\emph{Not claimed: the \(\mathcal I\)-envelope preserves mass.}  Only the
\(P\)-envelope does, by \eqref{eq:complex-p-projection}.  In
\Cref{ex:ch51-cantor} the \(\mathcal I\)-envelope raises the mass from
\(\tfrac12\) to \(1\).

\emph{Not claimed: potentials with equal Lelong numbers have equal
Monge--Amp\`ere measures, even up to pluripolar sets.}  They need not even
have equal masses.

\emph{Not claimed: any of this is special to non-algebraic behaviour.}
The potential of \Cref{ex:ch51-cantor} lives on \(\mathbf P^1\) in a
K\"ahler --- indeed ample --- class.  The pathology is analytic, not
geometric, and it cannot be removed by projectivity hypotheses.  This is
the reason the comparison theorems of
\Cref{chap:dxz-bj-comparison} are stated for \(\mathcal I\)-model
objects rather than for arbitrary potentials.

\section{The interface with the supplied complex book}
\label{sec:ch51-interface}

\begin{remark}[The exported interface]
\label{chk:complex-interface}
The later chapters of \cref{part:complex} use exactly the following facts.
\begin{enumerate}
\item Model potentials and test-curve levels always have positive mass.
\item The \(P\)-envelope governs nonpluripolar mass and geodesic full-mass
      classes.
\item The \(\mathcal I\)-envelope governs divisorial Lelong data and is
      compatible with K\"ahler perturbations.
\item \(P\)-comparison implies \(\mathcal I\)-comparison, not conversely.
\item Increasing suprema are handled by Choquet reduction; decreasing
      limits require an explicit convergence theorem.
\end{enumerate}
\end{remark}

\subsection{What is imported, item by item}
\label{sec:ch51-import-ledger}

The following table is the complete list of statements that this chapter
uses without proof, together with the exact source and the place of first
use.  Nothing else in \cref{part:complex,part:comparisons} may be
justified by an appeal to ``standard complex pluripotential theory''
without being added to this list.

\begingroup
\begin{longtable}{@{}p{.30\textwidth}p{.30\textwidth}p{.32\textwidth}@{}}
\toprule
imported statement & source & first used in\\
\midrule
\endfirsthead
\toprule
imported statement & source & first used in\\
\midrule
\endhead
Bedford--Taylor products of locally bounded psh functions; monotone
continuity; plurifine locality &
\cite{BT76,BT82,BT87} &
\eqref{eq:complex-npp-truncation},
\Cref{thm:complex-npp-ledger}\\
Negligible sets are pluripolar; locally pluripolar is pluripolar &
\cite[Theorem~7.1]{BT82}, \cite{Vu19} &
\Cref{prop:ch51-psh-operations}, item \cref{itm:ch51-op:sup}\\
Choquet's lemma &
\cite{DemBook} &
\Cref{cor:ch51-choquet}\\
Equivalence of the sup, mean and mass descriptions of the Lelong number &
\cite[Chapter~III]{DemBook}, \cite{Lel68,Dem93} &
\Cref{prop:ch51-lelong-three}\\
Siu semicontinuity and the Siu decomposition &
\cite{Siu74} &
\Cref{sec:ch51-divisorial}\\
Skoda integrability bounds &
\cite{Sko72} &
\Cref{prop:ch51-lelong-calculus}, item \cref{itm:ch51-lc:smooth}\\
Coherence of multiplier ideals &
\cite{Nad89,Nad90} &
\Cref{prop:ch51-ideal-formal}, item \cref{itm:ch51-if:coherent}\\
Subadditivity of multiplier ideals &
\cite{DEL00,OT87} &
\Cref{prop:ch51-ideal-formal}, item \cref{itm:ch51-if:product}\\
Strong openness &
\cite[Theorem~1.1]{GZ15} &
\eqref{eq:complex-strong-openness}\\
Valuative criterion for multiplier ideals &
\cite[Corollary~10.18]{Bou17} &
\Cref{prop:ch51-valuative-criterion}\\
Howald's monomial formula, analytic form &
\cite{Guenancia12,PAG2} &
\Cref{prop:ch51-howald}\\
Demailly regularization; quasi-equisingular refinement &
\cite{Dem92reg,DPS01} &
\Cref{prop:complex-quasi-equisingular}\\
Existence of a K\"ahler current with analytic singularities in a big
class; the ample locus &
\cite{Dem92reg,DPS01,Bou04} &
\Cref{prop:ch51-vtheta-bounded-iff}, item \cref{itm:ch51-vtheta:big}\\
Positivity of the volume characterizes bigness &
\cite{Bou02,Bou02b} &
\Cref{prop:ch51-mass-inequalities}, item \cref{itm:ch51-mi:minimal}\\
Divisorial Zariski decomposition; nonnegative intersection with curves of
zero generic Lelong number &
\cite{Bou04} &
\Cref{ex:ch51-blowup-mass}\\
Unrestricted monotonicity of nonpluripolar masses &
\cite[Theorem~1.1]{DDNL18mono};
\cite[Theorem~1.16]{BEGZ10} for small unbounded locus &
\Cref{thm:complex-npp-ledger}\\
Increasing-sequence continuity of nonpluripolar products &
\cite{DDNL18mono} &
\Cref{prop:complex-mass-continuity}\\
Contact property of the model envelope &
\cite{DDNL18mono,DDNLsurv} &
Step 3 of \Cref{thm:complex-p-projection}\\
Domination principle in relative full-mass classes &
\Cref{thm:appH-comparison} &
Step 4 of \Cref{thm:complex-p-projection}\\
Background-independence of \(\preceq_P\) &
\cite{DDNLmetric,DDNLsurv} &
\Cref{def:complex-p-preorder}\\
Capacity criterion for generalized Cantor sets &
\cite[p.~31]{Car83} &
\Cref{lem:ch51-polar-cantor}\\
\bottomrule
\end{longtable}
\endgroup

\subsection{What is proved here, item by item}
\label{sec:ch51-export-ledger}

The following statements are proved in this chapter from the imports
above, and are the ones that later chapters may cite.

\begin{enumerate}
\item\label{itm:ch51-exp:compactness}
      \(L^1\)-compactness of normalized potentials and convergence of
      decreasing \emph{nets}
      (\Cref{prop:complex-l1-compactness}), with the Hartogs and Choquet
      corollaries (\Cref{cor:ch51-hartogs,cor:ch51-choquet}).
\item\label{itm:ch51-exp:minimal}
      The structure of \(V_\theta\) and the four positivity notions
      (\Cref{def:ch51-positivity,prop:ch51-vtheta-bounded-iff}).
\item\label{itm:ch51-exp:lelong}
      The calculus of Lelong numbers
      (\Cref{prop:ch51-lelong-calculus,prop:ch51-divisorial-calculus})
      and their behaviour along increasing families
      (\Cref{lem:ch51-lelong-increasing}).
\item\label{itm:ch51-exp:ideals}
      The equivalence of multiplier ideals with divisorial Lelong numbers
      (\Cref{thm:complex-strong-openness}), the valuative criterion in
      this book's normalization
      (\Cref{prop:ch51-valuative-criterion}), and the rigidity of
      analytic singularities (\Cref{prop:ch51-analytic-rigid}).
\item\label{itm:ch51-exp:npp}
      The nonpluripolar ledger is \Cref{thm:complex-npp-ledger}, and the
      bidegree-one identity is \Cref{lem:ch51-npp-bidegree-one}.  The mass
      inequalities and increasing continuity of masses along nets are
      \Cref{prop:ch51-mass-inequalities,prop:complex-mass-continuity}.
\item\label{itm:ch51-exp:envelopes}
      The two projection theorems
      (\Cref{thm:complex-p-projection,thm:complex-i-projection}), the
      elementary calculus of both envelopes
      (\Cref{prop:ch51-p-elementary,prop:ch51-i-calculus}), and the
      reformulation of \(\preceq_P\)
      (\Cref{prop:ch51-p-preorder-reformulation}).
\item\label{itm:ch51-exp:hierarchy}
      The comparison hierarchy with both strictness statements
      (\Cref{cor:complex-comparison-hierarchy,thm:ch51-hierarchy-strict,%
      cor:ch51-quotient-maps}), and the no-go corollary
      \Cref{cor:ch51-no-i-invariant-mass}.
\item\label{itm:ch51-exp:examples}
      The worked examples
      \Cref{ex:ch51-p2-lelong,ex:ch51-p2-ideals,ex:ch51-blowup-mass,%
      ex:ch51-nef-zero-volume,ex:ch51-p1-envelope,ex:ch51-cantor}.
\end{enumerate}

\subsection{What is deferred, and to where}
\label{sec:ch51-deferred}

Relative energy, the \(I\)-functional, the metrics \(d_1\) and \(d_S\),
weak geodesics, geodesic rays, test curves and the Ross--Witt Nystr\"om
correspondence are \emph{not} developed here.  They are stated in
\Cref{sec:appH-relative-energy,sec:appH-geodesics,sec:appH-test-curves}
and proved in \Cref{chap:rays-test-curves}.  The reason for the division
is that all of them presuppose a positive-mass model potential as a
reference, so they belong after
\Cref{thm:complex-p-projection} and not before it.  The trace operator
belongs to \Cref{chap:operations-trace} for the same reason, with the
additional input of quasi-equisingular approximation from
\Cref{prop:complex-quasi-equisingular}.

Nothing about compact K\"ahler \emph{spaces}, as opposed to manifolds, is
asserted anywhere in this chapter; that material begins in
\Cref{chap:tropical-spectra-complex,chap:mp-dual-complexes}.

\section{Chapter dictionary}
\label{sec:ch51-dictionary}

\begingroup
\begin{longtable}{p{.24\textwidth}p{.36\textwidth}p{.30\textwidth}}
\toprule
object & content & where\\
\midrule
\endfirsthead
\toprule
object & content & where\\
\midrule
\endhead
\(\ddc\) &
\(\rmi\partial\bar\partial/2\pi\); \(\ddc\log|f|^2=[\divis(f)]\) &
\eqref{eq:complex-ddc-normalization},
\Cref{sec:ch51-normalizations}\\
metric potential &
\(|s|^2_h=\rme^{-\varphi}|s|^2_{h_0}\); norm weight is \(\varphi/2\) &
\Cref{sec:ch51-squared-norm}, \eqref{eq:appJ-complex-weight}\\
\(\theta\)-psh &
\(\varphi\) usc with \(\theta+\ddc\varphi\geq0\); bijection with closed
positive currents in \(\{\theta\}\) modulo constants &
\Cref{def:complex-theta-psh}, \eqref{eq:complex-current-potential}\\
compactness &
normalized potentials are \(L^1\)-compact; decreasing \emph{nets}
converge &
\Cref{prop:complex-l1-compactness},
\Cref{cor:ch51-hartogs}\\
\(V_\theta\) &
largest nonpositive potential; bounded iff \(\{\theta\}\) nef &
\eqref{eq:complex-vtheta},
\Cref{prop:ch51-vtheta-bounded-iff}\\
Lelong number &
\(\nu(\log|z|^2,0)=1\); slope of the spherical mean at \(-\infty\) &
\eqref{eq:complex-lelong-number},
\Cref{prop:ch51-lelong-three}\\
increasing limits &
\(\nu(\uscsup_i\varphi_i,E)=\inf_i\nu(\varphi_i,E)\); no decreasing
analogue &
\Cref{lem:ch51-lelong-increasing},
\Cref{warn:ch51-lelong-decreasing}\\
multiplier ideal &
weight \(\rme^{-\varphi}\);
\(\mathcal I(\lambda\log|z|^2)=(z^{\lfloor\lambda\rfloor})\) &
\eqref{eq:complex-multiplier-ideal},
\eqref{eq:complex-multiplier-example}\\
strong openness &
ideals at integral scales determine all divisorial Lelong numbers, and
conversely &
\Cref{thm:complex-strong-openness},
\Cref{prop:ch51-valuative-criterion}\\
monomial weights &
\(z^k\in\mathcal I(c\log|\mathfrak a|^2)\iff
k+\mathbf 1\in\Int(c\operatorname{Newt}\mathfrak a)\) &
\Cref{prop:ch51-howald},
\Cref{ex:ch51-p2-ideals}\\
analytic singularities &
\(\preceq\) and \(\preceq_{\mathcal I}\) agree between two such &
\Cref{def:complex-analytic-singularities},
\Cref{prop:ch51-analytic-rigid}\\
nonpluripolar product &
increasing limit of truncations; closed positive, plurifine local, no
mass on pluripolar sets &
\eqref{eq:complex-npp-truncation},
\Cref{thm:complex-npp-ledger}\\
bidegree one &
\(\langle\theta_\varphi\rangle
=\one_{\{\varphi>-\infty\}}\theta_\varphi\) &
\Cref{lem:ch51-npp-bidegree-one}\\
mass &
\(M_\theta(\varphi)=\int_X\langle\theta_\varphi^n\rangle\);
monotone in \(\preceq\); \(\leq\vol(\{\theta\})\) &
\eqref{eq:complex-npp-mass},
\Cref{prop:ch51-mass-inequalities}\\
volume versus \(\alpha^n\) &
equal for nef; \(\vol=1\) and \(\int\alpha^2=0\) on \(\operatorname{Bl}_p\mathbf P^2\)
&
\Cref{ex:ch51-blowup-mass},
\Cref{warn:ch51-mass-not-cohomological}\\
continuity &
increasing nets: masses converge; decreasing: they need not &
\Cref{prop:complex-mass-continuity},
\Cref{warn:ch51-continuity-rules}\\
rooftop &
\(\varphi\wedge_\theta\psi\); contact property &
\eqref{eq:complex-rooftop},
\Cref{lem:ch51-rooftop-contact}\\
\(P\)-envelope &
\(P_\theta[\varphi]\); idempotent, mass preserving, characterized by mass
&
\eqref{eq:complex-p-envelope},
\Cref{thm:complex-p-projection}\\
model potential &
\(P_\theta[\phi]=\phi\) \emph{and} \(M_\theta(\phi)>0\) &
\Cref{def:complex-p-envelope},
\Cref{warn:ch51-p-positive-mass}\\
\(\mathcal I\)-envelope &
\(P_\theta[\varphi]_{\mathcal I}\); idempotent, background-free
constraint, \(\geq P_\theta[\varphi]\) &
\eqref{eq:complex-i-envelope},
\Cref{thm:complex-i-projection}\\
ideal volume &
\(\vol(\theta,\varphi)=M_\theta(P_\theta[\varphi]_{\mathcal I})
\geq M_\theta(\varphi)\) &
\eqref{eq:complex-current-volume},
\eqref{eq:complex-mass-volume}\\
hierarchy &
\(\preceq\Rightarrow\preceq_P\Rightarrow\preceq_{\mathcal I}\); both
strict &
\Cref{cor:complex-comparison-hierarchy},
\Cref{thm:ch51-hierarchy-strict}\\
first strictness &
\(-\log(-\log|z|^2)\): full mass, \(P_\omega[\eta]=0\),
\(\eta\not\sim0\) &
\Cref{ex:complex-loglog-envelope},
\Cref{rem:ch51-loglog-explicit}\\
second strictness &
polar Cantor potential: \(\varphi\sim_{\mathcal I}0\),
\(M_\omega(\varphi)=\frac12\), \(M_\omega(0)=1\) &
\Cref{ex:ch51-cantor},
\Cref{cor:ch51-no-i-invariant-mass}\\
\bottomrule
\end{longtable}
\endgroup

Four separations should be carried forward.

\emph{Class positivity versus mass positivity.}  A K\"ahler class can
contain a potential of zero mass (\Cref{ex:ch51-p2-lelong}), and a nef
class can consist entirely of such potentials
(\Cref{ex:ch51-nef-zero-volume}).  Neither situation is excluded by any
hypothesis on \(\{\theta\}\) alone, and both invalidate probability
normalization.

\emph{Mass versus cohomology.}  \(\vol(\{\theta\})\) and
\(\int_X\{\theta\}^n\) agree for nef classes and not otherwise
(\Cref{ex:ch51-blowup-mass}).  Every computation that expands a
cohomology class must check nefness first.

\emph{Three singularity types.}  \(\sim\), \(\sim_P\) and
\(\sim_{\mathcal I}\) are pairwise distinct.  The mass is an invariant of
the second and not of the third; the ideal volume is an invariant of the
third.  \Cref{cor:ch51-quotient-maps} records the two surjections and
\Cref{thm:ch51-hierarchy-strict} their strictness.

\emph{Increasing versus decreasing.}  Every monotone statement in this
chapter is an increasing one: \Cref{lem:ch51-lelong-increasing} for
Lelong numbers, \Cref{prop:complex-mass-continuity} for masses,
\Cref{cor:ch51-choquet} for the reduction to sequences.  The decreasing
counterparts are false
(\Cref{warn:ch51-lelong-decreasing,warn:ch51-continuity-rules}) and are
replaced downstream by explicit convergence theorems with energy or
full-mass hypotheses.

\section{Source and provenance}
\label{sec:ch51-source}

\subsection{Source note}
\label{sec:ch51-source-note}

\begin{remark}[Source note]
\label{rem:complex-prerequisites-source}
The constructions and theorems in this chapter are the interfaces from
the supplied complex manuscript's chapters on psh functions,
nonpluripolar products, envelopes, and comparison of singularities.
The foundational nonpluripolar product is from \cite{BEGZ10};
the multiplier-ideal and envelope conventions are retained exactly from
that manuscript.  No claim in this chapter changes the zero-mass scope
of the supplied \(P\)-projection theorem.
\end{remark}

In more detail, and in the order of the chapter.  Quasi-psh functions,
compactness, and \(V_\theta\) follow \cite{GZ17} and
\cite[\S1]{BEGZ10}; the net formulation of
\Cref{prop:complex-l1-compactness} is this book's, for the reason given
in \Cref{warn:ch51-choquet}.  Lelong numbers, their calculus and Siu's
theorems follow \cite[Chapter~III]{DemBook} and \cite{Siu74,Dem93}.
Multiplier ideals follow \cite{Nad89,Nad90,DemBook}, with strong openness
\cite{GZ15} and the valuative criterion \cite{Bou17}; the monomial
formula is Howald's, in the analytic form of \cite{Guenancia12}.
Nonpluripolar products are \cite{BEGZ10}, with the unrestricted
monotonicity theorem of \cite{DDNL18mono}.  The two envelopes, the
projection theorems, and the preorders \(\preceq_P\),
\(\preceq_{\mathcal I}\) are the relative pluripotential theory of
Darvas--Di Nezza--Lu \cite{DDNL18mono,DDNL18big,DDNL18fullmass,%
DDNLmetric}, surveyed in \cite{DDNLsurv}; the supplied complex book
\cite{XiaBook} is the reference used for the exact form of the
statements, and \Cref{app:complex-reference} is its condensed
restatement.  The application of \(\mathcal I\)-model objects to
transcendental non-Archimedean metrics, which is what makes the
\(P\)-versus-\(\mathcal I\) distinction load-bearing rather than
taxonomic, is \cite{DXZ23,DX24,Xia23Operations} and is developed in
\Cref{chap:transcendental-metrics,chap:operations-trace}.

\subsection{What is supplied here and is not in the source}
\label{sec:ch51-additions}

\Cref{lem:ch51-lelong-increasing}: the identity
\(\nu(\uscsup_i\varphi_i,E)=\inf_i\nu(\varphi_i,E)\) for directed
families, proved from convexity of spherical means.  The statement is
folklore, and the sources use it in the special case of increasing
sequences; it is written out here because it is used four times below ---
in Steps 2 and 3 of \Cref{thm:complex-i-projection}, in
\Cref{prop:ch51-i-calculus}, and in
\Cref{prop:ch51-qes-preserves-i} --- and because
\Cref{warn:ch51-lelong-decreasing} shows that the corresponding
decreasing statement is false, so the direction cannot be left implicit.

\Cref{lem:ch51-npp-bidegree-one}: the identification
\(\langle\theta_\varphi\rangle
=\one_{\{\varphi>-\infty\}}\theta_\varphi\).  This is immediate from
plurifine locality and is not usually isolated, but every mass
computation on a curve in \cref{part:complex,part:comparisons} uses it,
including \Cref{ex:ch51-cantor} and the no-go argument in
\Cref{thm:verified-mass-no-go}.

\Cref{prop:ch51-analytic-rigid} and \Cref{rem:ch51-rigidity}: the exact
statement that \(\preceq\) and \(\preceq_{\mathcal I}\) agree
\emph{between two potentials with analytic singularities}, together with
the warning that this is a different statement from
\eqref{eq:complex-analytic-envelope-agreement}.  Conflating the two
produces the false principle that equal divisorial Lelong numbers imply
bounded difference.

\Cref{prop:ch51-mass-inequalities} and
\Cref{warn:ch51-mass-not-cohomological}: the mass inequalities collected
in one place with the nef/big distinction made explicit, and the
observation that \(\vol\) and \(\int\alpha^n\) are different numbers.

The full proofs of \Cref{thm:complex-p-projection} and
\Cref{thm:complex-i-projection}: the sources state these; the step
decomposition given here isolates exactly which step needs positive mass
(Steps 3 and 4 of the former, none of the latter) and exactly which step
needs resolution of singularities (Step 5 of the latter).  This matters
because \Cref{chap:transcendental-metrics} works in classes where the
mass may vanish and uses only the \(\mathcal I\)-envelope.

\Cref{ex:ch51-p2-lelong} and \Cref{ex:ch51-p2-ideals}: a single potential
on \(\mathbf P^2\) with two different analytic singularities, all Lelong
numbers, all monomial valuations, several multiplier ideals, both log
canonical thresholds, and the mass, computed.  The mass is \(0\) in a
K\"ahler class of volume \(9\), which is the cleanest available warning
against reading a mass off a cohomology class.

\Cref{ex:ch51-blowup-mass} and \Cref{ex:ch51-nef-zero-volume}: the
big non-nef class \(H+e\) on \(\operatorname{Bl}_p\mathbf P^2\) with
\(\vol=1\), \(\int\alpha^2=0\), \(\nu(V_\theta,E)=1\), and the mass
profile \(1-\tau^2\); and the nef class \(H-e\) with \(\vol=0\).  These
two examples are what make \Cref{rem:ch51-zero-volume} concrete.

\Cref{ex:ch51-p1-envelope}: the envelope \(P_\omega[\frac12G_0]\) on
\(\mathbf P^1\), computed in closed form, with both Monge--Amp\`ere
measures and both masses.  It shows in one page that the envelope moves
the function, preserves the ordinary type, preserves the mass, changes
the measure, and satisfies the contact property.

\Cref{lem:ch51-polar-cantor} and \Cref{ex:ch51-cantor}: the separation of
\(\sim_P\) from \(\sim_{\mathcal I}\), with the masses computed to be
\(\frac12\) and \(1\), and the family \(sG_\mu\) realizing every mass in
\([0,1]\) inside a single \(\mathcal I\)-class.  The supplied complex
book records that the two relations are different; the explicit witness,
and the resulting \Cref{cor:ch51-no-i-invariant-mass}, are supplied here
because \Cref{chap:operations-trace,chap:verified-extensions} lean on
exactly this fact and a reader who has not seen a witness will not
believe how badly \(\sim_{\mathcal I}\) forgets.

\subsection{Historical note}
\label{sec:ch51-history}

The three notions of singularity type in this chapter entered the subject
at three different moments and for three different reasons.

The oldest is the ordinary one.  Comparison of psh functions up to a
bounded error is already the organizing idea of Lelong's theory of
positive currents \cite{Lel68} and of Siu's analyticity theorem
\cite{Siu74}, and it is the notion under which Demailly's regularization
theorem \cite{Dem92reg} was first formulated.  Its defect was recognized
immediately: it is too fine to be preserved by the approximation
procedures that make the theory usable, and
\Cref{prop:complex-quasi-equisingular} produces approximants that are
faithful only in a weaker sense.

The second is the multiplier-ideal notion.  Multiplier ideals appear in
Nadel's vanishing theorem \cite{Nad89,Nad90}, built on the
H\"ormander--Bombieri--Skoda machinery and on Skoda's integrability
bounds \cite{Sko72}; Demailly and Kollár \cite{DK01} made the
semicontinuity of the associated thresholds a systematic tool, and
Demailly--Ein--Lazarsfeld \cite{DEL00} supplied the subadditivity that
makes the ideals behave like a valuation-theoretic invariant.  The
statement that closes the circle --- that the collection of all
multiplier ideals is equivalent to the collection of all divisorial
Lelong numbers --- became a theorem only with the strong openness theorem
of Guan and Zhou \cite{GZ15}, previously Demailly's conjecture; the
valuative form used here is Boucksom's \cite{Bou17}.  From the point of
view of \cref{part:complex} this is the notion that survives
analytification, which is why \Cref{chap:valuative-analytification} is
built on it.

The third is the model, or \(P\)-, notion, and it is the youngest.  It
emerged from the study of Monge--Amp\`ere equations with prescribed
singularities in big classes: the nonpluripolar calculus of
Boucksom--Eyssidieux--Guedj--Zeriahi \cite{BEGZ10} made
\(M_\theta(\varphi)\) available as an invariant, and the relative theory
of Darvas--Di Nezza--Lu \cite{DDNL18mono,DDNL18big,DDNL18fullmass,%
DDNL19log,DDNLmetric} identified \(P_\theta[\varphi]\) as the correct
projection, proved the mass identity \eqref{eq:complex-p-projection}, and
built the metric geometry of singularity types on it.  The unrestricted
monotonicity theorem of \cite{DDNL18mono} is the technical statement that
makes the whole scheme work, and it is imported here without proof.  The
survey \cite{DDNLsurv} is the current account.

That these three notions had to be separated, rather than merely ordered,
is a discovery of the transcendental non-Archimedean theory.  In
Darvas--Xia--Zhang \cite{DXZ23} and in \cite{Xia21,Xia23Operations} the
inverse system defining \(\PSH^{\NA}(X,\theta)\) is built from
\(\mathcal I\)-model test curves, because the \(\mathcal I\)-envelope is
the operator compatible with K\"ahler perturbation; but the invariants one
wants at the end --- masses, volumes, Duistermaat--Heckman measures ---
are \(P\)-invariants.  The trace operator of Darvas--Xia \cite{DX24},
recorded in \Cref{chap:operations-trace}, is where the mismatch becomes
visible as a nonformal theorem rather than as a choice of language: the
operator depends only on the input \(\mathcal I\)-class, yet its
\(\mathcal I\)-good outputs are unique up to \(P\)-equivalence.
\Cref{ex:ch51-cantor} explains why this conclusion cannot follow from the
abstract quotient hierarchy alone, and \Cref{thm:verified-mass-no-go}
explains why the restriction and goodness inputs are indispensable.

Two threads run out of this chapter.  Toward the analytic side,
\Cref{chap:rays-test-curves} builds geodesic rays and test curves on the
model envelope, and \Cref{chap:transcendental-metrics} assembles them
into transcendental non-Archimedean metrics.  Toward the comparison side,
\Cref{chap:valuative-analytification,chap:dxz-bj-comparison} turn the
multiplier-ideal package into a genuinely non-Archimedean object and
compare it with the Boucksom--Jonsson theory of \cref{part:bj}.  A reader
who wants to see the \(P\)-versus-\(\mathcal I\) distinction do work
immediately should read \Cref{chap:operations-trace} next; a reader who
wants the geometry first should read \Cref{chap:rays-test-curves}.
